% Faithful transcription — original French. Transcribed from the original first printing:
% G. Castelnuovo et F. Enriques, "Sur quelques récents résultats dans la théorie des surfaces
% algébriques", Mathematische Annalen, Band 48 (Leipzig: B. G. Teubner, 1897), pp. 241–316.
% Scan: EuDML http://eudml.org/doc/157819 (printed pp. 241–316; the volume is also at GDZ
% Göttingen, PPN235181684).
%
% \origpage uses the PRINTED page numbers. Running heads, printed page numbers and signature
% marks are page furniture and are not transcribed. The memoir is in French throughout, set in
% Antiqua in a German journal, so there is no Fraktur, no long-ſ and no ß anywhere in it.
%
% Cross-page rendering decisions are recorded in notation.md, which every transcribing batch
% followed: Max Noether's name is printed \emph{Nöther} and kept that way (R12); author names are
% letterspaced in the running text and are therefore \emph'd (R20); footnotes run *) **) ***) †)
% per page, referenced as math superscripts and set as complete units at the end of each page's
% text (R15).
%
% The scan is tall relative to its width (1:1.7), so the pages were transcribed from landscape
% HALF-page crops rather than one image per page: at one image per page the 1568px long-edge cap
% left under 1000px of text width on 72 of the 76 pages. See provenance.yaml.

\documentclass{article}
\usepackage{readmasters}

\begin{document}

\origpage{241}

\section*{Sur quelques récents résultats dans la théorie des surfaces algébriques.}

\subsection*{Par MM.\ G.\ Castelnuovo à Rome et F.\ Enriques à Bologne.}

La \emph{Géométrie sur une surface algébrique générale}, dont l'origine remonte à une remarque de \emph{Clebsch}${}^{*)}$ et à un Mémoire fondamental de M.\ \emph{Nöther}${}^{**)}$, s'est enrichie, dans ces derniers temps, de nouveaux résultats, grâce aux recherches des géomètres français et italiens. Ces recherches ont un même point de départ dans le Mémoire cité de M.\ \emph{Nöther}, mais elles suivent deux directions différentes que nous allons exposer brièvement.

En France on a considéré surtout les fonctions transcendantes qui sont attachées à l'équation d'une surface algébrique. M.\ \emph{Nöther} avait déjà introduit dans l'étude des surfaces certaines intégrales doubles (à différentielles algébriques) qui restent partout finies, et qui jouent le même rôle que les intégrales abéliennes de première espèce relatives aux courbes algébriques. Quelques ans plus tard M.\ \emph{Picard}${}^{***)}$ a donné une nouvelle impulsion à cette théorie, en considérant, à côté des intégrales doubles, certaines intégrales simples de différentielles totales qui sont aussi dignes de remarque. Il a étudié en particulier les intégrales simples qui restent partout finies. Les surfaces les plus connues ne possèdent pas vraiment de telles intégrales; mais M.\ \emph{Picard} a indiqué des classes de surfaces, dont les propriétés remarquables dépendent précisément de la présence de fonctions transcendantes de l'espèce nommée. M.\ \emph{Poincaré} avec une courte Note${}^{\dagger)}$, et M.\ \emph{Humbert}

\textbf{*)} Comptes Rendus de l'Ac.\ d.\ Sc.\ 1868.

\textbf{**)} \emph{Zur Theorie des eindeutigen Entsprechens}... I, II, Math.\ Annalen Bd.\ 2, 8, 1869, 1874.

\textbf{***)} Voir les deux Mémoires \emph{Sur les intégrales de différentielles totales algébriques}..., Journal de Mathém.\ 4${}^{\text{e}}$ s${}^{\text{e}}$, t.\ I, 1885; le \emph{Mémoire sur la théorie des fonctions algébriques}, J.\ d.\ M.\ 4${}^{\text{e}}$ s${}^{\text{e}}$, t.\ V; ainsi que plusieurs Notes dans les Comptes Rendus de l'Ac.\ d.\ Sc., Déc.\ 1884.

\textbf{†)} Comptes Rendus de l'Ac.\ d.\ Sc., Déc.\ 1884.

\origpage{242}
surtout avec plusieurs Mémoires importants${}^{*)}$, ont suivi la même direction, soit en approfondissant la théorie générale, soit en étudiant des classes particulières de surfaces${}^{**)}$. Mais le champ de recherches que M.\ \emph{Picard} a ouvert est loin d'être épuisé; car la question même de décider, en partant des caractères géométriques d'une surface, si la surface possède de telles intégrales simples, n'est pas encore résolue, d'une manière générale.

Les géomètres italiens, au contraire, se sont consacrés à l'étude des systèmes linéaires de courbes algébriques qui appartiennent à une surface algébrique, ou, en d'autres termes, à l'étude des fonctions rationnelles de trois quantités $x\,y\,z$ liées par une rélation algébrique $f(x, y, z) = 0$. C'était un ordre de recherches qui se présentait tout naturellement à leur attention; car en Italie, pendant les dernières années, on avait cultivé les deux théories auxiliaires qui servent de base aux développements nouveaux; à savoir \emph{la géométrie sur une courbe algébrique}, et \emph{la théorie des systèmes linéaires de courbes planes}.

La géométrie sur une courbe, qui a pris naissance dans les travaux immortels de \emph{Riemann} (1857), a été développée plus tard sous la forme géomètrique dans le Mémoire classique de MM.\ \emph{Brill} et \emph{Nöther}${}^{***)}$. Il restait pourtant encore à délivrer cette théorie des considérations projectives qu'on y faisait intervenir trop souvent; c'est ce qu'on a fait en Italie (\emph{Segre}, \emph{Castelnuovo}...) en raisonnant sur les courbes, sans faire attention ni à leurs ordres, ni aux espaces auxquels elles appartiennent, et en limitant l'usage des courbes adjointes, dont la considération peut être utile parfois, mais n'est point nécessaire pour le développement de cette théorie.

La théorie des systèmes linéaires de courbes planes, au point de vue des transformations birationnelles du plan, a commencé par les recherches bien connues de M.\ \emph{Cremona} (1863---1865) et de M.\ \emph{Nöther} (1870), et a été développée par les géomètres italiens \emph{Caporali}, \emph{Bertini}, \emph{Jung}, \emph{Guccia}, \emph{Martinetti}, \emph{Castelnuovo}... On a pu aborder par cette voie l'étude des propriétés du plan qui restent invariables par rapport aux transformations birationnelles; c'est donc la théorie d'une \emph{classe} particulière de surfaces (à savoir des surfaces \emph{rationnelles}), qui a pris naissance par ces recherches.

Mais on s'aperçut bientôt que la plupart des méthodes qu'on y

\textbf{*)} \emph{Théorie générale des surfaces hyperelliptiques}, Journal de Mathématiques, 4${}^{\text{e}}$ s${}^{\text{e}}$, t.\ IX; \emph{Sur quelques points de la théorie des courbes et des surfaces algébriques}, Journ.\ d.\ Math.\ 4${}^{\text{e}}$ s${}^{\text{e}}$, t.\ X; ainsi que plusieurs Notes dans les Comptes Rendus de l'Ac.\ d.\ Sc., 1893---95.

\textbf{**)} Consulter aussi sur ce sujet le Mémoire de M.\ \emph{Nöther}, \emph{Ueber die totalen algebraischen Differentialausdrücke}, Math.\ Annalen, 29, 1886.

\textbf{***)} \emph{Ueber die algebraischen Functionen}..., Math.\ Annalen, Bd.\ VII, 1873.

\origpage{243}
avait employées, pouvaient avoir une application plus étendue, car elles servaient aussi à l'étude de questions analogues sur les surfaces algébriques générales. On parvint ainsi à distinguer deux sortes de propriétés des systèmes linéaires de courbes situées sur une surface. Il y a en effet des propriétés qui ne dépendent pas de la surface, mais bien de la nature du système qu'on y envisage; tandis qu'il y a d'autres propriétés qui appartiennent à tous les systèmes linéaires de courbes situés sur la même surface, et dépendent seulement de la nature de celle-ci. Ce sont ces dernières propriétés qui offrent le plus grand intérêt, car elles fournissent des caractères d'une surface, qui se transportent à toutes les surfaces en correspondance birationnelle avec celle-là, des caractères \emph{invariants} (par rapport aux transformations nommées). C'est à l'aide de telles remarques que l'on retrouva sous une nouvelle lumière les invariants introduits par M.\ \emph{Nöther}, et qu'on parvint à en ajouter d'autres encore.

Ainsi se forma un ensemble de recherches qui sont étroitement liées entre elles, bien qu'elles se trouvent à présent éparses dans plusieurs travaux publiés tout dernièrement en Italie${}^{*)}$.

Réunir ces résultats, les présenter dans l'ordre logique (qui n'est pas toujours l'ordre chronologique), les rapprocher à des propriétés connues des courbes et du plan, pour en faire ressortir tantôt les analogies, tantôt les différences; voilà le but que nous nous sommes proposé en écrivant cette Monographie. Il ne faut donc y chercher des résultats nouveaux, mais bien une exposition nouvelle de résultats déjà publiés. Ce travail pourtant ne paraîtra pas superflu à nos lecteurs, à ce que nous pensons. C'est que les Mémoires originaux que nous résumons ici, ne sont pas d'une lecture facile. Cela dépend de plusieurs causes. Tout d'abord il y a les difficultés qui proviennent de la théorie elle-même; la plupart des propriétés des surfaces sont bien loin de présenter les caractères de simplicité, qu'on trouve en de propositions analogues relatives aux courbes. Mais le rapprochement avec celles-ci est très utile pour mettre en lumière l'esprit des résultats

\textbf{*)} Voici en ordre chronologique les travaux principaux dont nous exposons ici les résultats, pour ce qui concerne la théorie générale des surfaces algébriques; (pour la théorie des surfaces rationnelles nous donnerons ailleurs les citations relatives):

\emph{Enriques}, \emph{Ricerche di geometria sulle superficie algebriche}, Mem.\ dell'Acc.\ d.\ Scienze di Torino s${}^{\text{e}}$ 2${}^{\text{a}}$, t.\ 44, 1893.

\emph{Enriques}, \emph{Introduzione alla Geometria sopra le superficie algebriche}, Mem.\ della Società Italiana d.\ Scienze s${}^{\text{e}}$ III, t.\ X, 1896.

\emph{Castelnuovo}, \emph{Alcuni risultati sui sistemi lineari di curve appartenenti ad una superficie algebrica}, ibid., 1896.

\emph{Castelnuovo}, \emph{Sulle superficie di genere zero}, ibid., 1896.

\emph{Enriques}, \emph{Sui piani doppi di genere uno}, ibid., 1896.

\origpage{244}
nouveaux. Nous recourons sans cesse à ce moyen, et même pour le rendre plus aisé, nous avons exposé d'avance les propositions fondamentales de la géométrie sur une courbe, dans l'ordre qui nous semblait préférable.

Ainsi nous obtenons l'avantage d'exposer à nos lecteurs toutes les notions nécessaires pour aborder l'étude des surfaces. La lecture de cette Monographie n'exigera donc, nous le pensons, que la connaissance de ces notions générales de géomètrie, qui sont familières même à ceux qui se sont consacrés à tout autre ordre de recherches.

Mais pour atteindre ce but, nous nous sommes bornés à présenter seulement les énoncés des théorèmes, et les considérations qui servent le mieux à les mettre en lumière; quant aux démonstrations, nous avons pensé qu'il valait mieux les supprimer tout à fait, en renvoyant aux ouvrages originaux le lecteur désireux de les connaître. C'est qu'on ne réussit pas, pour le moment, à exposer ces démonstrations de manière à éviter toute difficulté. La conception qui inspire ces raisonnements peut bien être simple, dans la plupart des cas; mais le développement exige qu'on envisage tant de cas particuliers, qu'on rebute tant d'objections, que le lecteur fatigué par ce morcellement de la démonstration, perd de vue l'ensemble des résultats. C'est un défaut qu'on peut reprocher d'ordinaire aux théories en formation; et bien souvent le progrès consiste dans l'introduction de certaines conceptions plus amples, de certaines locutions plus convenables, qui permettent d'embrasser par un seul regard un grand nombre de cas qui semblaient distincts.

Le moment n'est pas encore venu pour introduire de tels perfectionnements dans la théorie des surfaces. Ce qui importe à présent, c'est de surmonter les obstacles qu'on rencontre sur le chemin, c'est d'arriver au sommet qui reluit devant nous. S'il y a quelqu'un qui ne se hasarde pas à nous suivre sur une voie désagréable, eh bien qu'il attende! Cependant, nous l'espérons, il lira avec intérêt la description du paysage qui s'offre à nos regards, du point où nous sommes parvenus; il la trouvera dans cette Monographie, qui lui est consacrée.

Mais s'il y a des courageux qui désirent de nous suivre, ils trouveront dans les pages suivantes tous les renseignements pour le faire. Et ils seront récompensés de leur peine, car le champ de recherches est bien loin d'être épuisé, et la région qui reste à explorer est encore très vaste, et peut amener à de brillantes découvertes. Il faut pourtant se rappeler que ce n'est pas d'un seul point de vue qu'on peut embrasser un horizon si vaste. C'est au contraire par les efforts réunis de plusieurs travailleurs suivant des voies différentes, c'est par l'application savante de toutes les ressources que la géométrie et

\origpage{245}
l'analyse mettent aujourd'hui à notre disposition, qu'on peut espérer d'approfondir la théorie des surfaces, et de l'enrichir de nouvelles découvertes.

Alors seulement cette branche sera digne de figurer à côté de la théorie admirable des fonctions algébriques d'une variable, dont la construction est une gloire des géomètres de notre siècle.

\section*{Chapitre I. Transformations birationnelles. Genre d'une courbe et d'une surface d'après Clebsch et M.\ Nöther.}

\subsection*{1. Transformations birationnelles.}

La géométrie projective (théorie des substitutions linéaires et homogènes entre deux groupes de variables) a trouvé son extension la plus naturelle dans la théorie des transformations birationnelles entre deux variétés algébriques.

On sait bien comment ces transformations sont définies. Soient $x_{1}, x_{2}, \ldots , x_{r}$, et $y_{1}, y_{2}, \ldots , y_{s}$ deux groupes de variables, que nous pouvons regarder comme les coordonnées homogènes d'un point $x$ dans un espace linéaire $X$ à $r-1$ dimensions, et d'un point $y$ dans un espace $Y$ à $s-1$ dimensions; établissons entre ces variables des relations de la forme
\[ \varrho\, y_{i} = f_{i}(x_{1}, x_{2}, \ldots , x_{r}) \qquad (i = 1, 2, \ldots , s) \tag{1} \]
où les $f_{i}$ sont des formes algébriques toutes du même degré dans les $x$, et $\varrho$ est un facteur de proportion; on aura alors une \emph{transformation rationnelle} entre l'espace $X$ et une certaine variété algébrique de l'espace $Y$. A tout point général $x$ de $X$ correspond un (seul) point $y$ de l'espace $Y$; tandis que $x$ décrit dans son espace une ligne, ou surface... algébrique $F$, le point $y$ décrit en $Y$ (en général) une ligne, ou surface... algébrique $F'$ qui correspond à $F$. Tout point général $y$ de $F'$ peut provenir cependant d'un, ou de plusieurs (même d'infinis) points $x$. Si c'est le premier cas qui arrive, alors, à l'aide des (1) et des équations qui définissent la $F$ en $X$, on peut exprimer à leur tour les variables $x$ en fonctions rationnelles des $y$, par des relations de la forme
\[ \sigma\, x_{k} = \varphi_{k}(y_{1}, y_{2}, \ldots , y_{s}) \qquad (k = 1, 2, \ldots , r). \tag{2} \]

Les (1) et (2) définissent une \emph{correspondance}, ou \emph{transformation birationnelle} entre $F$ et $F'$.

\origpage{246}

En particulier, on peut avoir une transformation birationnelle (ou \emph{de Cremona}) entre deux espaces linéaires (ayant les mêmes dimensions); les (1) et (2) définissent une telle transformation entre $X$ et $Y$, si $r = s$, et si en outre l'on peut deduire les (1) des (2), et celles-ci de celles-là, pour \emph{toutes} les valeurs des $x$ et des $y$.

\subsection*{2. Eléments exceptionnels d'une transformation.}

En revenant aux variétés $F$ et $F'$ en correspondance birationnelle, il y a deux cas à distinguer, suivant que les $F$, $F'$ sont des courbes, ou bien des variétés à plus d'une dimension. Dans le premier cas en effet, on a entre les courbes $F$, $F'$ une correspondance biunivoque \emph{sans exceptions} d'une nature essentielle. Mais si $F$, $F'$ sont par exemple des surfaces (à deux dimensions), il peut bien exister sur $F$ (ou $F'$) des points simples (en nombre fini), à chacun desquels vont correspondre sur $F'$ (ou $F$) tous les points d'une courbe${}^{*)}$. Sur une surface, une courbe qui par une transformation birationnelle convenable peut être changée en un point simple d'une autre surface, se distingue, par bien de propriétés, des autres courbes de la première surface; nous la nommerons \emph{courbe exceptionnelle}. Les courbes exceptionnelles (dans la plupart des cas) ne diffèrent pas des \emph{ausgezeichnete Curven}, considérées par M.\ \emph{Nöther}${}^{**)}$; c'est ce que nous verrons plus tard.

Supposons, en particulier, qu'au point $P$ de $F$ corresponde la courbe exceptionnelle $P'$ de $F'$; et soit $C$ une courbe quelconque de $F$ passant par $P$. Les points de $C$, autres que $P$, correspondent aux points d'une courbe $C'$ de $F'$; on pourrait vraiment regarder $C'$ comme la courbe correspondante de $C$ sur $F$. Mais d'ordinaire il convient de joindre à la courbe $C'$, la courbe exceptionnelle $P'$ (qui provient du point $P$ de $C$), et d'envisager la courbe ainsi composée comme la transformée de $C$. C'est là une convention que nous allons respecter, en général, dans la suite.

\subsection*{3. Répartition en classes des variétés algébriques.}

Par rapport aux transformations birationnelles, les variétés algébriques, ayant un même nombre de dimensions, se répartissent en classes; nous dirons, avec \emph{Riemann}, que deux variétés appartiennent

\textbf{*)} La projection stéréographique d'une surface du second ordre sur un plan donne un exemple bien simple de telles courbes.

\textbf{**)} Mém.\ cité, Mathem.\ Annalen 8, pag.\ 521.

\origpage{247}
à la même classe lorsqu'on peut établir une correspondance birationnelle entre elles. On aura donc, par exemple, à considérer la classe des \emph{courbes rationnelles} (ou unicursales) qui peuvent être transformées birationnellement en une droite, la classe des \emph{surfaces rationnelles} (ou unicursales, homaloïdes) qui sont en correspondance birationnelle avec un plan, etc.

Les variétés d'une même classe diffèrent entre elles, il est vrai, par leurs caractères projectifs (ordre, dimension de l'espace auquel elles appartiennent,...); mais elles ont plusieurs propriétés communes, dont la recherche forme le sujet de la \emph{Géométrie sur les variétés algébriques}; cette branche s'occupe donc des propriétés invariables par rapport aux transformations birationnelles. Il s'ensuit que lorsqu'on étudie une \emph{classe} de variétés, on peut se rapporter toujours à une variété abstraite, générale, de la classe, sans faire attention du tout à ses caractères projectifs; on pourrait même dire qu'on envisage le \emph{champ algébrique} dont les variétés de la classe (courbes, surfaces...) sont des \emph{images projectives}; il n'y aurait là pourtant qu'un nouveau mot, pas une conception nouvelle${}^{*)}$.

La considération de la variété générale d'une classe a ce grand avantage, qu'elle permet de faire abstraction de toutes les singularités de nature projective qu'une variété particulière peut présenter. Ainsi, dans la suite, lorsque nous parlerons des points de la variété générale d'une classe, nous allons supposer toujours qu'il s'agit de points \emph{simples}; on n'exclut pas naturellement que ces points ne peuvent correspondre à des points multiples sur des variétés particulières de la classe; mais on affirme simplement qu'un point multiple est une particularité projective de quelque représentant de la classe, pas de la classe elle-même; c'est, pour ainsi dire, un défaut de l'image, pas du modèle.

On va trouver pourtant, dans la suite, quelques paragraphes où l'on envisage vraiment une variété donnée par ses caractères projectifs (une courbe \emph{plane} avec ses points multiples, une surface \emph{de l'espace ordinaire} d'ordre connu, etc.); c'est que ces paragraphes n'appartiennent pas proprement à notre théorie, mais à des théories auxiliaires. On pourrait même les supprimer, si l'on ne redoutait pas la complication qui s'ensuivrait.

\subsection*{4. Genre d'une courbe et d'une surface algébrique.}

Pour effectuer la répartition en classes, dont nous venons de parler, il faut connaître des caractères numériques qui gardent la même valeur

\textbf{*)} A ce point de vue l'on peut regarder l'étude d'une classe de variétés comme l'étude d'un \emph{champ de rationnalité} (d'après \emph{Kronecker}) ou d'un \emph{corps de fonctions algébriques} (d'après \emph{Dedekind} et \emph{Weber}).

\origpage{248}
pour toutes les variétés d'une même classe. On sait bien comment \emph{Riemann}, en étudiant au point de vue de l'\emph{Analysis situs} la surface réelle qui représent les points (réels et imaginaires) d'une \emph{courbe algébrique}, est parvenu au plus important de ces caractères, qu'il a nommé \emph{genre de la courbe} (ou de la relation algébrique correspondante)${}^{*)}$. \emph{Clebsch} et \emph{Gordan}, en exposant (1866) la même conception sous la forme géométrique, ont donné du genre la définition bien connue qui suit.

Soit $C$ une courbe plane de l'ordre $n$; supposons par simplicité que $C$ n'ait d'autres points multiples, en dehors de $d$ points doubles. Considérons les courbes d'ordre $n-3$ qui passent simplement par les $d$ points nommés, courbes qu'on appelle \emph{adjointes} à $C$; \emph{le nombre des courbes adjointes qui sont linéairement distinctes est le genre $p$ de $C$.}

A côté de cette définition \emph{géométrique}, on peut donner une définition \emph{numérique} de $p$. Il suffit de remarquer que le nombre des paramètres qui entrent d'une façon homogène dans l'équation d'une courbe de l'ordre $n-3$, assujettie à passer par $d$ points donnés, est
\[
p = \frac{(n-1)(n-2)}{2} - d.
\]

On suppose vraiment ici que les $d$ points doubles de $C$ présentent autant de conditions \emph{distinctes} aux courbes d'ordre $n-3$, que l'on contraint à passer par eux; cela n'est pas évident \emph{a priori}, mais on a démontré plus tard (\emph{Brill} et \emph{Nöther}${}^{**)}$) que la supposition est toujours vérifiée, lorsque on se borne aux courbes irréductibles.

Même si l'on fait abstraction de cette remarque, il est bon d'observer que les deux définitions amènent également à des invariants d'une courbe plane; en effet, on a donné plusieurs démonstrations de l'invariance du genre, dont quelques-unes (\emph{Riemann}, \emph{Clebsch} et \emph{Gordan},...) se rapportent à la première définition, d'autres (\emph{Cremona}, \emph{Bertini}, \emph{Zeuthen},...) à la seconde.

L'extension de ces définitions aux surfaces algébriques est immédiate, dans les cas les plus simples; \emph{Clebsch}${}^{***)}$ et M.\ \emph{Nöther}${}^{\dagger)}$ l'ont indiqué.

Soit $F$ une surface algébrique d'ordre $n$ de l'espace ordinaire; supposons que $F$ n'ait d'autres singularités qu'une courbe double de l'ordre $d(\geqq 0)$ et genre $\pi$, et un certain nombre fini $t(\geqq 0)$ de points triples pour la surface, et triples aussi pour la courbe double nommée. Considérons maintenant les surfaces d'ordre $n-4$ qui passent simplement par la courbe double (et en conséquence deux fois, en général, par ses points triples), surfaces que nous appelerons \emph{adjointes} à $F$; \emph{le}

\textbf{*)} On pourrait faire remonter la notion du genre au théorème d'\emph{Abel}, bien que \emph{Abel} lui-même ne parle pas expressément de caractères invariants.

\textbf{**)} \emph{Ueber die algebraischen Functionen}..., Math.\ Annalen Bd.\ 7, 1873.

\textbf{***)} Comptes Rendus de l'Acad.\ d.\ Sc., Déc.\ 1868.

\textbf{†)} Mémoire cité, \emph{Zur Theorie}...; Mathem.\ Annalen 2, 8.

\origpage{249}
\emph{nombre de ces surfaces linéairement distinctes} est (d'après la démonstration de M.\ \emph{Nöther}) un caractère invariant (par rapport aux transformations birationnelles) de la surface $F$; on l'appelle le \emph{genre géométrique} de $F$, et on le désigne par $p_{g}$.

Or si nous remarquons qu'une surface d'ordre $m$ \emph{assez élevé} doit satisfaire à
\[ md - 2t - \pi + 1 \]
conditions pour contenir la courbe double de $F$, nous sommes portés (avec \emph{Cayley}${}^{*)}$) à former l'expression
\[ p_{n} = \frac{(n-1)(n-2)(n-3)}{6} - (n-4)d + 2t + \pi - 1, \]
dont la valeur est exactement $p_{g}$, si $m = n - 4$ est assez élevé pour que l'on puisse appliquer la formule qui précède. Que la dernière condition soit remplie ou non, \emph{l'expression de $p_{n}$ a toujours la propriété d'invariance;} c'est ce que MM.\ \emph{Zeuthen} et \emph{Nöther}${}^{**)}$ ont montré. La même expression est parfaitement analogue à celle considerée ci-dessus en parlant des courbes planes; on lui a donné le nom de \emph{genre numérique} de la surface $F$.

Mais ici s'arrête l'analogie avec les courbes planes; en effet nous ne pouvons plus affirmer l'égalité entre les valeurs de $p_{g}$ et $p_{n}$; au contraire \emph{il peut bien arriver que l'on ait $p_{g} \neq p_{n}$ et précisément $p_{g} > p_{n}$.} Par exemple une surface réglée ayant les sections planes de genre $p$, a (d'après \emph{Cayley}${}^{***)}$) $p_{g} = 0$, $p_{n} = -p$. On a porté plus tard d'autres exemples de surfaces avec $p_{g} > p_{n}$; nous y reviendrons dans la suite.

\subsection*{5. Remarques sur les définitions qui précèdent.}

M.\ \emph{Nöther} dans le Mémoire cité plusieurs fois, avance dans l'étude des surfaces, en considérant les courbes découpées sur $F$ par les surfaces adjointes d'ordre $n - 4$. Nous ne le suivrons pas, pour le moment, dans cette direction; nous aimons mieux montrer quelques difficultés auxquelles amènent les définitions des genres d'une surface.

Nous nous sommes bornés en les rappelant, à considérer des surfaces ayant des singularités bien simples; on pourrait vraiment, envisager aussi des cas plus compliqués, mais il s'en faut de beaucoup que l'on puisse toujours préciser d'une manière directe le comportement

\textbf{*)} \emph{On the Deficiency of certain Surfaces}, Mathem.\ Annalen, 3, 1871.

\textbf{**)} \emph{Zeuthen}, \emph{Etude géométrique}..., Mathemat.\ Annalen, 4, 1871; \emph{Nöther}, \emph{Zur Theorie}..., Math.\ Annalen, 8, pag.\ 526.

\textbf{***)} Mém.\ cité.

\origpage{250}
d'une surface adjointe dans un point singulier de la surface donnée${}^{*)}$. On n'est donc pas en mesure de définir le genre numérique d'une surface ayant des singularités supérieures, ni par suite de montrer dans ce cas la propriété d'invariance de $p_{n}$. Ce serait bien mieux si l'on pouvait donner de $p_{n}$, et même de $p_{g}$, une telle définition qui ne dependît pas des singularités de la surface considérée, et qui fît ressortir sur le champ l'invariance de ces caractères, et le lien entre eux. C'est ce que nous nous proposons d'obtenir. Mais pour plus de clarté il convient de partir des notions les plus simples, et de montrer comment la théorie des courbes et des surfaces algébriques peut se reconstruire sous un unique point de vue.

\section*{Chapitre II. Série linéaire de groupes sur une courbe. Système linéaire de courbes sur une surface.}

\subsection*{6. Série linéaire sur une courbe.}

Dans l'étude d'une \emph{classe} de variétés algébriques, le plus grand intérêt appartient aux \emph{fonctions rationnelles} des équations qui définissent la classe.

En nous bornant, pour le moment, aux courbes, si
\[ f(x, y) = 0 \]
est l'équation d'une courbe plane quelconque de la classe, à toute \emph{fonction rationnelle} $R(x, y)$ de $x$ et $y$ (liées par la $f = 0$) correspond une \emph{série linéaire} $\infty^{1}$ de groupes de points sur la courbe; chaque groupe de la série se compose des points (en nombre fixe) où la $R$ prend une valeur arbitrairement donnée, qui change d'un groupe à l'autre. Le nombre $n$ des points de chaque groupe est appelé \emph{l'ordre} de la $R$ ou de la série; celle-ci est désignée, d'ordinaire, par $g_{n}^{1}$. On peut dire aussi que les groupes de $g_{n}^{1}$ sont les groupes des zéros des fonctions rationnelles représentées par l'expression $R(x, y) + C$, où $C$ est un paramètre; toutes ces fonctions ont les mêmes pôles. Plus généralement considérons plusieurs fonctions rationnelles de notre courbe,
\[ R_{1}(x, y), R_{2}(x, y), \ldots , R_{r}(x, y) \]

\textbf{*)} La méthode analytique proposée à cet effet par M.\ \emph{Nöther} ne fournit pas une définition directe, permettant le calcul de $p_{n}$, car on doit recourir à une transformation de la surface. Voir à ce sujet \emph{Nöther}, Göttinger Nachrichten, 1871, Mathem.\ Annalen, 29, ainsi que \emph{Zeuthen}, \emph{Révision et extension des formules numériques}... Mathem.\ Annalen 10, pag.\ 544, 1876.

\origpage{251}
ayant toutes les mêmes $n$ pôles, et n'étant pas liées par aucune relation linéaire à coefficients constants; formons la fonction
\[ \lambda_{1} R_{1} + \lambda_{2} R_{2} + \cdots + \lambda_{r} R_{r} + C \]
qui dépend des paramètres $\lambda_{1}, \ldots , \lambda_{r}, C$; le groupe des zéros de celle-ci décrit (lorsque on fait varier les paramètres) \emph{une série linéaire $g_{n}^{r}$ d'ordre $n$, et dimension $r$;} $r$ points généraux de la courbe entrent dans un groupe de $g_{n}^{r}$, et dans un seul. Le groupe des pôles appartient aussi à la série, et ne présente aucune particularité dans celle-ci. Il est facile maintenant de définir la plus ample série linéaire d'ordre $n$, à laquelle appartient un groupe de $n$ points donnés, ou, comme nous dirons, la \emph{série complète} (\emph{Vollschaar}) qui est déterminée par ce groupe. Il suffit en effet de considérer \emph{toutes} les fonctions rationnelles d'ordre $n$, qui ont leurs pôles dans ce groupe; les groupes des zéros de ces fonctions forment la série demandée. Il est clair que la dimension de la dernière série ne peut surpasser $n$; si elle atteint cette valeur, la courbe dont il s'agit est rationnelle, car dans le cas contraire on ne saurait fixer arbitrairement les zéros et les pôles d'une même fonction rationnelle.

Les mêmes considérations peuvent être présentées aussi sous la forme géométrique${}^{*)}$. Soit $C$ une courbe algébrique (plane ou de l'espace ordinaire, ou d'un hyperespace); découpons-la par un système linéaire $\infty^{r}$ $(r \geqq 0)$ de variétés (courbes, surfaces...)
\[ \lambda_{1} f_{1} + \lambda_{2} f_{2} + \cdots + \lambda_{r+1} f_{r+1} = 0, \]
où les $f$ sont des formes algébriques toutes du même degré (à trois, quatre... variables). Si chaque $f$ a $n$ points d'intersection avec la $C$, et s'il n'y a dans le système aucune variété qui contient la $C$, les groupes des intersections formeront une série linéaire $g_{n}^{r}$. Plus généralement si avec les groupes de la $g_{n}^{r}$ nous formons des groupes de $n \pm n'$ points, en ajoutant aux premiers groupes les mêmes $n'$ points fixés arbitrairement sur $C$, ou en les retranchant s'ils entrent dans chaque groupe de $g_{n}^{r}$, nous dirons que les nouveaux groupes forment une série linéaire $g_{n \pm n'}^{r}$.

Les variétés $f$ qui nous ont permis de construire la série $g_{n}^{r}$ ne sont pas pourtant liées avec elle (au point de vue de la géométrie sur la courbe), et n'entrent pas du tout dans l'étude de la série. On peut

\textbf{*)} Ainsi que nous l'avons dit dans l'introduction, nous présentons ici la théorie géométrique des séries linéaires (qui est due à MM.\ \emph{Brill} et \emph{Nöther}, Mém.\ cité) suivant le plan adopté dans quelques travaux italiens. On trouvera les développements et les citations relatives dans la Monographie de M.\ \emph{Segre}, \emph{Introduzione alla geometria sopra un ente algebrico}..., Annali di Matematica s${}^{\text{e}}$.\ II, tomo XXII, 1894. Dans le même volume on pourra aussi consulter avec profit une Monographie de M.\ \emph{Bertini}, où le même sujet est exposé au point de vue de \emph{Brill} et \emph{Nöther}.

\origpage{252}
bien construire celle-ci par d'autres systèmes de variétés, soit sur la même courbe $C$, soit sur une courbe $C'$ de la même classe (car une transformation birationnelle entre $C$ et $C'$ change la $g_{n}^{r}$ de $C$ en une $g_{n}^{r}$ de $C'$).

On peut même définir la série linéaire $g_{n}^{r}$ sur une courbe algébrique quelconque, sans mentionner expressément les variétés qui servent à la découper. En effet \emph{c'est une $g_{n}^{r}$ tout système $\infty^{r}$ de groupes de $n$ points sur la courbe, dont les groupes dépendent rationnellement de $r$ paramètres, de manière que la condition de contenir un point arbitrairement fixé sur la courbe se traduise par une relation linéaire entre les paramètres.} Il s'ensuit naturellement que $r$ points généraux de la courbe définissent un groupe; mais nous verrons tout à l'heure comment on peut aussi (avec quelques restrictions) tirer les prémisses de la conclusion.

Lorsqu'une série linéaire $g_{n}^{r}$ $(r \geqq 0)$ est donnée sur une courbe, il peut se faire qu'il existe sur la même courbe une nouvelle série $g_{n}^{s}$ du même ordre et de dimension $s > r$, qui contienne tous les groupes de $g_{n}^{r}$, et par suite la $g_{n}^{r}$ elle-même; il peut arriver au contraire qu'une telle série $g_{n}^{s}$ n'existe pas. Or on dit \emph{qu'une série est complète lorsqu'elle n'est pas contenue dans une autre série du même ordre et de dimension supérieure;} la série est \emph{incomplète} dans le cas contraire${}^{*)}$.

Dans ce dernier cas, en formant des séries de l'ordre $n$ et de plus en plus amples, qui contiennent une série incomplète $g_{n}^{r}$ donnée, on parviendra nécessairement à une série $g_{n}^{s}$ $(r < s \leqq n)$ complète contenant celle-là. Or on démontre que cette dernière série ne dépend pas du procédé qu'on suit pour la construire; elle dépend seulement de la série $g_{n}^{r}$ d'où l'on part. En d'autres termes:

\emph{Une série linéaire $g_{n}^{r}$ $(r \geqq 0)$ ou bien est complète, ou bien elle est contenue dans une série complète du même ordre $n$, laquelle est entièrement déterminée par la $g_{n}^{r}$}${}^{**)}$.

Il s'ensuit par exemple que \emph{si d'une série complète $g_{n}^{r}$ on déduit une nouvelle série $g_{n-\nu}^{\varrho}$ en imposant aux groupes de celle-là la condition de contenir $\nu$ points quelconques de la courbe, et en les supprimant ensuite dans ces groupes, la nouvelle série sera encore complète.} On l'appelle \emph{série résiduelle} des $\nu$ points par rapport à $g_{n}^{r}$; sa dimension $\varrho$ est égale à $r - \nu$, ou supérieure, selon que les $\nu$ points présentent des conditions toutes indépendantes, ou non, aux groupes de $g_{n}^{r}$ qui les contiennent.

Ce sont les séries complètes qui offrent le plus grand intérêt. Toutefois la théorie des surfaces va nous amener souvent à considérer aussi des séries incomplètes. Un caractère d'une telle série $g_{n}^{r}$ auquel il faut avoir égard, c'est la difference $s - r$ que l'on obtient en

\textbf{*)} On voit bien que cette définition est parfaitement d'accord avec celle que nous avons donnée, en partant de considérations analytiques.

\textbf{**)} \emph{Segre}, l.\ c., §~14.

\origpage{253}
retranchant la dimension de notre série de la dimension de la série complète $g_{n}^{s}$ qui contient celle-là. Nous appelons cette différence le \emph{défaut (deficienza, difetto di completezza, Defect)} de la $g_{n}^{r}$.

\emph{Remarque.} --- Sur la définition d'une série linéaire on peut faire une remarque, qui, du reste, n'a pas d'importance pour la suite de notre Monographie. Appelons \emph{involution d'ordre $n$ et dimension $r$} sur une courbe algébrique tout système algébrique de $\infty^{r}$ groupes de $n$ points, tels que $r$ points généraux de la courbe entrent dans un seul groupe du système. Alors nous pouvons dire que toute série linéaire est une involution; sera-t-il vrai aussi le théorème réciproque? Pas toujours, car il y a des involutions $\infty^{1}$ qui ne sont pas linéaires; celles, par exemple, qui sont découpées par les génératrices d'une surface réglée irrationnelle sur les courbes appartenantes à la surface. Mais, \emph{en dehors} des involutions $\infty^{1}$, et \emph{des involutions $\infty^{r}$ dont tout groupe est formé par $r$ groupes arbitraires d'une involution $\infty^{1}$} (d'ordre $n \geqq 1$), on démontre que \emph{toute involution de dimension $> 1$ est une série linéaire}${}^{*)}$. On a donc là une nouvelle définition de ces séries.

\subsection*{7. Système linéaire de courbes sur une surface.}

Les fonctions rationnelles d'une équation algébrique à deux variables nous ont portés à considérer les séries linéaires sur une courbe; de même, en partant d'une équation à trois variables, on parviendrait à envisager les systèmes linéaires de courbes sur une surface.

Pour les définir au point de vue géométrique il suffit de considérer une surface algébrique $F$ (de l'espace ordinaire ou d'un hyperespace) et un système linéaire de variétés (à deux, ..., dimensions)
\[
\lambda_{1}f_{1} + \cdots + \lambda_{r+1}f_{r+1} = 0;
\]
les courbes que ces variétés découpent sur la $F$ forment un \emph{système linéaire} dont la \emph{dimension} est $r$, si aucune des variétés nommées ne contient la $F$.

Au sujet des variétés $f$, il peut se faire qu'elles contiennent une même courbe fixe de $F$, courbe qu'on est libre de regarder comme faisant partie de la courbe générale du système, ou bien comme étrangère à celle-ci (qui alors se réduira à la partie variable de l'intersection $fF$). Il peut arriver encore que les $f$ passent par des points fixes de $F$ (en dehors des courbes fixes); ces points, en nombre fini, qui appartiennent à toutes les courbes du système, sont appelés \emph{points-base} de celui-ci.

\textbf{*)} Ce théorème a été donné en même temps (juin 1893) par MM.\ \emph{Humbert} (Comptes Rendus de l'Ac.\ d.\ Sc.; voir aussi Journal de Mathématiques, 4${}^{\text{e}}$ série t.\ X) et \emph{Castelnuovo} (Atti dell'Accad.\ d.\ Sc.\ di Torino).

\origpage{254}

Un système linéaire $\infty^{r}$ de courbes sur une surface a la propriété que par $r$ points généraux de la surface passe une seule courbe du système. Cette propriété peut même servir de définition pour un système linéaire de dimension $r > 1$. Il existe, c'est vrai, des systèmes pas linéaires $\infty^{1}$ de courbes satisfaisant à la condition que par tout point général de la surface passe une seule courbe du système; \emph{ce sont les faisceaux irrationnels}, dont le système des génératrices d'une surface réglée irrationnelle fournit un exemple bien simple. \emph{Mais tout système $\infty^{r}$, avec $r > 1$, dont une seule courbe est entièrement déterminée par $r$ points de la surface, est un système linéaire, si l'on exclut seulement le cas où la courbe générale du système est composée par plusieurs courbes d'un faisceau irrationnel}${}^{*)}$.

La courbe générale d'un système linéaire peut être \emph{irréductible}, ou bien se composer de plusieurs parties; dans le premier cas on dit que le système est \emph{irréductible}, et \emph{réductible} dans le second cas. L'étude de ce dernier cas se réduit au premier (le plus intéressant) à l'aide du théorème:

\emph{La courbe générale d'un système linéaire réductible $\infty^{r}$ est formée:}

1) \emph{ou bien par une partie fixe (commune à toutes les courbes du système), jointe à une courbe qui varie dans un système linéaire irréductible $\infty^{r}$;}

2) \emph{ou bien par plusieurs $(\geqq r)$ courbes qui appartiennent à un même faisceau (rationnel ou non), auxquelles peuvent s'ajouter parfois des parties fixes.}

L'exclusion des systèmes réductibles, qui rendrait plus simples nos considérations, n'est pas toujours possible; car il y a des opérations sur les systèmes linéaires qui, même en partant d'un système irréductible, peuvent porter à des systèmes réductibles. Mais, à l'aide de conventions appropriées, on peut réunir sous un même énoncé la plupart des propriétés des uns et des autres systèmes.

Bornons-nous, pour le moment, aux systèmes irréductibles; soit $|C|$ un tel système, $C$ sa courbe générale.

Les caractères les plus importants du système (qui restent invariables par les transformations birationnelles de la surface) sont les suivants:

1) la \emph{dimension} $r$ de $|C|$;

2) le \emph{degré} $n$, c'est à dire le nombre des intersections variables de deux courbes générales du système (pour $r = 1$, on a $n = 0$).

3) le \emph{genre} $\pi$, qui est le genre de la courbe $C$.

On a aussi à considérer souvent une série linéaire qui est liée invariablement avec le système $|C|$, et qui va jouer dans la suite un

\textbf{*)} \emph{Enriques}, \emph{Una questione sulla linearità}..., Rendic.\ della R.\ Accad.\ d.\ Lincei, luglio 1893; la démonstration est rapportée avec quelques simplifications dans la Monographie citée de M.\ \emph{Segre}, n${}^{\text{o}}$ 27.

\origpage{255}
rôle important; c'est la \emph{série caractéristique $g_{n}^{r-1}$}, qui est découpée sur une $C$ par les autres courbes de $|C|$${}^{*)}$. Ici l'on a $r - 1 \leqq n$, d'où l'on conclut que la dimension d'un système irréductible ne peut surpasser le degré augmenté d'une unité.

\emph{Remarque.} --- La propriété d'un système $|C|$, situé sur une surface $F$, d'être réductible ou irréductible, ainsi que les caractères $r$, $n$, $\pi$ de $|C|$, se transportent sans modifications à un système $|C'|$ d'une surface $F'$, lorsqu'il y a une transformation birationnelle qui change $F$ en $F'$ et $|C|$ en $|C'|$. Toutefois pour que cette remarque subsiste sans exceptions, il faut faire la convention suivante. Il peut arriver qu'un point-base $P$ de $|C|$ sur $F$, ayant l'ordre $i \geqq 1$ pour la courbe $C$ générale (point que nous supposerons simple pour $F$), donne, par la transformation, une courbe exceptionnelle $P'$ de $F'$. D'après le n${}^{\text{o}}$ 2, on devrait regarder cette courbe $P'$, comme faisant partie de la courbe $C'$ transformée de $C$; mais alors $|C'|$ serait réductible, lors même que $|C|$ ne l'est pas. Pour éviter le désaccord on convient dans ce cas, de faire abstraction de la courbe $P'$ (correspondante à un point-base de $|C|$), et de regarder $C'$ comme le lieu des points correspondants aux points de $C$, excepté le point $P$. On va respecter au contraire la règle du n${}^{\text{o}}$ 2 lorsque $P$ (qui se transforme en la courbe $P'$) n'est pas un point-base de $|C|$, mais bien un point général de $F$.

Outre l'exemple que nous avons porté, il y en aurait d'autres analogues à discuter. Nous n'y insistons pas ici; le lecteur qui voudra approfondir ces questions, pourra le faire ailleurs${}^{**)}$; il verra pourtant de ces quelques mots, qu'il y a dans la théorie des surfaces des difficultés et des distinctions qu'on ne réussit pas à éviter.

\subsection*{8. Systèmes linéaires complets.}

Si nous nous proposons de transporter aux systèmes linéaires sur une surface la notion de série complète établie pour une courbe, nous rencontrons tout d'abord une distinction qui ne se présentait pas alors. En effet si l'on cherche un système $|D|$ qui contienne le système donné $|C|$, il faut fixer quel est le caractère de $|C|$ qui doit se transmettre au nouveau système. On peut exiger que le genre $\pi$ seulement reste

\textbf{*)} L'importance de cette série dans l'étude des systèmes linéaires de courbes planes a été remarquée par MM.\ \emph{Segre} (\emph{Sui sistemi lineari}..., Rendic.\ Circolo Mat.\ di Palermo, tomo I, 1887) et \emph{Castelnuovo} (\emph{Ricerche generali sopra i sistemi lineari}..., Memorie dell'Acc.\ d.\ Sc.\ di Torino, s${}^{\text{e}}$ II, t.\ 42). Quant à la considération de cette série sur les surfaces on n'a qu'à consulter les travaux de MM.\ \emph{Enriques} et \emph{Castelnuovo}, cités dans l'introduction.

\textbf{**)} \emph{Enriques}, \emph{Introduzione}..., cap.\ I.

\origpage{256}
invariable; mais on peut aussi exiger que ce soit le degré $n$ (et par suite le genre) qui ne change pas. Bien que la première condition soit de nature plus générale que la seconde, c'est à celle-ci qu'il faut assigner plus d'intérêt; le progrès de la théorie des surfaces l'a montré. Par conséquent nous dirons qu'\emph{un système linéaire irréductible de courbes est complet (par rapport au degré), lorsqu'il n'y a sur la surface d'autre système du même degré et de dimension supérieure, qui contienne le premier système}${}^{*)}$.

En voici un exemple fort simple. Considérons sur un plan le système linéaire de toutes les courbes ayant des multiplicités données en certains points-base; c'est un système \emph{complet}. En d'autres termes: \emph{un système linéaire de courbes planes entièrement défini par ses points-base est complet}${}^{**)}$.

La définition que nous venons de donner ci-dessus, peut se présenter d'une manière plus générale, qui permet de l'étendre aussi aux systèmes réductibles. A cet effet il faut expliquer d'abord la locution: \emph{une courbe $C_{0}$ (simple ou composée) est contenue totalement dans un système $|C|$ donné.}

Cela revient à dire que la courbe $C_{0}$ est par elle-même \emph{une} courbe de $|C|$ (on exclut ainsi qu'elle soit \emph{une partie} d'une telle courbe), et en outre que $C_{0}$ a le même ordre de multiplicité de la courbe générale de $|C|$ en tout point-base du système. Pareillement on dit qu'un système est contenu \emph{totalement} dans un autre, lorsque la courbe générale du premier système est contenue totalement dans le second, et les points-base de celui-ci sont aussi les points-base de celui-là. On démontre alors facilement que deux systèmes dont l'un est contenu totalement dans l'autre, ont le même \emph{degré}; et \emph{viceversa}, si de deux systèmes ayant le même degré, l'un est contenu dans l'autre, on peut affirmer qu'il y est contenu \emph{totalement}${}^{***)}$. Il résulte que la définition de système complet peut maintenant s'énoncer comme il suit:

\emph{Un système linéaire est complet s'il n'est pas contenu totalement dans un autre système.}

Sous cette forme la définition a une signification, et montre ce

\textbf{*)} La notion de système complet (par rapport au degré), \emph{sistema normale} a été introduite par M.\ \emph{Enriques} (\emph{Ricerche di geometria}... I, 2; \emph{Introduzione}...): l'extension aux variétés $\infty^{r}$ se trouve dans la Monographie citée de M.\ \emph{Segre} (n${}^{\text{o}}$ 26). M.\ \emph{Enriques} dans le premier Mémoire parle aussi des systèmes complets par rapport au genre, qu'il appelle \emph{sistemi completi}. Nous n'y recourons pas dans ce Mémoire.

\textbf{**)} Si parmi les points-base il y en a des simples, en faisant abstraction de ceux-ci l'on obtient un système plus ample que le système donné, qui a pourtant le même genre de celui-ci, mais le degré supérieur. On voit donc qu'un système complet par rapport au degré peut bien n'être pas complet par rapport au genre.

\textbf{***)} \emph{Enriques}, \emph{Introduzione}... n${}^{\text{o}}$ 7.

\origpage{257}
que l'on doit entendre par \emph{système complet}, lors même que le système donné est réductible, ou s'il a la dimension \emph{un} ou \emph{zéro}; dans ces derniers cas pourtant on doit désigner quels sont les points qu'il faut regarder comme points-base du système complet.

C'est dans ce sens qu'il faut entendre le théorème suivant, qui est parfaitement analogue à un théorème sur les séries linéaires${}^{*)}$:

\emph{Tout système linéaire $\infty^{r}$ $(r \geqq 0)$ est complet, ou bien il est contenu totalement dans un système complet qui est entièrement déterminé par celui-là.}

On a encore cet autre théorème qui rappelle aussi un théorème sur les séries:

\emph{Si d'un système linéaire complet on déduit un nouveau système en imposant aux courbes du premier la condition de passer par des points fixes avec des ordres de multiplicité donnés, le dernier système est aussi complet.}

Dans la suite nous allons considérer, autant que possible, des systèmes complets.

\subsection*{9. Opérations élémentaires sur les séries et sur les systèmes de courbes.}

Nous appelons par ce nom les opérations d'addition et de soustraction des séries et des systèmes, que nous allons définir; elles correspondent à la multiplication et à la division des fonctions rationnelles qui appartiennent à un même champ de rationnalité.

Soient d'abord deux séries linéaires complètes $g_{m}^{q}$, $g_{n}^{r}$ sur une même courbe; formons tous les groupes de $m + n = p$ points que l'on obtient en réunissant un groupe arbitraire de $g_{m}^{q}$ à un groupe arbitraire de $g_{n}^{r}$; les groupes ainsi formés appartiennent, on le démontre facilement, à une même série complète $g_{p}^{s}$, que l'on appelle \emph{somme des deux séries données}. On n'exclut pas le cas où les deux séries $g_{m}^{q}$, $g_{n}^{r}$ coïncident; on parvient alors à la série \emph{double} d'une série donnée. Et de même l'on définit la série somme de plusieurs séries, et la série multiple (k.\ uple) d'une série donnée${}^{**)}$.

Nous sommes maintenant en mesure de définir l'opération inverse, la soustraction des séries, autant qu'elle est possible. Considérons à

\textbf{*)} \emph{Enriques}, \emph{Introduzione}... n${}^{\text{o}}$ 11.

\textbf{**)} \emph{Segre}, \emph{Introduzione}... n${}^{\text{o}}$ 59. Vraiment on entendait quelques fois par \emph{somme de deux séries} $g_{m}^{q}$, $g_{n}^{r}$ (complètes, ou non) la série d'ordre $m + n$ et de la plus petite dimension, qui contient les groupes composés d'un groupe de $g_{m}^{q}$ et d'un groupe de $g_{n}^{r}$ (\emph{voir} \emph{Castelnuovo}, \emph{Sui multipli di una serie lineare}..., Rendic.\ Circolo Mat.\ di Palermo, t.\ VII). Mais dans cette Monographie nous adoptons la définition exposée ci-dessus.

\origpage{258}
cet effet deux séries complètes $g_{p}^{s}$, $g_{n}^{r}$ $(p > n)$ sur la courbe, et supposons qu'un certain groupe $G_{n}$ de la seconde série soit contenu dans $\infty^{q}$ groupes $(q \geqq 0)$ de la première. Les groupes des $p - n = m$ points qui restent, forment comme nous savons, une série complète $g_{m}^{q}$ \emph{résiduelle} de $G_{n}$ par rapport à $g_{p}^{s}$. Or cette série $g_{m}^{q}$ a une propriété remarquable; en effet elle est résiduelle (par rapport à $g_{p}^{s}$) non seulement de $G_{n}$, mais aussi de tout autre groupe de $g_{n}^{r}$. On dit par conséquent que $g_{m}^{q}$ \emph{est la série résiduelle de} $g_{n}^{r}$ \emph{par rapport à} $g_{p}^{s}$; mais l'on voit de même que $g_{n}^{r}$, à son tour, est résiduelle de $g_{m}^{q}$ par rapport à $g_{p}^{s}$; la dernière série est la somme des deux premières. Cette propriété donne lieu à l'énoncé qui suit${}^{*)}$:

\emph{Deux séries complètes $g_{p}^{s}$, $g_{n}^{r}$ étant données sur une même courbe, si un groupe de la $g_{n}^{r}$ est contenu dans quelque groupe de la prémière série, il en est de même pour tout autre groupe de la $g_{n}^{r}$; il existe alors une troisième série $g_{m}^{q}$ entièrement déterminée, qui additionnée à la seconde $g_{n}^{r}$ donne la première série $g_{p}^{s}$; les deux séries $g_{m}^{q}$, $g_{n}^{r}$ sont résiduelles l'une de l'autre par rapport à $g_{p}^{s}$.}

A ce théorème l'on pourrait donner le nom de \emph{théorème du reste} pour les séries; ce nom cependant a été dejà adopté par MM.\ \emph{Brill} et \emph{Nöther} pour désigner un théorème qui a bien d'analogie avec le précédent, mais qui comprend aussi d'autres notions dont nous allons parler plus tard.

Venons maintenant aux systèmes linéaires de courbes sur une surface; l'analogie nous permet d'étendre sur le champ les résultats qui précèdent. Si $|C_{1}|$ et $|C_{2}|$ sont deux systèmes complets de courbes sur une même surface, les courbes $C_{1} + C_{2}$, que l'on peut former en réunissant deux courbes prises arbitrairement dans les deux systèmes, appartiennent à un même système complet déterminé $|C| = |C_{1} + C_{2}|$ qui s'appelle \emph{somme de} $|C_{1}|$ \emph{et} $|C_{2}|$. En particulier on peut parler du système \emph{double, triple}... d'un système donné.

Partons à présent du système $|C|$, et soit $C_{1}$ une courbe qui entre dans quelque courbe (réductible) de $|C|$; les courbes $C_{2}$ qui avec la $C_{1}$ donnent lieu à des courbes de $|C|$ forment un système complet $|C_{2}|$ \emph{résiduel de la} $C_{1}$ \emph{par rapport à} $|C|$. Or si la $C_{1}$ appartient à un système complet $|C_{1}|$ (déterminé par elle), on démontre que $|C_{2}|$ est aussi résiduel (par rapport à $|C|$) de toute autre courbe de $|C_{1}|$; c'est \emph{le système} $|C_{2}| = |C - C_{1}|$ \emph{résiduel de} $|C_{1}|$ \emph{par rapport à} $|C|$, de même que $|C_{1}|$ est résiduel de $|C_{2}|$ par rapport à $|C|$. On a donc le théorème${}^{**)}$:

\emph{Si un système complet $|C|$ contient une courbe d'un second système}

\textbf{*)} \emph{Segre}, \emph{Introduzione}... n${}^{\text{o}}$ 58.

\textbf{**)} \emph{Enriques}, \emph{Introduzione}... n${}^{\text{o}}$ 13.

\origpage{259}
\emph{complet $|C_{1}|$, le premier système contient toute courbe du second; il existe alors un système complet $|C_{2}|$ entièrement déterminé, tel que $|C| = |C_{1} + C_{2}|$; chacun des deux systémes $|C_{1}|$ et $|C_{2}|$ est résiduel de l'autre par rapport à $|C|$.}

Il faut remarquer pourtant que même si $|C|$ et $|C_{1}|$ sont supposés irréductibles, il peut bien se faire que $|C_{2}| = |C - C_{1}|$ soit réductible.

Pour terminer ces considérations sur le système somme de deux systèmes donnés, nous ferons connaître ici deux relations bien simples, auxquelles il faut pourtant recourir souvent dans les recherches sur les surfaces. Ces formules expriment le degré $n$ et le genre $\pi$ du système somme $|C| = |C_{1} + C_{2}|$ à l'aide des caractères analogues $n_{1}$, $\pi_{1}$ et $n_{2}$, $\pi_{2}$ des deux systèmes $|C_{1}|$ et $|C_{2}|$; on a précisément
\[
\begin{aligned}
n &= n_{1} + n_{2} + 2i,\\
\pi &= \pi_{1} + \pi_{2} + i - 1,
\end{aligned}
\]
où $i$ désigne le nombre des intersections variables d'une $C_{1}$ avec une $C_{2}$. On suppose ici que $|C_{1}|$ et $|C_{2}|$ sont irréductibles; mais il est utile de remarquer que s'il n'en était pas ainsi, on pourrait de même définir le genre et le degré des systèmes réductibles de sorte que les relations précédentes fussent encore vraies${}^{*)}$.

\section*{Chapitre III. Courbes adjointes à une courbe plane. --- Surfaces sous-adjointes à une surface de l'espace ordinaire.}

Dans le Chapitre précédent nous avons consideré des séries complètes et des systèmes complets de courbes, mais nous n'avons pas indiqué le moyen de \emph{completer} une série ou un système donné. C'est ce que nous allons faire maintenant, en nous bornant (dans ce Chapitre) aux séries situées sur les courbes \emph{planes}, et aux systèmes de courbes sur les surfaces de \emph{l'espace ordinaire}. Tous les autres cas peuvent se ramener à ceux-ci, à l'aide d'une transformation birationnelle convenable de la courbe ou de la surface donnée.

\subsection*{10. Courbes adjointes à une courbe plane.}

Soit d'abord $C$ une courbe \emph{plane} de l'ordre $m$, qui aura en général des points multiples. On appelle \emph{courbe adjointe} à $C$ une courbe qui passe $\alpha - 1$ fois par tout point multiple d'ordre $\alpha$ de la courbe $C$ $(\alpha = 2, 3 \ldots)$${}^{**)}$. Les courbes adjointes à $C$ d'un ordre donné (qui

\textbf{*)} \emph{Enriques}, \emph{Introduzione}... n.\ 15, 16.

\textbf{**)} \emph{Brill} et \emph{Nöther}, \emph{Ueber die algebraischen Functionen}..., l.\ c.

\origpage{260}
existent certainement si cet ordre est assez élevé) forment un système linéaire (complet) de courbes planes. Or ce système a une propriété bien remarquable, d'où s'ensuit l'application des courbes adjointes à la théorie des séries linéaires. En effet${}^{*)}$:

\emph{Les courbes adjointes à $C$ d'un ordre quelconque donné découpent sur la courbe $C$ une série linéaire complète;} il faut entendre ici que tout groupe de la série est formé par les intersections de $C$ avec une courbe adjointe en dehors des points multiples. Il s'ensuit (n${}^{\text{o}}$ 9) qu'on obtient de même sur $C$ une série complète, si l'on oblige les courbes adjointes à passer par des points fixés arbitrairement \emph{sur} $C$, et l'on se borne à considérer les intersections qui tombent en dehors de ces points (et des points multiples).

A l'aide de la dernière observation on peut tout de suite \emph{construire sur la courbe $C$, la série linéaire complète qui est définie par un groupe de points, ou par une série incomplète donnée}. Par le groupe $G$ donné (qui dans le second cas sera un groupe quelconque de la série) on fait passer une courbe adjointe quelconque, ce qui est toujours possible si l'ordre de celle-ci est assez élevé. La courbe découpe sur $C$, en dehors des points multiples et du groupe $G$, un groupe de points $G'$. Conduisons par $G'$ toutes les courbes adjointes de l'ordre choisi; celles-ci découperont sur $C$ la série complète définie par $G$.

Il y a de l'arbitraire dans cette construction, car on dispose de l'ordre des courbes adjointes, et aussi de la courbe, parmi celles-ci, que l'on conduit d'abord par $G$; mais de quelque manière que l'on s'arrange, toujours on parviendra enfin à la même série. C'est précisément cette affirmation qui constitue le \emph{Restsatz} de MM.\ \emph{Brill} et \emph{Nöther}${}^{**)}$.

\subsection*{11. Série caractéristique d'un système linéaire de courbes planes.}

La propriété de couper sur une courbe plane donnée $C$ une série complète n'appartient pas seulement aux systèmes de courbes adjointes,

\textbf{*)} \emph{Brill} et \emph{Nöther}, l.\ c.\ pag.\ 275; voir aussi \emph{Segre}, \emph{Introduzione}... n${}^{\text{o}}$ 77.

\textbf{**)} Le \emph{Restsatz} provient donc de deux propositions bien distinctes, dont l'une appartient proprement à la \emph{Géométrie sur la courbe}, et se rapporte aux opérations d'addition et de soustraction des séries (n${}^{\text{o}}$ 9); tandis que l'autre proposition appartient à la \emph{Géometrie du plan}, et exprime la propriété des systèmes de courbes adjointes que nous avons énoncée ci-dessus. On pourrait même remarquer que le \emph{Restsatz} (comme il est énoncé par MM.\ \emph{Brill} et \emph{Nöther}) dit un peu moins que la réunion de nos deux propositions, car il se borne à affirmer que la série découpée par les courbes adjointes de la manière indiquée, ne dépend pas de celles-ci, tandis que nous affirmons en outre que la même série est complète; mais il faut d'ailleurs remarquer que cette dernière partie se déduit sur le champ de la première. Voir \emph{Brill} et \emph{Nöther}, l.\ c.\ pag.\ 273; \emph{Segre}, \emph{Introduzione}..., n${}^{\text{o}}$ 78.

\origpage{261}
mais aussi à d'autres systèmes. Par exemple les courbes d'un ordre arbitraire, qui passent avec la multiplicité $\alpha$ par chaque point multiple d'ordre $\alpha$ de $C$, découpent elles-mêmes sur $C$ une série complète; c'est là une propriété qui s'ensuit aisément de la propriété des courbes adjointes. Le cas le plus intéressant se présente si l'ordre des courbes que nous considérons, est l'ordre de $C$; en effet dans ce cas la série nommée n'est autre chose que la \emph{série caractéristique} du système complet des courbes planes ayant l'ordre de $C$ et les mêmes multiplicités dans les points multiples de celle-ci. Nous pouvons donc énoncer le théorème suivant qui va nous être utile${}^{*)}$:

\emph{La série caractéristique d'un système linéaire complet de courbes planes est toujours complète.}

\subsection*{12. Surfaces sous-adjointes à une surface de l'espace ordinaire.}

Nous allons voir maintenant comment les théorèmes qui précèdent peuvent s'étendre aux surfaces de l'espace ordinaire.

Soit $F$ une telle surface ayant des singularités (courbes multiples et points multiples) quelconques. On peut toujours construire une surface $\Psi$ d'un ordre assez élevé, qui remplit la condition de couper sur tout plan général une courbe adjointe à la section de $F$ faite par le même plan. On précise, de cette manière, le comportement de la $\Psi$ le long des courbes multiples de $F$; ainsi lorsqu'il s'agit de courbes multiples ordinaires, la $\Psi$ va passer $\alpha - 1$ fois par toute courbe de $F$ qui a la multiplicité $\alpha$ pour $F$. On ne dit rien au contraire sur le passage de $\Psi$ par les points multiples \emph{isolés} de $F$, dont la présence sur $F$ ne s'ensuit pas (pour ainsi dire) de l'existence des courbes multiples; de sorte que si $F$ possède, par exemple, un point multiple qui n'appartient à aucune courbe multiple, la $\Psi$ ne passera pas, en général, par ce point. En vertu de ces explications, nous dirons que la $\Psi$ est une surface \emph{adjointe à $F$ le long des courbes multiples}, ou, tout court, que la $\Psi$ est \emph{sous-adjointe (subaggiunta) à $F$}${}^{**)}$. Nous réservons le nom de surfaces \emph{adjointes} à certaines surfaces, dont nous aurons à nous occuper plus tard, qui tout en satisfaisant aux conditions de la $\Psi$, vont encore passer convenablement par les points isolés de $F$. Pour le moment les surfaces sous-adjointes suffisent à notre but.

\textbf{*)} \emph{Brill} et \emph{Nöther}, l.\ c.\ pag.\ 308; \emph{Castelnuovo}, \emph{Ricerche generali sopra i sistemi lineari}... l.\ c.\ n${}^{\text{o}}$ 18.

\textbf{**)} \emph{Enriques}, \emph{Introduzione}... n${}^{\text{o}}$ 17; M.\ \emph{Nöther} (\emph{Zur Theorie}..., Math.\ Annalen 8, pag.\ 509) appelle \emph{adjungirte Flächen} les surfaces que nous appelons \emph{sous-adjointes}, et réserve le nom de \emph{surfaces} $\varphi$ aux surfaces d'ordre $n - 4$ qui, d'après cette Monographie, sont \emph{adjointes à $F$}. Le développement de la théorie des surfaces a conseillé ce changement de locutions.

\origpage{262}

Les surfaces d'un ordre donné sous-adjointes à $F$ forment un système linéaire, qui jouit d'une propriété analogue à celle des courbes adjointes. En effet${}^{*)}$:

\emph{Les surfaces d'un ordre quelconque donné qui sont sous-adjointes à une surface, découpent sur celle-ci un système de courbes complet:} la courbe générale du système est formée par l'intersection de la surface sous-adjointe avec $F$, en dehors des courbes multiples de $F$. Le théorème subsiste encore si l'on oblige les surfaces sous-adjointes à passer par des points ou des courbes assignées sur $F$, et l'on retranche ces courbes fixes des courbes du système (n${}^{\text{o}}$ 9).

A l'aide de ces remarques on peut toujours construire sur $F$ le système complet déterminé par une courbe donnée $C$, ou par un système linéaire incomplet. Il suffit en effet de conduire par la courbe $C$ (qui dans le second cas sera choisie arbitrairement dans le système) une surface sous-adjointe $\Psi$ d'un ordre assez élevé pour que la construction soit possible. La $\Psi$ va découper en général sur $F$ une nouvelle courbe $C'$, en dehors de $C$ et des courbes multiples de $F$; que l'on conduise maintenant par $C'$ toutes les surfaces adjointes du même ordre de $\Psi$; celles-ci vont découper sur $F$ un système linéaire complet auquel appartient la courbe $C$; c'est le système cherché, ou tout au plus, il en diffère seulement en cela, que le système cherché peut avoir quelques points-base en plus. Il suffit, dans ce dernier cas, d'imposer ces nouveaux points-base aux courbes du système que nous venons de construire, ou (ce qui revient au même) d'imposer aux surfaces $\Psi$ des passages ou des contacts convenables avec $F$ dans les mêmes points.

Le système complet auquel on parvient ne dépend ni de l'ordre des surfaces sous-adjointes $\Psi$ que l'on emploie, ni de la surface, parmi les $\Psi$, qu'on conduit d'abord. Cette affirmation constitue le \emph{Restsatz} pour les surfaces, d'après M.\ \emph{Nöther}${}^{**)}$.

\emph{Remarque.} --- Le théorème sur la série caractéristique d'un système linéaire de courbes planes peut s'étendre aussi aux systèmes linéaires de surfaces de l'espace ordinaire. Enonçons le résultat, bien qu'on n'ait pas dans la suite l'occasion d'y recourir:

\emph{Si un système linéaire de surfaces est entièrement déterminé par les courbes et les points-base, le système de courbes, qui est découpé sur une des surfaces par les autres surfaces du système, est complet.}

\textbf{*)} \emph{Enriques}, \emph{Introduzione}... n${}^{\text{o}}$ 35 (vraiment on y considère les surfaces adjointes, mais la chose est vraie aussi pour les surfaces sous-adjointes). Le même théorème s'ensuit aussi du \emph{Restsatz} de M.\ \emph{Nöther} dont nous allons parler tout à l'heure.

\textbf{**)} \emph{Zur Theorie}..., Math.\ Annalen 8, pag.\ 509; il y aurait à répéter ici une remarque analogue à celle, que nous avons faite à propos des courbes planes.

\origpage{263}

\section*{Chapitre IV. Système adjoint à un système donné.}

Nous allons définir dans ce Chapitre une opération qui va jouer un rôle fondamental dans la théorie des systèmes linéaires situés sur une surface; c'est l'opération \emph{d'adjonction}. Elle fournit le moyen de déduire de tout système linéaire donné, un nouveau système, \emph{l'adjoint} au premier; et par la considération des liens qui passent entre ces deux systèmes, elle nous porte à établir les principaux invariants d'une surface par rapport aux transformations birationnelles.

Suivant notre méthode nous commençons par des considérations relatives aux séries linéaires sur les courbes.

\subsection*{13. Série canonique sur une courbe.}

Parmi les courbes adjointes à une courbe plane $C$ d'ordre $m$, il faut considérer en particulier celles qui ont l'ordre $m - 3$. Supposons qu'il existe de telles courbes, et désignons par $p(>0)$ le nombre de celles qui sont linéairement indépendantes. Alors les dites courbes découpent sur la $C$ une série linéaire complète dont la dimension est $p - 1$, et l'ordre est (d'après MM.\ \emph{Brill} et \emph{Nöther}${}^{*)}$) $2p - 2$; c'est une $g_{2p-2}^{p-1}$. Or cette série jouit d'une propriété fondamentale. Si en effet on transforme birationnellement la courbe $C$ en une autre courbe plane $C'$ d'ordre $m'$, la série $g_{2p-2}^{p-1}$ de $C$ se transforme dans la série qui est découpée à son tour sur $C'$ par les courbes d'ordre $m' - 3$ adjointes à $C'$. En d'autres termes \emph{la série découpée sur une courbe plane d'ordre $m$ par les courbes adjointes d'ordre $m - 3$ jouit de la propriété d'invariance} (par rapport aux transformations birationnelles). Nous aurons souvent l'occasion de nommer cette série; nous la désignerons dans la suite par le nom de \emph{série canonique} sur la courbe $C$ (ou sur toute autre courbe de la même \emph{classe}). On parlera donc aussi de la série canonique sur une courbe gauche; ce sera la série correspondante à celle découpée par les adjointes d'ordre $m - 3$, sur une courbe plane d'ordre $m$ en correspondance birationnelle avec la courbe donnée.

La dimension de la série canonique augmentée d'une unité, est aussi un invariant de la courbe considérée; c'est le \emph{genre} $p$ de la courbe, dont nous avons dejà rappelé la définition dans un cas particulier (n${}^{\text{o}}$ 4).

\textbf{*)} \emph{Ueber die algebraischen Functionen}... pag.\ 280; voir aussi \emph{Segre}, \emph{Introduzione}..., n${}^{\text{o}}$ 73.

\origpage{264}

Le caractère $p$, en vertu de sa définition, ne saurait descendre au-dessous de zéro; il est nul si la courbe correspondante, supposée plane, d'ordre $m$, n'a pas de courbes adjointes d'ordre $m - 3$. On sait bien d'ailleurs que la condition $p = 0$ définit entièrement une \emph{classe} de courbes; ce sont les \emph{courbes rationnelles}.

Ce cas excepté, on connaît des propriétés de la série canonique qu'il convient de rappeler ici, car elles sont d'un usage continuel${}^{*)}$. La série canonique est la seule série qui ait, sur une courbe de genre $p$, la dimension $p - 1$ avec l'ordre $2p - 2$. La même série n'a pas de points fixes qui soient communs à tous ses groupes. Par suite si l'on fixe un point quelconque de la courbe, on déduit de la série canonique $g_{2p-2}^{p-1}$ une nouvelle série $g_{2p-3}^{p-2}$; celle-ci à son tour n'a pas \emph{en général} de points fixes. Tout au plus, dans un cas particulier, elle peut avoir un point fixe; il arrive alors que tous les groupes de la série canonique qui contiennent un point fixe quelconque, contiennent en conséquence un second point fixe qui est déterminé par le premier; alors il existe sur la courbe une série $g_{2}^{1}$, telle que $p - 1$ quelconques de ses groupes forment un groupe de la série canonique $g_{2p-2}^{p-1}$; c'est le cas des courbes \emph{hyperelliptiques}.

\emph{Remarque.} --- De la série canonique on peut donner une définition analytique, bien connue, qui a l'avantage de s'étendre à la fois à toutes les courbes d'une même classe, qu'elles soient planes ou gauches: \emph{les groupes de la série canonique sont les groupes des zéros des différentielles abéliennes de la première espèce, qui appartiennent à la courbe donnée.}

\subsection*{14. Système adjoint à un système de courbes planes.}

Revenons à la courbe plane $C$ d'ordre $m$, et supposons maintenant qu'elle soit une courbe générale d'un système linéaire $|C|$ de courbes planes, système complet, c'est-à-dire déterminé entièrement par ses points-base. Ces points ont les mêmes ordres de multiplicité $(\geqq 1)$ pour la $C$ et pour les autres courbes du système; et tous les points multiples de la $C$ tombent parmi eux. Il s'ensuit que les courbes adjointes à $C$, de l'ordre $m - 3$, (et plus généralement d'un ordre quelconque) sont de même adjointes à toute autre courbe du système $|C|$. Nous dirons que ces courbes de l'ordre $m - 3$ forment \emph{le système} $|C'|$ \emph{adjoint au système} $|C|$; on dit \emph{adjoint}, tout court, sans désigner l'ordre des courbes $C'$, car on n'a pas d'occasion de recourir à des systèmes adjoints d'ordre différent.

\textbf{**)} \emph{Brill} et \emph{Nöther}; \emph{Segre} --- l.\ c.

\origpage{265}

Le système (complet) $|C'|$, adjoint au système $|C|$, est lié avec celui-ci de telle sorte, que si une transformation birationnelle (ou de Cremona) de tout le plan change les systèmes $|C|$ et $|C'|$ en $|D|$ et $|D'|$ (d'un autre plan), \emph{c'est encore $|D'|$ le système adjoint à $|D|$}, pourvu que l'on introduise convenablement en $|D'|$ les courbes exceptionnelles qui proviennent de la transformation${}^{*)}$; dans bien de questions il n'est pas même nécessaire de faire attention à ce dernier avertissement. C'est pourquoi nous pouvons dire que \emph{le lien qui passe entre un système linéaire de courbes planes et son système adjoint, est invariable par rapport aux transformations birationnelles du plan.}

Il vient de là que la considération d'un système de courbes planes avec les systèmes adjoints successifs (à savoir, l'adjoint $|C'|$ que nous venons de définir, l'adjoint de l'adjoint, ou second adjoint $|C''|$, etc.) peut rendre bien de services dans l'étude de la Géometrie du plan (ou des surfaces rationnelles), lorsque on prend comme groupe fondamental (dans le sens de M.\ \emph{Klein}) le groupe des transformations de Cremona. C'est M.\ S.\ \emph{Kantor} qui semble avoir aperçu, le premier, le profit que l'on pouvait tirer de cette méthode. En effet il a montré dans une Note des Comptes Rendus de l'Ac.\ d.\ Sc.\ (1885) comment on peut aborder par cette voie le problème de la détermination des transformations planes \emph{cycliques} de Cremona (transformations dont une puissance convenable donne l'identité); et il a continué cette recherche dans d'autres travaux plus récents${}^{**)}$. Plus tard M.\ \emph{Castelnuovo} a profité des systèmes adjoints pour étudier les systèmes linéaires des courbes planes${}^{***)}$, et M.\ \emph{Enriques} s'en est servi pour la détermination des groupes continus de transformations birationnelles du plan${}^{\dagger)}$.

La propriété d'invariance du système adjoint, que nous venons de rappeler, permet aussi de définir le système $|C'|$ adjoint à un système $|C|$ situé sur une surface rationnelle; il faudra en effet choisir $|C'|$ de telle sorte, que dans une représentation birationnelle quelconque de la surface sur le plan, $|C|$ et $|C'|$ correspondent à deux systèmes de courbes planes, dont le second soit l'adjoint au premier.

Or lorsque la surface rationnelle $F$ est donnée par ses caractères projectifs (à savoir, une surface de l'espace ordinaire, d'un ordre $n$ connu, etc.), on peut demander de construire directement sur $F$ le système $|C'|$ adjoint

\textbf{*)} \emph{Nöther}, \emph{Rationale Ausführung der Operationen}... §~31, Math.\ Annalen 23, 1883; \emph{Castelnuovo}, \emph{Ricerche generali sopra i sistemi lineari}... n${}^{\text{o}}$ 27.

\textbf{**)} \emph{Premiers fondements pour une théorie}..., 4${}^{\text{me}}$ partie, Atti dell'Accad.\ d.\ Sc.\ fis.\ e mat.\ di Napoli, serie II, vol.\ 4${}^{\text{o}}$. Voir aussi «Journal für die Mathem.\ 114», et «Acta Mathem.\ 19».

\textbf{***)} \emph{Ricerche generali sopra i sistemi lineari}..., l.\ c.

\textbf{†)} \emph{Sui gruppi continui}..., Rendic.\ della R.\ Accad.\ d.\ Lincei, 1893.

\origpage{266}
à un système donné $|C|$, sans avoir recours à la représentation plane. Supposons, pour plus de simplicité, que $|C|$ soit le système $\infty^{3}$ des sections planes de $F$; existera-t-il un système de surfaces découpant sur $F$ le système de courbes $|C'|$ adjoint à $|C|$? Certainement; on voit même${}^{*)}$ que cette propriété appartient à certaines surfaces $\Phi$ de l'ordre $n-3$, qui rencontrent tout plan général dans les courbes d'ordre $n-3$ adjointes à la section plane correspondante de $F$.

Il résulte de là que ces surfaces adjointes $\Phi$ sont comprises parmi les surfaces $\Psi$ du même ordre sous-adjointes à $F$. Or il y a lieu de se demander si les $\Phi$ coïncident avec les $\Psi$, ou bien s'il ne faut imposer aux $\Psi$ de nouvelles conditions pour obtenir les $\Phi$.

Dans certains cas on peut vérifier effectivement que les $\Phi$ et les $\Psi$ sont la même chose. Ainsi cela arrive si la surface $F$ n'a pas de \emph{points multiples $P$ isolés} ou \emph{propres}, c'est-à-dire des points $P$ tels que tout plan conduit par $P$ coupe sur la $F$ une courbe ayant le genre \emph{inférieur} au genre de la section plane générale de $F$. La surface générale du troisième ordre, la surface du quatrième ordre ayant une droite double, ou une section conique double, nous donnent des exemples de telles surfaces. La même particularité (coïncidence des $\Phi$ avec les $\Psi$) peut du reste se présenter aussi dans d'autres cas plus étendus${}^{**)}$.

Il ne faut penser pourtant qu'il en soit toujours ainsi. Au contraire, on démontre que les surfaces adjointes ont un point de l'ordre $i-2$ en tout point multiple (ordinaire) de l'ordre $i$ de $F$; tandis que les surfaces sous-adjointes peuvent ne passer du tout par ce point, ou bien y passer avec un ordre inférieur. D'où il résulte que le système des surfaces adjointes $\Phi$ peut bien avoir une dimension inférieure que le système des surfaces sous-adjointes $\Psi$.

Vérifions cette remarque sur un exemple. Considérons à cet effet la surface $F$ du quatrième ordre ayant un point triple $P$ (point multiple \emph{isolé}), et désignons par $|C|$ le système des sections planes de la $F$. On voit aisément que la $F$ peut se représenter sur le plan, de manière qu'au système $|C|$ corresponde un système $|c|$ de courbes planes du quatrième ordre rencontrant une même cubique (image de $P$) en 12 points fixes. Or les surfaces $\Psi$ d'ordre $4-3=1$ sous-adjointes à $F$, sont, d'après la définition, les plans de l'espace. Ceux-ci pourtant ne découpent pas sur $F$ le système $|C'|$ adjoint à $|C|$; car $|C'|$ est découpé par les seuls plans $\Phi$ qui passent par le point $P$; on le voit tout de suite si l'on recourt à la représentation plane, où $|C'|$ est représenté par le système $|c'|$ des droites adjointes aux courbes $|c|$. Il s'ensuit donc que du

\textbf{*)} \emph{Enriques}, \emph{Ricerche di geometria}... III, 5; \emph{Introduzione}... IV; \emph{Humbert}, \emph{Sur la théorie générale des surfaces unicursales}, Math.\ Annalen 45, 1894.

\textbf{**)} Voir à ce propos les Mémoires cités de MM.\ \emph{Enriques} et \emph{Humbert}.

\origpage{267}
système $\infty^{3}$ des $\Psi$ \emph{sous-adjointes} on déduit le système $\infty^{2}$ des $\Phi$ \emph{adjointes}, en imposant aux $\Psi$ la condition de passer par le point $P$.

Les considérations qui précèdent regardent le cas des points multiples ordinaires. Mais dès qu'il s'agit de singularités supérieures, le comportement des surfaces adjointes ne peut plus se définir d'une manière si simple, bien qu'on puisse recourir toujours à la représentation de la surface sur le plan.

Ce dernier moyen va pourtant faire défaut, lorsqu'on passe des surfaces rationnelles aux surfaces irrationnelles. En effet, même pour celles-ci, on peut définir un système de surfaces \emph{adjointes} $\Phi$, comprises parmi les sous-adjointes $\Psi$, et telles que les $\Phi$ découpent sur la surface un système \emph{adjoint} au système des sections planes, jouissant de propriétés analogues à celles que nous venons de voir au sujet des surfaces rationnelles. Le comportement des $\Phi$ en un point multiple isolé ordinaire de la nouvelle surface pourra s'établir par la même loi que nous avons exposée dans le cas des surfaces rationnelles; mais les points multiples singuliers présentent ici une difficulté encore plus grande; car on ne peut plus recourir à la représentation sur le plan.

On voit d'après cela l'opportunité d'établir une nouvelle définition du système adjoint, définition qui s'étende à toutes les surfaces algébriques, sans faire attention à leurs caractères projectifs (ordre, singularités, etc.), et qui permette, en conséquence, d'approfondir l'étude du lien (invariable par rapport aux transformations birationnelles) qui passe entre un système de courbes et son adjoint. Pour atteindre ce but il convient d'abord d'examiner les propriétés du système adjoint, dans le cas où nous avons tout-à-fait défini ce système, c'est-à-dire sur le plan. Nous chercherons ensuite d'étendre nos conclusions aux systèmes situés sur une surface quelconque.

\subsection*{15. Sur quelques propriétés caractéristiques du système adjoint.}

Reprenons donc un système linéaire $|C|$ de courbes planes d'ordre $m$ et genre $\pi$, et le système adjoint $|C'|$ dont les courbes ont l'ordre $m-3$; et tâchons d'établir le lien qui passe entre $|C|$ et $|C'|$, sans faire attention aux caractères projectifs des deux systèmes.

Nous remarquons d'abord que $|C'|$ \emph{découpe la série canonique} $g_{2\pi-2}^{\pi-1}$ \emph{sur toute courbe générale de} $|C|$.

Cette propriété pourtant ne suffit pas toujours à définir le système $|C'|$. Pour le voir, on n'a qu'à reprendre le système $\infty^{3}$ de courbes du quatrième ordre qui passent par 12 points d'une cubique, système qui représente sur le plan la surface du quatrième ordre ayant un

\origpage{268}
point triple; en effet le système lui-même a la propriété de découper sur chacune de ses courbes la série canonique $g_{4}^{2}$, tandis que le système adjoint $(\infty^{2})$ est formé par les droites du plan.

Malgré cela il nous convient de considérer le plus ample système $|C_{1}|$ découpant la série canonique sur la courbe générale de $|C|$; on démontre que ce système $|C_{1}|$ est unique et comprend (à des courbes fixes près) toute courbe qui rencontre les $C$ en de groupes canoniques. $|C_{1}|$ ne sera pas toujours le système adjoint à $|C|$; nous l'appelons \emph{système sous-adjoint à} $|C|$, car on voit aisément que sur une surface rationnelle $F$, dont les sections planes sont représentées par $\infty^{3}$ courbes de $|C|$, le système $|C_{1}|$ est découpé par les surfaces sous-adjointes à $F$ (surfaces de l'ordre $M-3$, si $M$ est l'ordre de $F$${}^{*)}$)).

Le système sous-adjoint $|C_{1}|$ contient toujours le système adjoint $|C'|$ (et peut même parfois coïncider avec celui-ci). Il s'agit maintenant d'examiner par quelles conditions on peut déduire le système adjoint du système sous-adjoint. Si nous reprenons encore une fois la surface rationnelle $F$ que nous venons de considérer, nous voyons sur le champ que la question proposée ne diffère pas de celle que nous avons rencontrée ci-dessus: établir les conditions qu'il faut imposer à une surface sous-adjointe $\Psi$, pour qu'elle devienne une surface adjointe $\Phi$. Ces conditions dépendaient de l'existence sur $F$ de certains points multiples \emph{isolés} (ou \emph{propres}), qui n'ont pas d'effet sur les surfaces sous-adjointes. Or remarquons qu'un point multiple isolé de $F$ est représenté sur le plan par une \emph{courbe fondamentale $\gamma$} du système $|C|$, courbe qui n'est pas rencontrée en de points variables par les $C$; c'est de plus une courbe fondamentale que nous dirons \emph{propre}, en faisant attention à la propriété que le système résiduel $|C-\gamma|$ a le genre inférieur au genre de $|C|$ (en relation avec le fait que le point multiple $P$ est \emph{isolé}). On a donc à rechercher sur le plan quelles sont les conditions qu'une courbe fondamentale propre du système $|C|$ impose aux courbes adjointes. Mais cette recherche (ainsi que son analogue sur les surfaces adjointes) ne manque pas de présenter parfois de sérieuses difficultés.

Voici comment on réussit, par une voie indirecte, à les éviter. Mettons en relation notre système $|C|$ avec un autre système linéaire quelconque $|D|$ de courbes du même plan. Envisageons encore le système $|C+D|$ somme de $|C|$ et $|D|$, et le système $|(C+D)'|$ adjoint a celui-là. La courbe générale du système $|C+D|$ a l'ordre $m+n$ (si $m$ et $n$ sont les ordres de $C$ et $D$), et a l'ordre de multiplicité $\mu+\nu$ en tout point qui est \emph{base} d'ordre $\mu$ et $\nu$ pour $|C|$ et $|D|$.

\textbf{*)} Par exemple, sur la surface du quatrième ordre ayant un point triple, le système sous-adjoint au système des sections planes, est ce même système.

\origpage{269}
Par suite la courbe générale du système $|(C+D)'|$ a l'ordre $m+n-3$ et la multiplicité $\mu + \nu - 1$ dans le point nommé. Or si nous comparons ce système $|(C + D)'|$ avec le système $|C' + D)|$ somme de $|C'|$ et de $|D|$, nous voyons que les deux systèmes sont formés de courbes du même ordre $(m - 3) + n$, qui ont la même multiplicité $(\mu - 1) + \nu$ en tout point qui est \emph{base} $(\mu \geqq 1)$ pour $|C|$. Mais une distinction entre les deux systèmes se présente en tout point (s'il en existe), qui est base pour $|D|$, mais pas pour $|C|$ $(\mu = 0, \nu > 0)$; car un tel point a l'ordre $\mu + \nu - 1$ pour le système $|(C + D)'|$, et l'ordre $\mu + \nu$ pour le système $|C' + D|$. D'après cela nous pouvons énoncer le résultat suivant:

\emph{Si $|C|$ et $|D|$ sont deux systèmes de courbes planes, le système $|(C + D)'|$ adjoint à leur somme contient le système $|C' + D|$; et précisément si l'on veut déduire $|C' + D|$ de $|(C + D)'|$, il faut augmenter d'une unité l'ordre de tout point qui est base pour $|D|$, mais ne l'est pas pour $|C|$.}

Or il y a lieu de se demander: cette propriété subsiste-t-elle aussi pour les systèmes sous-adjointes? Pas toujours. Et alors, est-ce qu'on ne pourra se servir de la même propriété pour séparer le système adjoint du système sous-adjoint qui le contient?

Il en est précisément ainsi. Partons d'un système $|C|$ quelconque, et choisissons un système auxiliaire $|D|$ qui n'ait pas de courbes fondamentales propres; (cela est possible d'une infinité de manières). Maintenant que l'on choisisse parmi les courbes du système $|C_{1}|$ sous-adjoint à $|C|$, celles $C'$ qui, additionnées à une courbe $D$ quelconque, donnent des courbes sous-adjointes à $|C + D|$; on peut démontrer que ces courbes $C'$ forment un système linéaire qui ne dépend pas de $|D|$; \emph{$|C'|$ est le système adjoint à $|C|$}. On peut tirer de cette remarque une nouvelle définition du système adjoint à un système donné de courbes planes. Or, et c'est là un résultat fort intéressant, cette définition peut se transporter aux systèmes linéaires situés sur une surface algébrique \emph{quelconque}. On parvient ainsi à l'énoncé suivant, dont le lecteur trouvera ailleurs la démonstration${}^{*)}$.

\subsection*{16. Système adjoint à un système linéaire donné sur une surface algébrique.}

\emph{Lorsqu'on connaît sur une surface algébrique un système linéaire irréductible $|C|$, on peut construire un second système $|C'|$, appelé système adjoint à $|C|$, qui est tout-à-fait défini par les propriétés suivantes:}

\textbf{*)} A la définition du système adjoint ici rapportée, est consacré le Chap.\ III de la \emph{Introduzione}... de M.\ \emph{Enriques}.

\origpage{270}

1) \emph{les courbes de $|C'|$ découpent des groupes de la série canonique sur toute courbe générale de $|C|$;}

2) \emph{quel que soit le système $|D|$, le système $|C' + D|$ découpe des groupes de la série canonique sur toute courbe générale de $|C + D|$ (supposée irréductible). Ces groupes toutefois contiennent $i$ points réunis dans chaque point qui a l'ordre $i$ pour $|D|$, et n'est pas base pour $|C|$.}

Il faut bien remarquer que la série $g_{2\pi-2}$ découpée par le système adjoint $|C'|$ sur la courbe $C$ générale (supposée de genre $\pi$), n'est pas nécessairement complète; sur quelques surfaces (au contraire de ce qui arrive sur le plan et sur les surfaces rationnelles) la dite série peut bien avoir une dimension inférieure à $\pi - 1$. Elle peut même faire défaut, avec le système adjoint lui-même (en dehors du cas évident $\pi = 0$). Ainsi l'on démontre${}^{*)}$ qu'il n'y a pas de système adjoint au système linéaire des sections planes d'une surface réglée de genre $\pi$ quelconque.

Il convient souvent d'énoncer, sous forme de théorème, la définition que nous venons de donner du système adjoint, ou du moins la partie essentielle de la même définition. On parvient ainsi à la \emph{propriété fondamentale du système adjoint}:

\emph{Soient $|C|$ et $|D|$ deux systèmes linéaires de courbes sur une même surface; le système $|C'|$ adjoint au premier, additionné avec le second, donne un système $|C' + D|$ qui est contenu dans le système $|(C + D)'|$ adjoint à $|C + D|$; il en diffère seulement en cela, qu'il peut exister des points en chacun desquels l'ordre de multiplicité de $|C' + D|$ est supérieur à l'ordre correspondant de $|(C + D)'|$.} Ce sont précisément les points qui sont base, d'ordre $\nu$ par exemple, pour $|D|$, tandis qu'il ne sont pas base pour $|C|$; en effet dans chacun de ces points $|C' + D|$ a l'ordre de multiplicité $\nu$, tandis que $|(C + D)'|$ a l'ordre de multiplicité $\nu - 1$. Si en particulier ces points font défaut, de sorte que tout point-base de $|D|$ est aussi point-base de $|C|$, on a
\[
|C' + D| = |(C + D)'|.
\]

D'après le théorème qui précède, on peut, par une simple addition, construire le système $|(C + D)'|$ adjoint à $|C + D|$, dès qu'on connaît le système $|C'|$ adjoint à $|C|$. Inversement pour déduire le système adjoint à $|C|$ du système adjoint à $|C + D|$, on n'a qu'à retrancher le système $|D|$ du dernier système adjoint. De ces observations il résulte aisément que, dès qu'on connaît le système $|C'|$ adjoint à un système particulier $|C|$ sur une surface, on peut construire le système adjoint à un système arbitrairement donné sur la surface, à l'aide de simples opérations d'addition et de soustraction (avec des impositions ou suppressions de points multiples, s'il est nécessaire).

\textbf{*)} \emph{Enriques}, \emph{Introduzione}..., 30.

\origpage{271}

\emph{Remarque.} --- La définition qui précède, comme on l'a vu, n'établit pas d'une manière directe les caractères de l'opération d'\emph{adjonction}, par laquelle on passe d'un système donné $|C|$ au système adjoint $|C'|$. Au contraire on définit cette opération par son comportement en relation avec une autre opération connue d'avance (l'addition). Ce comportement est représenté par la relation
\[
|(C + D)'| = |C' + D|
\]
qui subsiste toujours, à des multiplicités dans certains points près. Or dans l'Analyse on a souvent l'occasion de recourir à de telles définitions; toutes les fois qu'on définit une fonction, non pas directement, mais à l'aide d'une \emph{équation fonctionnelle}. Ce rapprochement avec des théories connues peut servir à mettre en lumière la voie que nous venons de suivre.

\subsection*{17. Définition directe du système adjoint.}

Dans les cas les plus simples, on peut donner une définition directe du système adjoint à un système donné, définition qui s'ensuit de celle qui précède. Bornons-nous à l'énoncé suivant:

\emph{Soit $|C|$ un système qui ne possède pas de courbes fondamentales propres} (c'est-à-dire des courbes $\gamma$ qui ne sont pas rencontrées en de points variables par les $C$, et telles que $|C - \gamma|$ a le genre inférieur au genre de $|C|$); \emph{alors, en général, toute courbe qui découpe un groupe de la série canonique sur les courbes $C$, appartient au système adjoint à $|C|$.} Les mots \emph{en général} se rapportent aux observations suivantes.

Pour démontrer le théorème on suppose d'abord que les courbes $C$, qui passent par un point quelconque de la surface, ne vont passer en conséquence par d'autres points mobiles avec celui-ci. On rencontre ensuite une seconde restriction; elle se rapporte au cas où la surface appartenant à la \emph{classe} des surfaces réglées, les courbes $C$ découperaient en un point les génératrices. Tandis que la première restriction est peut-être superflue, la seconde est certainement nécessaire. En effet le système $|C|$ des sections planes d'une surface réglée n'a pas de système adjoint, mais il y a bien des courbes qui découpent sur les $C$ des groupes de la série canonique.

En revenant à la définition directe du système adjoint, il y aurait maintenant à considérer le cas d'un système $|C|$ doué de courbes fondamentales propres. Pour ne fatiguer le lecteur nous n'insisterons pas sur ce sujet; on pourra consulter pour ces développements le n${}^{\text{o}}$ 29 de la \emph{Introduzione}... de M.\ \emph{Enriques}.

\origpage{272}

\subsection*{18. Surfaces adjointes à une surface donnée.}

La définition du système adjoint à un système donné (n${}^{\text{o}}$ 16) vient d'être énoncée sous une forme qui ne subit aucun changement; lorsqu'on transforme birationnellement la surface considérée; un système $|C|$ et son adjoint $|C'|$ se transforment en deux systèmes $|D|$ et $|D'|$ dont le second est encore l'adjoint au premier (à des courbes \emph{exceptionnelles} près). C'est pourquoi cette définition se prête mieux que toute autre dans les recherches des propriétés d'invariance d'une surface algébrique; nous allons le voir bientôt.

Mais d'abord il nous convient de faire quelques remarques au point de vue projectif.

Soit $|C|$ le système des sections planes d'une surface $F$ d'ordre $n$ de l'espace ordinaire; et soit $|C'|$ le système adjoint à $|C|$. Alors on peut construire un système linéaire de surfaces $\Phi$ d'ordre $n - 3$, sous-adjointes à $F$, qui découpent sur $F$ (en dehors des courbes multiples de $F$) les courbes de $|C'|$. Ces surfaces $\Phi$ sont tout-à-fait déterminées par les courbes $C'$; on les appelle \emph{surfaces adjointes à $F$}, de l'ordre $n - 3$; nous en avons dejà parlé dans le cas où $F$ était une surface rationnelle, et nous avons dit alors que cette définition pourrait s'étendre aux surfaces générales.

Sans insister nouvellement sur le même sujet, bornons-nous à remarquer que les considérations précédentes fournissent une définition indirecte des surfaces $\Phi$, à l'aide du système $|C'|$ que celles-ci découpent sur $F$. Il résulte de la définition que les $\Phi$ sont comprises parmi les surfaces sous-adjointes $\Psi$ du même ordre $n - 3$, et qu'elles en diffèrent seulement par leur comportement dans les points multiples isolés de la $F$. De sorte que si la $F$ ne possède de tels points, les surfaces adjointes coïncident avec les sous-adjointes; elles sont alors tout-à-fait déterminées par la condition de découper sur tout plan général des courbes adjointes à la section correspondante de $F$. C'est ce qui arrive par exemple si la $F$ est une surface générale de l'ordre $n$, ou si elle possède une courbe double et un nombre fini de points d'ordre $k = 3, 4, \ldots$ satisfaisant à la condition d'avoir l'ordre $\dfrac{k(k-1)}{2}$ pour la courbe double.

Si au contraire la $F$ possède des points multiples \emph{isolés ordinaires}, de l'ordre $k$, on démontre que les $\Phi$ passent avec l'ordre $k - 2$ par chacun d'eux. De sorte que l'on peut énoncer le résultat suivant:

\emph{Si une surface $F$ n'a d'autres singularités que des courbes multiples ordinaires de l'ordre $h = 2, 3 \ldots$, et des points multiples ordinaires de l'ordre $k = 2, 3 \ldots$, les surfaces adjointes sont tout-à-fait déterminées par les conditions de passer $h - 1$ fois par ces courbes et $k - 2$ fois}

\origpage{273}
\emph{par ces points.} En particulier un point double, isolé, ordinaire, ne présente aucune condition aux surfaces adjointes.

L'énoncé qui précède nous donne une définition directe des surfaces adjointes à une surface douée de singularités ordinaires. On pourrait chercher d'étendre cette définition aux surfaces ayant des points multiples singuliers; mais on rencontre quelques difficultés si l'on veut comprendre, par un même énoncé, tous les cas que ces points peuvent présenter. C'est pourquoi nous avons suivi à ce propos une voie indirecte, en définissant les surfaces $\Phi$ adjointes à $F$ à l'aide du système $|C'|$, qu'elles découpent sur $F$. Mais on aurait pu aborder le problème du comportement des surfaces $\Phi$ en un point singulier en s'appuyant sur des considérations analytiques, comme M.\ \emph{Nöther} a montré${}^{*)}$. Cette dernière voie (qui au point de vue historique précède la nôtre) permettrait alors de définir le système adjoint $|C'|$ à l'aide des $\Phi$, tandis que nous avons suivi le chemin inverse.

\emph{Remarque.} --- Ainsi que nous l'avons fait sur le plan, de même l'on pourrait définir sur une surface quelconque les courbes $C_{1}$ \emph{sous-adjointes} à un système $|C|$, par l'unique condition de découper sur toute courbe $C$ des groupes de la série canonique. Ces courbes forment un système linéaire $|C_{1}|$ (\emph{sous-adjoint} à $|C|$), avec une seule exception relative aux surfaces réglées. Le système $|C_{1}|$ contient le système adjoint $|C'|$.

Si $|C|$ est le système des sections planes d'une surface $F$ de l'ordre $n$ de l'espace ordinaire, on démontre que $|C_{1}|$ est découpé par les surfaces d'ordre $n-3$ sous-adjointes à $F$; le cas excepté où $F$ est une surface réglée de genre $\pi > 0$, car alors il y a des $C_{1}$ composées de $2\pi-2$ génératrices, mais il n'existe aucune surface de l'ordre $n-3$ sous-adjointe à $F$, qui rencontre $F$ dans ces génératrices seulement (en dehors des courbes multiples).

Du reste, nous n'aurons pas à considérer, dans la suite, le système sous-adjoint.

\subsection*{19. Surfaces adjointes d'autres ordres.}

Lorsqu'on a défini, par les moyens qu'on vient de voir, les surfaces $\Phi$ d'ordre $n-3$ adjointes à $F$, on peut aussi construire des \emph{surfaces adjointes d'ordre} $n-3+i$ $(i \gtreqqless 0)$, c'est-à-dire des surfaces qui ont

\textbf{*)} \emph{Zur Theorie}... Math.\ Annalen 8, pag.\ 510; voir aussi une Note du même auteur dans les Götting.\ Nachrichten (1871), et le Mémoire \emph{Ueber die totalen algebraischen Differentialausdrücke}, Math.\ Annalen 29. M.\ \emph{Nöther} s'occupe en particulier des surfaces adjointes de l'ordre $n-4$ ($\varphi$-\emph{Flächen}), que nous allons étudier ensuite. Voir aussi (pour des surfaces particulières) le Mémoire cité de M.\ \emph{Humbert}, et l'autre Mémoire du même auteur \emph{Théorie générale des surfaces hyperelliptiques}, Journal de Mathém.\ 4${}^{\text{e}}$ s${}^{\text{e}}$, t${}^{\text{e}}$ IX, pag.\ 394, 1893.

\origpage{274}
cet ordre, et qui se comportent comme les $\Phi$ le long des courbes multiples et dans les points multiples de $F$. Les nouvelles surfaces sont aussi déterminées par la condition de découper sur $F$ les courbes du système $|C'+iC|$, où $|C|$ est le système des sections planes de $F$, et $|C'|$ est le système adjoint à celui-ci. Or d'après le théorème fondamental sur le système adjoint, $|C'+iC|$ n'est autre chose que le système adjoint au système $|(i+1)C|$, multiple de $|C|$. On a dans cette remarque un moyen pour définir le système $|C'+iC|$ (et par suite les surfaces adjointes d'ordre $n-3+i$ qui le découpent) toutes les fois que $|C'|$ fait défaut; car même alors le système $|(i+1)C|$ peut bien avoir un système adjoint; et il l'a certainement (on le démontre) lorsque $i$ surpasse une certaine valeur.

Nous allons rencontrer dans la suite ces systèmes $|C'+iC|$; il nous faudra alors récourir à la remarque, que nous faisons dès maintenant. Nous avons déjà remarqué (n${}^{\text{o}}$ 17) que le système $|C'|$ adjoint à un système $|C|$ de genre $\pi$, découpe sur la courbe $C$ générale une série $g_{2\pi-2}$, qui est contenue dans la série canonique $g_{2\pi-2}^{\pi-1}$, mais qui peut bien être incomplète et avoir par conséquent un certain \emph{défaut}; elle peut même manquer (l'exemple des surfaces réglées nous l'a prouvé), et dans ce cas le défaut atteint son \emph{maximum} $\pi$. De même la série que le système $|C'+iC|$ $(i > 0)$ découpe sur $|C|$ est une série $g_{2\pi-2+in}$ (où $n$ désigne le degré de $|C|$); si elle était complète, sa dimension serait $\pi-2+in$; mais il peut se faire qu'elle aussi soit incomplète, et alors on aura à considérer son défaut. Or, et c'est là la remarque que nous nous proposions de faire, la considération de ce défaut n'a lieu que pour les premières valeurs de $i$; car \emph{on peut toujours fixer une telle limite $l$, que pour les valeurs de $i > l$ la série découpée par $|C'+iC|$ sur $C$ résulte complète.}

\subsection*{20. Surfaces adjointes à une courbe gauche.}

Dans le n${}^{\text{o}}$ 18 nous avons appris à construire des surfaces qui découpent sur une surface donnée $F$ d'ordre $n$ le système adjoint aux sections planes de celle-ci. De même, à l'aide du théorème fondamental, on peut construire des surfaces qui découpent sur $F$ le système adjoint à un système de courbes $|C|$ \emph{quelconque} donné sur $F$. A cet effet on fait passer d'abord par la courbe $C$ générale une surface $G$ d'un ordre $m$ assez élevé pour que cette construction soit possible. Appelons $K$ la courbe intersection résiduelle de la surface $G$ avec $F$; on aura alors $|K+C| = |mC|$. Or le système $|C'+(m-1)C|$ adjoint à celui-ci est découpé sur $F$ par les surfaces d'ordre

\origpage{275}
\[ n-3+(m-1) = m+n-4 \]
adjointes à $F$ (n${}^{\text{o}}$ 19). Donc, d'après le théorème fondamental:

\emph{Si nous imposons à ces surfaces adjointes d'ordre $m+n-4$ la condition de passer par la courbe $K$, et d'avoir avec $F$ un contact d'ordre $i-1$ dans tout point-base d'ordre $i$ pour le système $|C|$, nous allons découper sur $F$ le système adjoint à $|C|$.}

Le théorème qu'on vient d'énoncer, nous fournit tout d'abord le moyen de découper sur une courbe gauche $C$, intersection partiale de deux surfaces $F$ et $G$ ayant les ordres $n$, $m$, la série canonique, ou du moins une partie de cette série. Nous voyons qu'il faut construire à ce but les surfaces d'ordre $m+n-4$, qui passent par l'intersection résiduelle $K$ de $F$ et $G$, et se comportent d'une manière convenable dans les points multiples de ces surfaces. Or c'est là précisément un théorème de M.\ \emph{Nöther} sur les surfaces \emph{adjointes} à une courbe gauche${}^{*)}$.

Mais il résulte encore du théoreme qui précède, que toute courbe adjointe à $|C|$ sur la surface $F$ peut être découpée par les dites surfaces adjointes, ce qui permet, pour ainsi dire, l'inversion du théorème de M.\ \emph{Nöther}, autant qu'elle est possible.

\section*{Chapitre V. Invariants d'une surface par rapport aux transformations birationnelles.}

En nous appuyant sur les résultats qui précèdent nous sommes maintenant en mesure d'aborder la question fondamentale de notre théorie; à savoir, la recherche de quelques invariants d'une surface par rapport aux transformations birationnelles; invariants qui peuvent être des nombres ou des variétés géométriques (courbes ou systèmes de courbes).

A cette recherche il convient cependant de faire précéder quelques remarques d'un caractère général.

A la notion des invariants on est parvenu tout d'abord en considérant les expressions transcendantes liées à la surface (\emph{Clebsch}, \emph{Nöther} 1868---69; \emph{Picard} 1884). Ensuite on s'est proposé de définir les mêmes invariants par la voie géométrique (ou algébrique). Les géomètres qui ont abordé cette question (\emph{Cayley}, \emph{Nöther}, \emph{Zeuthen}), procédaient de la manière suivante. On considérait une surface $F$ de l'espace ordinaire, et on faisait attention à ses caractères projectifs (ordre, courbes et points multiples). On transformait ensuite birationnellement

\textbf{*)} \emph{Zur Theorie}... pag.\ 510, Math.\ Annalen 8; voir aussi le Mémoire, \emph{Ueber die algebraischen Raumcurven}, Denkschriften der Berliner Akad.\ 1882, pag.\ 12.

\origpage{276}
la $F$ en une autre surface $F'$ du même espace, et en comparant les caractères projectifs de $F$ aux caractères analogues de $F'$, on formait des expressions numériques ou des fonctions, qui n'étaient pas modifiées par la transformation.

Nous allons maintenant remplacer cette méthode par une autre qui présente sur celle-là quelque avantage, bien qu'elle n'en diffère pas d'une manière essentielle; elle nous permet pourtant de faire abstraction des caractères projectifs des surfaces, et d'éviter par là les difficultés qui proviennent des points singuliers et de leur comportement dans les transformations.

A cet effet reprenons la surface générale d'une \emph{classe} donnée (n${}^{\text{o}}$ 3), surface dont les caractères projectifs n'ont pour nous aucun intérêt, et partons d'un système linéaire de courbes $|C|$ de la surface. De $|C|$ nous pouvons déduire, par quelque opération, de nouveaux systèmes qui sont liés avec $|C|$; tels sont, par exemple, le système adjoint $|C'|$, l'adjoint $|C''|$ de $|C'|$, etc. Que l'on cherche maintenant des rélations entre ces systèmes, ou entre leurs caractères (genre, dimension, degré); si l'on parvient à une rélation qui ne dépend pas du système $|C|$ d'où l'on part, on aura là un \emph{invariant de la surface}.

Les exemples que nous allons donner, montreront la fécondité de cette méthode.

\subsection*{21. Système canonique.}

L'exemple le plus important nous est fourni par l'application immédiate du théorème fondamental sur le système adjoint.

Ce théorème en effet nous dit, que si $|C|$ et $|D|$ sont deux systèmes de courbes sur la surface, on parvient au système $|(C+D)'|$ adjoint à $|C+D|$ de deux manières différentes: on peut additionner à $|D|$ le système $|C'|$ adjoint à $|C|$, mais on peut de même additionner à $|C|$ le système $|D'|$ adjoint à $|D|$. On a donc l'égalité symbolique${}^{*)}$
\[ |C'+D| = |C+D'|, \]
d'où l'on déduit
\[ |C'-C| = |D'-D|. \]
Cette rélation a une signification géométrique seulement dans le cas où $|C'|$ contient $|C|$, car elle nous dit alors que $|D'|$ aussi contient $|D|$, et en outre que le système $|K| = |C'-C|$ résiduel de $|C|$ par

\textbf{*)} Au sujet des ordres de multiplicité des points-base qui appartiennent aux deux systèmes considérés dans cette rélation, il y aurait une observation à faire, car ces ordres peuvent différer d'une unité; qu'il nous soit permis de l'omettre, attendu qu'ici nous nous proposons seulement d'esquisser la démonstration. Voir la démonstration complète dans le n${}^{\text{o}}$ 38 de l'\emph{Introduzione}... de M.\ \emph{Enriques}.

\origpage{277}
rapport à $|C'|$, est de même le résiduel de tout autre système $|D|$ par rapport à son adjoint $|D'|$. En conclusion:

\emph{Le résiduel $|K|$ d'un système linéaire $|C|$ par rapport à son système adjoint, ne dépend pas du système $|C|$ d'où l'on part.}

En d'autres termes, $|K|$ est un système de courbes qui dépend seulement de la surface considérée, ou plus précisément de la \emph{classe} de surfaces à laquelle celle-là appartient; c'est donc un \emph{invariant} de la dite surface. Nous l'appellerons \emph{système canonique} (et courbe canonique chacune de ses courbes), car, ainsi que nous allons le voir, il a la plus grande analogie avec la série canonique définie sur les courbes.

La dimension du système canonique $|K|$ est aussi un \emph{caractère invariable} de la surface. Cette dimension augmentée d'une unité s'appelle \emph{premier genre} de la surface (\emph{Flächengeschlecht} d'après M.\ \emph{Nöther})${}^{*)}$; il convient d'ajouter \emph{genre géométrique}, pour le distinguer du genre numérique que nous allons introduire bientôt; on le désigne par $p_{g}$.

Si $p_{g} = 0$, il n'y a pas de système canonique; en d'autres termes, il n'y a sur la surface aucun système $|C|$, qui soit contenu dans son adjoint $|C'|$. C'est ce qui arrive, par exemple, sur les surfaces rationnelles, sur les surfaces réglées, ou plus généralement sur toute surface contenant un système $|C|$ de courbes dont la série caractéristique est non-spéciale.

Si $p_{g} = 1$, la courbe générale d'un système quelconque $|C|$ appartient à \emph{une seule} courbe du système adjoint $|C'|$. Le reste $C' - C$, s'il y en a, fournit (à des courbes exceptionnelles près) la seule courbe canonique. Mais dans certains cas il peut se faire que $|C|$ coïncide avec $|C'|$; alors il n'y a pas de reste, et la courbe canonique a, pour ainsi dire, l'ordre zéro.

Si $p_{g} \geqq 1$, le système canonique $|K|$ peut se construire à l'aide d'un seul système de courbes $|C|$ de la surface, et du système adjoint $|C'|$; cela résulte de la définition. En rappelant les propriétés du système adjoint, on trouve, pour les courbes canoniques, la propriété suivante:

\emph{Toute courbe canonique découpe sur la courbe générale $C$ d'un système quelconque $|C|$ un groupe de points, qui joint aux points-base de $|C|$ et à un groupe de la série caractéristique de $|C|$, fournit un groupe de la série canonique de $C$.}

La propriété énoncée est caractéristique pour les courbes canoniques; il suffit même, pour définir ces courbes, qu'elle soit vérifiée par rapport à un seul système $|C|$, pourvu que ce système satisfasse aux restrictions dont on parle au n${}^{\text{o}}$ 17.

\textbf{*)} \emph{Zur Theorie}... pag.\ 520, ibid.

\origpage{278}

\subsection*{22. Construction du système canonique à l'aide des surfaces adjointes de l'ordre $n-4$.}

Lorsqu'une surface $F$ est donnée par ses caractères projectifs dans l'espace ordinaire, et l'on désigne par $n$ son ordre, il convient, dans la construction du système canonique $|K|$, de recourir au système $|C|$ des sections planes de $F$. On a à retrancher le système $|C|$ du système adjoint $|C'|$, qui est découpé sur $F$ par les surfaces adjointes d'ordre $n-3$. Le système résiduel sera naturellement déterminé par les surfaces adjointes d'ordre $n-4$. On a donc ce résultat:

\emph{Sur une surface d'ordre $n$ de l'espace ordinaire le système canonique est découpé par les surfaces adjointes de l'ordre $n-4$}${}^{*)}$. Le nombre de telles surfaces qui sont linéairement indépendantes, fournit donc le genre $p_{g}$ de $F$. Ici, comme d'ordinaire, on envisage l'intersection d'une surface adjointe avec $F$, en dehors des courbes multiples de $F$. Mais dans cette intersection on peut trouver parfois des courbes qui appartiennent à toutes les surfaces adjointes d'ordre $n-4$, et qu'on ne doit pas regarder comme des courbes canoniques. On démontre en effet que si une courbe $\gamma$ de $F$ est \emph{exceptionnelle}, c'est-à-dire, si elle correspond, par une transformation birationnelle, à un point $P'$ d'une surface $F'$ (n${}^{\text{o}}$ 2), cette courbe $\gamma$ appartient à toutes les surfaces d'ordre $n-4$ adjointes à $F$.

Or au sujet des courbes exceptionnelles $\gamma$ de $F$, il y a une distinction à faire. En effet le point $P'$ de $F'$, qui correspond à la courbe $\gamma$, n'appartient pas, en général, à toutes les courbes canoniques de $F'$ (n'est pas point-base du système canonique sur $F'$), et dans ce cas la courbe exceptionnelle $\gamma$, que nous dirons de la \emph{première espèce}, ne doit pas être regardée comme faisant partie des courbes canoniques de $F$; par suite il faut retrancher la $\gamma$ de l'intersection de $F$ avec une surface adjointe générale de l'ordre $n-4$, pour obtenir une courbe canonique. On remarquera encore qu'une courbe exceptionnelle de la première espèce n'a aucune intersection variable avec les courbes canoniques${}^{**)}$.

Mais il peut arriver au contraire, dans quelques cas particuliers, que $P'$ soit un point-base du système canonique sur $F'$; alors la courbe exceptionnelle $\gamma$ correspondante sur $F$, courbe exceptionnelle de la \emph{seconde espèce}, jouit de la propriété d'invariance, et doit être regardée comme faisant partie de toutes les courbes canoniques. Les courbes

\textbf{*)} \emph{Nöther}, \emph{Zur Theorie}... pag.\ 514, ibid.

\textbf{**)} Pour avoir un exemple d'une telle courbe exceptionnelle il suffit de considérer la droite qui joint deux points triples d'une surface du cinquième ordre; ici les surfaces adjointes de l'ordre $5-4=1$ sont les plans conduits par cette droite.

\origpage{279}
résiduelles variables découpent à present $\gamma$ en de points variables (un au moins).

Voici comment on peut obtenir un exemple d'une courbe exceptionnelle de la seconde espèce. Qu'on envisage une surface du cinquième ordre ayant deux contacts avec elle-même (\emph{points tacnodals})${}^{*)}$; chacun de ces points impose une condition aux surfaces adjointes, de sorte que les surfaces adjointes de l'ordre $5-4=1$ sont les plans passant par ces points. Les courbes canoniques sont donc des courbes du cinquième ordre (genre 2), qui ont un point-base $P'$ (en dehors des points singuliers) dans l'intersection ultérieure de la surface avec la droite joignant les deux points singuliers. Toute transformation birationnelle de la surface qui change le point $P'$ en une courbe $\gamma$, donne lieu à une courbe exceptionnelle de la seconde espèce.

En laissant de côté l'exemple pour revenir à la remarque qui précède, on peut maintenant préciser l'énoncé du dernier théorème:

\emph{Les surfaces d'ordre $n-4$ adjointes à une surface de l'ordre $n$ découpent sur celle-ci le système canonique, en dehors des courbes multiples et des courbes exceptionnelles de la première espèce.}

La considération des surfaces adjointes de l'ordre $n-4$ pourrait même servir à définir le système canonique. C'est précisément la voie que \emph{Clebsch} a indiquée, et que M.\ \emph{Nöther} a suivie dans le Mémoire cité. Ce savant a dû, naturellement, surmonter la difficulté de démontrer le caractère d'invariance du système qu'il venait à construire ainsi. Nous avons évité cette difficulté en définissant le système canonique à l'aide de la propriété fondamentale du système adjoint, laquelle peut même remplacer la propriété d'invariance du système canonique, pour les surfaces $(p_{g} = 0)$ qui ne possèdent pas un tel système.

La construction du système canonique à l'aide des surfaces adjointes de l'ordre $n-4$, montre que ce système a la plus grande analogie avec la série canonique découpée sur une courbe plane de l'ordre $n$ par les courbes adjointes de l'ordre $n-3$; on voit ainsi la raison du nom que nous avons attribué à ce système.

\subsection*{23. Système canonique réductible.}

Le système canonique appartenant à une surface donnée peut être réductible, soit par la présence de quelque courbe fixe commune à

\textbf{*)} Si 1 et 2 (points fondamentaux du système des coordonnées) sont les deux points singuliers, l'équation de la surface est de la forme
\[
a x_{1}^{3} x_{2}^{2} + b x_{2}^{3} x_{1}^{2} + x_{1}^{2} x_{2}^{2} f_{1} + x_{1} x_{2} (x_{1} g_{2} + x_{2} h_{2} + i_{3}) + x_{1} j_{4} + x_{2} k_{4} + l_{5} = 0
\]
où $a$ et $b$ sont des constantes, et $f_{1}, g_{2} \ldots l_{5}$ sont des formes algébriques dans les variables $x_{3}$, $x_{4}$ ayant les degrés $1, 2 \ldots 5$.

\origpage{280}
toute courbe canonique, soit par la division de la courbe canonique générale en plusieurs courbes variables dans un faisceau (n${}^{\text{o}}$ 7).

Le premier cas se présente, nous le savons, toutes les fois qu'il y a sur la surface quelque courbe exceptionnelle de la seconde espèce. Mais il y a lieu de rechercher si le système canonique ne peut contenir d'autres courbes fixes, en dehors de ces courbes exceptionnelles. Il paraît qu'il en puisse contenir en effet, bien qu'il ne soit pas facile de trouver des exemples à ce sujet. Sans nous arrêter sur la question, nous nous bornons à montrer, par un exemple, que, pour les surfaces ayant $p_{g} = 1$, la courbe canonique n'est pas nécessairement exceptionnelle, ainsi qu'on a été porté à le croire; au contraire il peut se faire que cette courbe (ou quelqu'une de ses parties) ait le genre $> 0$. Considérons en effet une surface du cinquième ordre passant par une section conique, et ayant trois contacts avec elle-même en trois points $A$, $B$, $C$ de cette courbe. L'unique surface adjointe de l'ordre $5-4=1$ est ici le plan $ABC$, qui rencontre la surface le long de la section conique nommée et d'une courbe ultérieure du troisième ordre passant simplement par $A$, $B$, $C$. Or il est clair que cette dernière courbe (ayant le genre $> 0$) n'est certainement pas exceptionnelle.

Il peut encore arriver, nous avons dit, que la courbe canonique générale (en dehors peut-être de parties fixes) se décompose en plusieurs courbes (au moins $p_{g} - 1$) variables dans un même faisceau (qui peut être rationnel ou bien irrationnel). Si le faisceau n'a pas de points-base, et s'il n'existe aucune courbe exceptionnelle de la seconde espèce, on démontre (d'après M.\ \emph{Nöther}${}^{*)}$) que \emph{chacune des courbes partielles a le genre} 1. La restriction énoncée est nécessaire, car il y a des surfaces de genre $p_{g} = 2$ ayant un faisceau de courbes canoniques de genre $> 1$ (avec des points-base). On n'a qu'à rappeler l'exemple de la surface du cinquième ordre cité au n${}^{\text{o}}$ 22. Inversement \emph{si une surface $(p_{g} \geqq 2)$ contient un faisceau de courbes du genre} 1, \emph{toute courbe canonique se decompose en $i \geqq p_{g} - 1$ courbes de ce faisceau.}

\subsection*{24. Les caractères $p^{(1)}$ et $p^{(2)}$ d'une surface.}

Pour obtenir des nombres qui jouissent de la propriété d'invariance par rapport à une surface donnée, on n'a qu'à considérer les autres caractères du système canonique, à côté de la dimension $p_{g} - 1$ dont nous avons déjà parlé. Ainsi se présentent le genre et le degré du système canonique.

\textbf{*)} \emph{Zur Theorie}... pag.\ 523, ib.\ Voir aussi \emph{Castelnuovo}, \emph{Osservazioni intorno alla geometria sopra una superficie}, Rendic.\ Istituto Lombardo 1891. Nota I, §~2.

\origpage{281}

Le genre (de la courbe canonique générale), dans la supposition $p_{g} \geqq 1$, nous fournit l'invariant qu'on appelle \emph{second genre} (\emph{Curvengeschlecht} d'après M.\ \emph{Nöther}${}^{*)}$) et qu'on désigne par $p^{(1)}$. On peut le calculer directement pour les seules surfaces qui ont des vraies courbes canoniques (d'ordre $> 0$)${}^{**)}$; pour les autres surfaces on ne saurait le définir qu'à l'aide d'une convention; c'est ce que nous allons faire tout à l'heure, du moins dans certains cas.

Mais d'abord parlons aussi du \emph{troisième genre} $p^{(2)}$${}^{***)}$, qui est le degré du système canonique, c'est-à-dire le nombre des intersections variables de deux courbes canoniques; pour que cette définition ait un sens, par elle-même, il faut supposer qu'on ait $p_{g} > 1$ (pour $p_{g} = 2$ ou aura $p^{(2)} = 0$), et en outre que la courbe canonique générale soit irréductible.

A ce dernier caractère pourtant, on ne fait pas attention d'ordinaire, car M.\ \emph{Nöther} à montré${}^{\dagger)}$ qu'il est lié à $p^{(1)}$ par la relation bien simple
\[
p^{(2)} = p^{(1)} - 1.
\]

Mais au sujet de cette relation une remarque est nécessaire; car elle n'est pas vraie sans restriction, si par $p^{(1)}$ et $p^{(2)}$ on entend le genre de la courbe canonique (supposée irréductible), et le nombre des intersections \emph{variables} de deux courbes canoniques. C'est ce que les exemples suivants nous montreront.

Reprenons à cet effet la surface du cinquième ordre qui a deux points de contact avec elle-même (surface dont nous avons parlé au n${}^{\text{o}}$ 22). Les courbes canoniques sont ici les sections de la surface déterminées par les plans passant par les points singuliers; elles sont par suite $\infty^{1}$ courbes du genre $p^{(1)} = 2$; et pourtant on a $p^{(2)} = 0$, car deux courbes de ce faisceau n'ont aucune intersection \emph{variable}. La relation précédente $p^{(2)} = p^{(1)} - 1$ est donc en défaut.

Un autre exemple, relatif à une valeur plus élevée de $p_{g}$, ne sera pas inutile. Désignons par $1, 2, \ldots, 8$ les huit intersections de trois surfaces du second ordre, et construisons une surface $F^{6}$ d'ordre 6, qui ait un point multiple d'ordre trois dans $1, 2, \ldots, 7$ et un point simple dans 8; cela est toujours possible. Or on voit que les surfaces d'ordre $6 - 4 = 2$ adjointes à $F^{6}$, sont les surfaces du second ordre qui passent par $1, 2, \ldots, 7$, et en conséquence par 8. Ici l'on a $p_{g} = 3$, $p^{(1)} = 4$, $p^{(2)} = 2 = p^{(1)} - 2$.

\textbf{*)} \emph{Zur Theorie}... pag.\ 520, ibid.

\textbf{**)} Ainsi par exemple pour la surface générale d'ordre 5, on a
\[
p_{g} = 4, \quad p^{(1)} = 6.
\]

\textbf{***)} \emph{Nöther} l.\ c.

\textbf{†)} l.\ c.

\origpage{282}

On voit par ces exemples que la relation $p^{(2)} = p^{(1)} - 1$ est en défaut, et doit être remplacée par l'inégalité $p^{(2)} < p^{(1)} - 1$, toutes les fois que le système canonique possède sur la surface des points-base en dehors des points multiples, ou, ce qui revient au même, s'il contient des courbes exceptionnelles de la seconde espèce.

\subsection*{25. Définition numérique de $p^{(1)}$ et $p^{(2)}$.}

Les nombres invariants d'une surface dont nous avons parlé jusqu'ici, à savoir $p_{g}$, $p^{(1)}$ et $p^{(2)}$, ont été introduits par la considération directe du système canonique et de ses caractères géométriques. Or à côté de ces invariants \emph{géométriques}, il convient, dans la théorie des surfaces, de considérer d'autres invariants définis par des expressions numériques formées avec les caractères d'un système quelconque de courbes sur la surface. Ces nouveaux invariants, que nous appellerons \emph{numériques} ou \emph{virtuels}, sont tels, d'ordinaire, qu'ils coïncident avec les invariants géométriques correspondants, lorsque sont remplies certaines conditions de \emph{régularité} de la surface. Ils fournissent au contraire de nouveaux caractères d'invariance pour les surfaces \emph{irregulières}, et parfois ils peuvent remplacer les invariants géométriques dans les cas où la définition de ceux-ci n'a plus de sens. Le plus important de ces invariants numériques est le (premier) \emph{genre numérique} $p_{n}$, qui correspond à $p_{g}$; nous y consacrerons quelques uns des paragraphes suivants. Occupons-nous d'abord des caractères numériques $p_{n}^{(1)}$, $p_{n}^{(2)}$, qui se présentent à côté de $p^{(1)}$ et $p^{(2)}$.

Partons de l'hypothèse que le système canonique $|K|$ de notre surface existe $(p_{g} > 0)$, et soit irréductible; et considérons sur la surface un système quelconque $|C|$ irréductible, ainsi que son système adjoint $|C'|$. Pour plus de simplicité supposons encore que $|C|$ et $|C'|$ sur notre surface n'aient pas de points-base, et que la surface ne possède aucune courbe exceptionnelle: on aura alors exactement
\[
|K| = |C' - C|.
\]
On peut calculer le genre $p_{n}^{(1)}$ et le degré $p_{n}^{(2)}$ de $|K|$ en fonction des caractères analogues $\pi$, $n$ de $|C|$ et $\pi'$, $n'$ de $|C'|$, pourvu que l'on recoure aux rélations du n${}^{\text{o}}$ 9. Ainsi l'on trouvera aisément
\[
\left\{
\begin{aligned}
p_{n}^{(1)} &= \pi' + n - 2\pi + 3,\\
p_{n}^{(2)} &= n' + n - 4\pi + 4.
\end{aligned}
\right.
\tag{1}
\]

Dans les hypothèses où nous nous sommes placés, les deux caractères $p_{n}^{(1)}$ et $p_{n}^{(2)}$ ne diffèrent pas de $p^{(1)}$ et $p^{(2)}$, et jouissent par suite de la propriété d'invariance. Mais il faut remarquer que les relations précédentes

\origpage{283}
gardent un sens déterminé, lors même qu'on fait abstraction des restrictions imposées au système canonique $|K|$. Ce système peut bien faire défaut $(p_{g} = 0)$, ou, quoique existant, il peut différer du système $|C' - C|$ par la présence de courbes exceptionnelles, il peut même être réductible; néammoins on pourra toujours construire les expressions (1) à l'aide d'un système $|C|$ et de son adjoint $|C'|$. Mais on ne peut plus affirmer maintenant que les valeurs (1) coïncident avec les valeurs de $p^{(1)}$ et $p^{(2)}$ qui peuvent même n'avoir aucun sens. Il y a donc à examiner si, même dans ce cas, les expressions (1) gardent la propriété d'invariance. Il en est bien ainsi. On démontre en effet le théorème suivant${}^{*)}$:

\emph{Soit $F$ une surface ne possédant aucune courbe exceptionnelle, et $|C|$ un système quelconque qui n'ait pas de points-base sur $F$; les valeurs $p_{n}^{(1)}$, $p_{n}^{(2)}$ calculées à l'aide des caractères du système $|C|$, ne changent pas si l'on remplace ce système par un autre satisfaisant à la même condition;} $p_{n}^{(1)}$, $p_{n}^{(2)}$ sont par suite deux \emph{invariants} de la surface${}^{**)}$. \emph{Ils satisfont toujours à la relation}
\[
p_{n}^{(2)} = p_{n}^{(1)} - 1.
\]

Que l'on ait au contraire une surface $F$ possédant des courbes exceptionnelles; on pourra de même définir les caractères $p_{n}^{(1)}$, $p_{n}^{(2)}$ de $F$, pourvu que l'on parvienne à transformer birationnellement la surface $F$ en une autre $F'$ qui n'ait pas de telles courbes; car il suffira ensuite d'envisager, par exemple, sur $F'$ le système des sections planes. Or cette transformation de $F$ en $F'$ est possible dans bien de cas, même si le système canonique fait défaut $(p_{g} = 0)$. Pas toujours pourtant; il y a en effet des surfaces (par exemple, les surfaces rationnelles, et les surfaces réglées) qui possèdent certainement des courbes exceptionnelles, en quelque manière qu'on les transforme. Pour ces surfaces les définitions de $p_{n}^{(1)}$ et $p_{n}^{(2)}$ que nous avons données, sont en défaut. On pourrait vraiment les modifier de sorte qu'elles s'étendent aussi à ces cas; mais il ne convient pas de nous arrêter ici sur des recherches qui sont encore incomplètes.

\textbf{*)} \emph{Enriques}, \emph{Introduzione}... n${}^{\text{o}}$ 41.

\textbf{**)} Sur la surface du cinquième ordre ayant deux contacts avec elle-même (n${}^{\text{o}}$ 22, 24) un système tel que $|C|$ est fourni par les sections planes de la surface $(\pi = 6, n = 5; \pi' = 12, n' = 16)$; on trouvera $p_{n}^{(1)} = 2$, $p_{n}^{(2)} = 1$, tandis qu'on avait $p^{(1)} = 2$, $p^{(2)} = 0$.

Un autre exemple digne de remarque nous est donné par la surface du sixième ordre qui passe deux fois par les six arêtes d'un tetraèdre. Cette surface dont nous aurons à parler plus tard, n'a pas de système canonique $(p_{g} = 0)$ et par suite $p^{(1)}$ et $p^{(2)}$ n'ont aucun sens; pourtant si l'on part du système des sections planes de la surface $(\pi = 4, n = 6; \pi' = 4, n' = 6)$ on trouve $p_{n}^{(1)} = 1$, $p_{n}^{(2)} = 0$.

\origpage{284}

\subsection*{26. Le genre numérique.}

Nous allons maintenant définir un nouvel invariant numérique d'une surface, qui jouit aussi d'une grande importance.

Partons d'un système $|C|$ situé sur la surface, et de son système adjoint $|C'|$. Nous avons remarqué jadis (n${}^{\text{o}}$ 16) que la série $g_{2\pi-2}$ découpée par $|C'|$ sur la courbe générale $C$ de $|C|$, n'est pas nécessairement complète; au contraire elle peut avoir un certain défaut $\delta_{0}$, et alors sa dimension se réduit à $\pi - 1 - \delta_{0}$ (et la dimension de $|C'|$ se réduit à $p_{g} + \pi - 1 - \delta_{0}$).

De même les séries découpées par les systèmes $|C' + C|$, $|C' + 2C| \ldots$ sur la courbe peuvent avoir des défauts positifs $\delta_{1}$, $\delta_{2}, \ldots$ La succession des nombres positifs $\delta_{0}$, $\delta_{1}$, $\delta_{2}, \ldots$ n'est pas pourtant illimitée, car nous savons (n${}^{\text{o}}$ 19) qu'on peut toujours choisir un nombre $i$ assez élevé, pour que les systèmes $|C' + iC|$, $|C' + (i+1)C| \ldots$ découpent sur $C$ une série complète. Il y a donc lieu à considérer la somme
\[
\Delta = \delta_{0} + \delta_{1} + \cdots + \delta_{i-1}
\]
formée avec les défauts des séries découpées sur la courbe générale d'un système $|C|$ par les systèmes $|C'|$, $|C' + C| \ldots$ qui sont adjoints à $|C|$, $|2C| \ldots$.

Eh bien, cette somme $\Delta$, qui est tout-à-fait déterminée par le système $|C|$, ne change pas lorsqu'on remplace le système $|C|$ par un autre système quelconque de la surface; en d'autres termes, $\Delta$ est un \emph{nouvel invariant} de la surface${}^{*)}$. Bien que $\Delta$ lui-même entre dans plusieurs questions, toutefois en nous conformant à l'usage établi, nous introduirons au lieu de $\Delta$ un nouveau caractère, le \emph{genre numérique} $p_{n}$, qui est défini par la relation
\[
p_{g} - p_{n} = \Delta.
\]

\emph{On obtient donc le genre numérique d'une surface en retranchant du genre géométrique la somme des défauts des séries, qui sont découpées sur la courbe générale d'un système linéaire par les systèmes adjoints à celui-ci et à ses multiples.}

D'après cette définition on a
\[
p_{n} \leqq p_{g};
\]
$p_{n}$ peut même être négatif, tandis que $p_{g}$ ne l'est jamais.

Par exemple une surface réglée ayant les sections planes de genre $\pi$, a, comme nous avons remarqué, $p_{g} = 0$; de plus, si $|C|$ est le système des sections planes, on trouve $\delta_{0} = \pi$, $\delta_{1} = \delta_{2} = \cdots = 0$, et par

\textbf{*)} \emph{Enriques}, \emph{Introduzione}... n${}^{\text{o}}$ 40.

\origpage{285}
suite $\Delta = \pi$, $p_{n} = -\pi$; \emph{le genre numérique d'une surface réglée est le genre des sections planes de celle-ci précédé du signe négatif.}

La définition de $p_{n}$ que nous venons de donner, peut aussi se présenter sous la forme projective; nous verrons par là que ce nouveau caractère ne se distingue pas du genre numérique introduit par \emph{Cayley}, \emph{Zeuthen} et \emph{Nöther}${}^{*)}$, et cela donnera raison du nom que nous avons attribué à $p_{n}$.

Il suffit, à cet effet, de supposer que le système $|C|$, dont nous nous servons pour calculer $\Delta$, soit formé par les sections planes d'une surface $F$ de l'espace ordinaire, dont nous envisageons, à present, les caractères projectifs, l'ordre $n$, etc. Les systèmes $|C'|, |C'+C|, \ldots, |C'+iC|$ sont alors découpés sur $F$ par les surfaces $\Phi^{\nu}$ d'ordre $\nu = n-3$, $n-2, \ldots n-3+i$ adjointes à $F$. La remarque d'après laquelle les systèmes $|C'+iC|, |C'+(i+1)C| \ldots$ découpent sur $C$ une série complète (pour $i$ assez élevé), se traduit maintenant dans la propriété du système des $\Phi^{\nu}$ de découper sur un plan général \emph{toutes} les courbes d'ordre $\nu$ adjointes à la section de $F$, lorsque $\nu \geqq n-3+i$; tandis que pour les valeurs de $\nu < n-3+i$ les $\Phi^{\nu}$ découpent sur le plan une partie seulement des courbes adjointes, si $\delta_{i-1}, \delta_{i-2}, \ldots$ sont supérieurs à zéro. Or les courbes d'ordre $\nu$ adjointes à la section plane générale de $F$ forment un système qui a la dimension
\[
\varrho_{\nu} = \binom{\nu+2}{2} - 1 - k,
\]
où $k$ est une constante qui dépend seulement des singularités de la section nommée, c'est-à-dire des courbes multiples de la $F$; $\nu$ peut avoir une valeur quelconque $\geqq n-3$. Il s'ensuit (par un calcul bien simple) que la dimension du système des $\Phi^{\nu}$ pour $\nu \geqq n-3+i$ est
\[
r_{\nu} = \binom{\nu+3}{3} - 1 - k\nu + k', \tag{1}
\]
où $k'$ est une nouvelle constante qui dépend des courbes et des points multiples isolés de $F$. Le second membre de la (1), pour les valeurs de $\nu < n-3+i$, ne fournit plus (en général) la dimension exacte, \emph{effective}, du système des $\Phi^{\nu}$ (car alors il faut faire attention aux défauts $\delta_{i-1}, \delta_{i-2}, \ldots$); il fournit seulement un caractère qu'on pourrait appeler \emph{dimension virtuelle} (numérique), et qui coïncide avec la dimension \emph{effective} seulement dans le cas où $\delta_{i-1} = \delta_{i-2} = \cdots = 0$.

Précisément on trouve, sans difficulté, que la dimension effective du système des $\Phi^{n-4}$ est
\[
r_{n-4} = \binom{n-1}{3} - 1 - k(n-4) + k' + \sum_{0}^{i-1} \delta_{h},
\]

\textbf{*)} Voir les citations au n${}^{\text{o}}$ 4.

\origpage{286}
d'où l'on déduit sur le champ (en ayant égard à la définition $\displaystyle p_{g} - p_{n} = \sum_{0}^{i-1} \delta_{h}$)
\[
p_{n} = \binom{n-1}{3} - k(n-4) + k'.
\]
Cette relation s'énonce par le théorème:

\emph{Le genre numérique d'une surface n'est autre chose que la dimension virtuelle du système des surfaces adjointes d'ordre $n-4$, augmentée d'une unité.}

C'est précisément cette dimension augmentée de 1, qui d'après les auteurs cités, fournit la définition de $p_{n}$.

Au sujet des considérations qui précèdent il ne sera pas inutile d'ajouter quelques remarques.

Nous sommes parvenus à la rélation (1) sans préciser les singularités (courbes et points multiples) de la surface $F$, qu'on avait à considérer; $k$ et $k'$ dépendent de ces singularités, il est vrai, mais les expressions de $k$ et $k'$ en fonction de celles-ci n'ont aucune importance, du moins dans les considérations théoriques que nous venons de faire. Il en serait autrement si l'on voulait calculer en effet $p_{n}$ pour une surface donnée. Mais même dans ce cas on pourrait éviter la determination directe de $k$ et $k'$, si l'on parvenait à déterminer par une voie quelconque les dimensions des systèmes des $\Phi^{\nu}$ en correspondance aux valeurs assez élevées de $\nu$; car nous venons de voir, en conclusion, que \emph{lorsque $\nu$ parcourt la série croissante des nombres entiers} (à partir d'une certaine limite), \emph{ces dimensions parcourent une progression arithmétique du troisième ordre, laquelle prolongée aussi aux valeurs inférieures de $\nu$, contient $p_{n} - 1$ parmi ses termes.}

A la même conclusion arrivent aussi \emph{Cayley} et \emph{Nöther} directement, en calculant les expressions de $k$ et $k'$ au moyen des singularités de la surface, et en les substituant dans la formule (1) (\emph{formule de postulation})${}^{*)}$.

Dès que ce calcul est fait, une voie se présente pour démontrer l'invariance de $p_{n}$; il suffit en effet de refaire un calcul analogue avec les caractères d'une surface transformée de $F$, et de comparer les résultats. C'est la voie que les auteurs cités ont suivie. Nous avons remarqué plus haut les difficultés qu'on y rencontre; nous venons de voir maintenant comment on peut les éviter.

\textbf{*)} Voir à ce propos: \emph{Nöther}, \emph{Sulle curve multiple di superficie algebriche}, Annali di Matem.\ serie II, t.\ 5${}^{\text{o}}$, (1871) pag.\ 172.

\origpage{287}

\subsection*{27. Deux nouvelles définitions du genre numérique; la série caractéristique d'un système linéaire.}

La définition que nous avons donnée plus haut du caractère $p_{g} - p_{n}$, à l'aide d'un système $|C|$ et du système adjoint $|C'|$, peut même se présenter sous la forme suivante.

La série que $|C'|$ découpe sur la courbe générale de $|C|$ peut avoir, nous l'avons dit, un certain défaut $\delta_{0}$. Ce défaut ne reste pas invariable en général, lorsqu'on remplace le système $|C|$ et son adjoint $|C'|$, par un nouveau système $|D|$ et par son adjoint $|D'|$. Considérons toutes les valeurs que $\delta_{0}$ peut prendre en correspondance avec les différents systèmes situés sur la surface. On démontre que ces valeurs ont un \emph{maximum}, qui est naturellement un invariant de la surface; mais ce n'est pas un nouvel invariant; au contraire ce maximum est précisément $p_{g} - p_{n}$. On a donc le théorème${}^{*)}$:

\emph{Le défaut de la série découpée sur la courbe générale d'un système linéaire par le système adjoint, peut dépendre du système que l'on choisit sur la surface; mais dans tous les cas ce défaut atteint un maximum, qui est égal à la différence entre les genres géométrique et numérique de la surface.}

Si l'on compare la définition de $p_{g} - p_{n}$ qui résulte ainsi, avec celle que nous avons donnée ci-dessus, on voit qu'on peut choisir toujours sur une surface un système $|D|$ tellement, que les séries découpées sur la courbe $D$ par $|D'+D|, |D'+2D| \ldots$ résultent complètes${}^{**)}$- On peut même démontrer qu'on parvient à un tel système $|D|$, en prenant un multiple assez élevé d'un système $|C|$ quelconque de la surface.

Dans la dernière définition de $p_{g} - p_{n}$ on fait attention à la série découpée par $|C'|$ sur $C$. Or sur la courbe $C$ on trouve une autre série remarquable, que nous avons définie plus haut (n${}^{\text{o}}$ 7); c'est la \emph{série caractéristique} $g_{n}^{r-1}$ découpée sur la courbe générale du système $|C|$ $\infty^{r}$, par les autres courbes du système. Nous allons maintenant nous occuper de cette série; sera-t-elle toujours complète? ou bien dans le cas contraire, qu'est-ce qu'on peut dire de son défaut?

On suppose naturellement ici que le système $|C|$ soit complet; car dans l'hypothèse contraire, si $|C|$ était contenu dans un système $\infty^{r+\delta}$ du même degré $n$, la série caractéristique $g_{n}^{r-1}$ de $|C|$ serait contenue

\textbf{*)} \emph{Enriques}, \emph{Introduzione}... n${}^{\text{o}}$ 40. On peut vérifier ce théorème sur les surfaces \emph{hyperelliptiques} $(p_{g} = 1, p_{n} = -1)$; voir à ce sujet \emph{Humbert}, \emph{Théorie générale des surfaces hyperelliptiques} n${}^{\text{o}}$ 130, 144; l.\ c.

\textbf{**)} \emph{Enriques}. l.\ c.; pour les surfaces hyperelliptiques voir \emph{Humbert}, l.\ c., n${}^{\text{o}}$ 144. Pour les surfaces reglées voir ci-dessus, n${}^{\text{o}}$ 26.

\origpage{288}
à son tour dans la série caractéristique $g_{n}^{r+\delta-1}$ du nouveau système, et aurait par suite un défaut $\geqq \delta$. Or si $|C|$ est complet, la série caractéristique peut être complète, ou bien ne l'être pas; cela dépend de la surface sur laquelle le système est situé. C'est le premier cas qui se présente, par exemple, si la surface est rationnelle; nous l'avons remarqué plus haut pour les systèmes de courbes planes (n${}^{\text{o}}$ 11), et il s'agit ici naturellement d'une propriété invariable par rapport aux transformations birationnelles. La même propriété est verifiée sur les surfaces générales de l'ordre $n$, etc. Considérons au contraire une surface réglée irrationnelle, qui ne soit pas un cône, par exemple la surface du quatrième ordre ayant deux droites directrices doubles; les sections planes de cette surface forment bien un système linéaire complet, et pourtant la série caractéristique $g_{4}^{2}$ n'est pas complète; au contraire elle a le défaut 1. Quel est donc le lien qui passe entre le défaut de la série caractéristique et les invariants d'une surface donnée?

Si l'on considère plusieurs systèmes linéaires complets sur une même surface, et l'on fait attention aux défauts de leurs séries caractéristiques, on trouvera parfois que ces défauts diffèrent l'un de l'autre. Que l'on prenne le maximum parmi les défauts qui correspondent à tous les systèmes de la surface; ce maximum, qui est fini et existe toujours (comme on peut le démontrer), est naturellement un invariant de la surface. Sera-t-il un invariant nouveau? On pourrait le penser. Eh bien, il n'en est pas ainsi; car on montre que ce maximum est précisément égal à la différence $p_{g} - p_{n}$. On a donc le théorème suivant${}^{*)}$:

\emph{Le défaut de la série caractéristique d'un système complet situé sur une surface, peut dépendre du système que l'on choisit sur la surface; mais dans tous les cas ce défaut atteint un maximum, qui est égal à la différence entre les genres géométrique et numérique de la surface.}

\subsection*{28. Les surfaces régulières.}

Les deux théorèmes qui précèdent donnent sur le champ les propriétés fondamentales des surfaces qui ont $p_{g} = p_{n}$, surfaces que nous appellons \emph{régulières}. On a en effet:

a) \emph{Sur une surface régulière le système adjoint à un système linéaire quelconque (irréductible) découpe sur la courbe générale de celui-ci la}

\textbf{*)} \emph{Castelnuovo}, \emph{Alcuni risultati}... n${}^{\text{o}}$ 7; on y trouvera la démonstration seulement pour le cas $p_{g} = p_{n}$; l'Auteur se propose d'exposer prochainement la démonstration complète. On peut vérifier le théorème sur les surfaces hyperelliptiques (\emph{Humbert}, l.\ c., n${}^{\text{o}}$ 114), et sur les surfaces réglées (\emph{Segre}, Mathem.\ Annalen, 34, n${}^{\text{o}}$ 3.).

\origpage{289}
\emph{série canonique complète;} il a la dimension $p_{g} + \pi - 1$, si $\pi$ est le genre de la dernière courbe.

b) \emph{Sur une surface régulière la série caractéristique de tout système linéaire complet est elle-même complète.}

Chacun des deux théorèmes peut servir à caractériser une surface régulière. Pour ce qui concerne le premier, on peut même dire davantage, car \emph{en général la connaissance d'un seul système $\infty^{r} (r \geqq 3)$, sur les courbes duquel le système adjoint découpe la série canonique complète, est suffisante pour affirmer que la surface est régulière}${}^{*)}$. Les mots \emph{en général} se rapportent aux restrictions suivantes qu'on rencontre dans la démonstration: 1) le système n'a pas de courbes fondamentales propres; 2) les courbes du système qui passent par un point de la surface, ne passent pas en conséquence par d'autres points variables avec celui-ci. Est-ce qu'il y a là des restrictions nécessaires?

\subsection*{29. Sur quelques classes de surfaces irrégulières.}

La nécessité de considérer les deux caractères $p_{g}$ et $p_{n}$ d'une surface, dépend de cela qu'il y a de nombreux exemples de surfaces (\emph{irrégulières}) ayant $p_{g} > p_{n}$. Nous avons déjà cité à ce propos la surface réglee à sections de genre $\pi$, pour laquelle on a $p_{g} = 0$, $p_{n} = -\pi$. Mais on connaît maintenant deux familles très étendues de surfaces irrégulières.

La première famille, qui comprend aussi les surfaces réglées, se compose des surfaces qui contiennent un faisceau irrationnel de courbes, c'est à dire un système $\infty^{1}$ pas linéaire, tel que par tout point général de la surface passe une seule courbe du système. \emph{La connaissance d'un faisceau irrationnel sur une surface suffit toujours pour affirmer que la surface est irrégulière.} En d'autres termes: \emph{sur une surface régulière tout faisceau de courbes est un système linéaire}${}^{**)}$.

Pour construire la surface la plus générale contenant un faisceau irrationnel, il suffit de transformer une surface réglée irrationnelle à l'aide d'une transformation $(1, n)$ rationnelle en un seul sens (exprimant les coordonnées d'un point de la surface réglée par des fonctions rationnelles des coordonnées d'un point de la surface transformée${}^{***)}$).

\textbf{*)} \emph{Castelnuovo}, \emph{Alcuni risultati}... n${}^{\text{o}}$ 6.

\textbf{**)} \emph{Castelnuovo}, \emph{Alcuni risultati}... n${}^{\text{o}}$ 10. La connaissance d'un faisceau irrationnel suffit aussi pour reconnaître que la surface possède des intégrales de différentielles totales de première espèce.

\textbf{***)} On trouvera des exemples de telles surfaces dans une Note de M.\ \emph{Castelnuovo}, \emph{Osservazioni intorno alla geometria sopra una superficie}, Rendic.\ Istituto Lombardo, serie II, vol.\ 24, pag.\ 135.

\origpage{290}

Une seconde famille de surfaces irrégulières s'obtient en considérant les surfaces, dont les points correspondent, élément par élément, aux couples de points d'une courbe quelconque de genre $p > 0$. Ainsi, si l'on part d'une courbe elliptique $(p = 1)$, on parvient à la surface réglée elliptique $(p_{g} = 0, p_{n} = -1)$; si l'on part de la courbe de genre $p = 2$, on obtient une surface hyperelliptique $(p_{g} = 1, p_{n} = -1)$${}^{*)}$; etc.

Il existe certainement d'autres familles de surfaces irrégulières, mais jusqu'à présent aucune, à ce que nous croyons, n'a formé le sujet d'une étude particulière.

\subsection*{30. Sur de nouveaux systèmes invariants qui appartiennent à une surface.}

Nous avons parlé jusqu'à présent d'un seul système invariant situé sur une surface donnée, à savoir du système canonique $|K|$. Mais on peut tout de suite introduire de nouveaux systèmes invariants, car il suffit de considérer les multiples $|K_{i}| = |iK|$ du système canonique. Il ne vaudrait pas même la peine de faire attention à ces nouveaux systèmes, s'il n'arrivait parfois de rencontrer des systèmes analogues sur des surfaces qui ne possèdent pas de système canonique $(p_{g} = 0)$. Dans ce cas pourtant on ne peut définir $|K_{i}|$ comme un multiple de $|K|$, car celui-ci n'existe plus; mais il suffit de recourir à la définition de $|K|$, laquelle s'étend sur le champ. Nous avons démontré que si $|C|$ est un système quelconque sur une surface, et $|C'|$ est son adjoint, le système résiduel $|K| = |C'-C|$ ne dépend pas du système $|C|$ d'où l'on part; de même l'on voit directement que le système $|K_{i}| = |iC'-iC|$ ne dépend pas de $|C|$ (à des courbes exceptionnelles près), et peut donc s'obtenir aussi en retranchant le système $|iD|$ du système $|iD'|$, où $|D'|$ est l'adjoint de $|D|$${}^{**)}$. Le système $|K_{i}|$ jouit donc de la propriété d'invariance; on va l'appeler \emph{système i-canonique}; $i$ est un nombre entier, positif, quelconque. Outre que par la relation $|K_{i}| = |iC'-iC|$, ce système est aussi défini par la relation $|K_{i}| = |C^{(i)}-C|$ (à des courbes exceptionnelles près), où $C^{(i)}$ est le système que l'on obtient de $|C|$ en appliquant de suite $i$ fois l'opération d'adjonction. Le système \emph{i-canonique} nous fournit de nouveaux caractères invariants de la surface; tels sont par exemple le nombre $P_{i}$ des courbes i-canoniques qui sont linéairement indépendantes, le genre de ces courbes, etc. Pour $i = 1$, $P_{i} = P_{1}$ nous donne le genre géométrique $p_{g}$ de la surface; pour $i = 2$ on a un caractère $P_{2}$ qu'on pourrait appeler \emph{bigenre}, et qui va jouer (comme on verra) un

\textbf{*)} \emph{Humbert}, \emph{Théorie générale des surfaces hyperelliptiques}, Journal de Mathém., 4${}^{\text{e}}$ série, t${}^{\text{e}}$ IX; voir aussi pour $p = 3$ deux Notes du même auteur dans les Comptes Rendus de l'Ac.\ d.\ Sc., Février 1895.

\textbf{**)} \emph{Enriques}, \emph{Introduzione}..., n${}^{\text{o}}$ 39.

\origpage{291}
rôle fondamental dans la théorie des surfaces rationnelles; etc. Quelques-uns de ces nouveaux caractères peuvent s'exprimer en fonction des anciens; pas tous pourtant; c'est ce que nous allons reconnaître au sujet de $P_{2}$.

Mais remarquons d'abord que si une surface $F$ de l'ordre $n$, est donnée (par ses caractères projectifs) dans l'espace ordinaire, le système i-canonique sur $F$ peut être découpé par des surfaces de l'ordre $i(n-4)$, qui se comportent d'une façon \emph{convenable} le long des courbes multiples et dans les points multiples de $F$. Préciser le comportement des surfaces \emph{i-adjointes} dans le cas général, exigerait beaucoup de détails relatifs aux singularités élevées; qu'on nous permette de nous borner au cas $i = 2$, dans l'hypothèse que la surface $F$ ne possède d'autres singularités qu'une courbe double, et des points d'ordre $k = 3, 4, \ldots$ qui sont d'ordre $\dfrac{k(k-1)}{2}$ pour la courbe double. \emph{Dans ce cas le système bicanonique est découpé sur $F$ par les surfaces d'ordre $2(n-4)$ qui passent deux fois par la courbe double de $F$.}

Revenons maintenant aux relations entre le système canonique $|K|$ et le système i-canonique $|K_{i}|$. Si le premier système existe, le second existera de même, quel que soit $i$, car on a $|K_{i}| = |iK|$, et en conséquence $P_{i} \geqq p_{g}$; donc si $p_{g} > 0$, on a aussi $P_{i} > 0$. Supposons, au contraire, qu'on ait $p_{g} = 0$; $|K|$ n'existe plus. Alors il y a à distinguer deux cas. Il peut arriver en effet que $|K_{i}|$ n'existe non plus, quel que soit $i$; c'est ce qu'on peut vérifier tout de suite sur le plan, sur une surface réglée quelconque, plus généralement sur toute surface contenant un système $|C|$ de genre $\pi$, dont deux courbes ont plus que $2\pi - 2$ intersections variables; dans cette hypothèse, en effet, le système $|iC|$ ne peut être contenu dans le système $|iC'|$. Mais il peut arriver au contraire que $|K_{i}|$ existe, en relation avec quelques valeurs de $i$, lors même que $|K|$ n'existe pas; en d'autres termes, il peut arriver que l'on ait $p_{g} = 0$, $P_{i} > 0 (i > 1)$. C'est là une propriété singulière, qui n'a pas son analogue dans la théorie des courbes algébriques. Qu'elle se présente parfois, c'est ce que nous allons reconnaître sur quelques exemples, en nous bornant à la valeur $i = 2$.

Un premier exemple, bien simple, est fourni par \emph{la surface du sixième ordre qui passe deux fois par les six arêtes d'un tetraèdre}, et en conséquence trois fois par les sommets de celui-ci${}^{*)}$; la surface a le genre géomètrique $p_{g} = 0$, car il n'existe pas de surface du second ordre passant simplement par les six arêtes. Il existe au contraire une surface bi-adjointe de l'ordre 4, qui est formée par les quatre faces du tetraèdre; par suite l'on a $P_{2} = 1$. La dernière surface ne découpe pas la surface donnée en

\textbf{*)} \emph{Enriques}, \emph{Introduzione}..., n${}^{\text{o}}$ 39; \emph{Castelnuovo}, \emph{Sulle superficie di genere zero}, n${}^{\text{o}}$ 15.

\origpage{292}
dehors des lignes multiples; on dira donc que la courbe bicanonique correspondante a l'ordre zéro.

Voici un second exemple${}^{*)}$. Considérons dans l'espace une droite $d$, une section conique $k$ ne rencontrant pas la droite, et quatre points $A_{1}$, $A_{2}$, $A_{3}$, $A_{4}$ en de positions générales. Construisons maintenant une surface de l'ordre 7 qui passe trois fois par $d$ et deux fois par $k$, avec la condition ultérieure que deux nappes de la surface touchent en $A_{i}$ le plan $A_{i}d$ $(i = 1, 2, 3, 4)$. On démontre tout d'abord qu'une telle surface existe; ensuite qu'elle a le genre $p_{g} = 0$ et le bigenre $P_{2} = 2$; enfin que les courbes bicanoniques sont $\infty^{1}$ courbes du quatrième ordre et de genre 1, découpées par les plans passant par la droite $d$.

\subsection*{31. Sur quelques relations entre les caractères invariants.}

Les caractères principaux d'une surface que nous venons de définir $(p_{g}, p_{n}, P_{2}, p^{(1)})$, quoique indépendants deux à deux, satisfont à quelques inégalités, qu'il est bon de connaître. Nous avons remarqué déjà les deux suivantes, qui découlent de la définition,
\[ p_{g} \geqq p_{n}, \quad P_{2} \geqq p_{g}. \]
On peut même dire davantage, si l'on fait attention à ce que le système bicanonique est l'adjoint du système canonique; en effet, si l'on suppose que celui-ci soit irréductible, on obtient, lorsque $p_{g} > 1$,
\[ p_{g} + p^{(1)} \geqq P_{2} \geqq p_{n} + p^{(1)}. \]

Cette relation, qui devient une égalité pour les surfaces régulières, ne subsiste pas toujours si $p_{g} = 1$, même si l'on recourt à la définition virtuelle de $p^{(1)}$. Précisément, en nous bornant au cas $p_{g} = p_{n} = 1$${}^{**)}$, on trouve
\[ P_{2} = p^{(1)} + 1, \quad \text{ou} \quad P_{2} = p^{(1)}, \]
selon qu'il existe, ou n'existe pas, une courbe canonique d'ordre $> 0$ (en dehors des courbes exceptionnelles).

Par suite le minimum de $p^{(1)}$ pour les surfaces régulières de genre 1, est 1. Il y a deux espèces de telles surfaces pour lesquelles le minimum est atteint; elles correspondent aux valeurs $P_{2} = 2$, $P_{2} = 1$${}^{***)}$. Les surfaces de la seconde espèce sont surtout remarquables. On peut les représenter sur une surface $F$, d'un certain ordre $n$, ne possédant

\textbf{*)} \emph{Castelnuovo}, l.\ c.

\textbf{**)} \emph{Voir} pour ce cas \emph{Enriques}, \emph{Sui piani doppi di genere uno}, n${}^{\text{os}}$ 4, 5.

\textbf{***)} Un exemple de surfaces de la première espèce est fourni par la surface du cinquième ordre du n${}^{\text{o}}$ 23. Quant aux surfaces de la seconde espèce, on a tout d'abord la surface générale du quatrième ordre, ensuite la surface du sixième ordre qui passe doublement par l'intersection de deux surfaces de l'ordre 2 et 3; etc.

\origpage{293}
aucune courbe exceptionnelle; la $F$ n'est pas rencontrée par la surface adjointe de l'ordre $n-4$, en dehors des courbes multiples. Sur $F$ tout système linéaire complet $|C|$, qui n'a pas de points-base, jouit d'une propriété remarquable; \emph{si $\pi$ est le genre de $|C|$, la dimension de $|C|$ a la même valeur $\pi$, et deux courbes $C$ se rencontrent en $2\pi-2$ points $(r=\pi, n=2\pi-2)$. Le système $|C|$ est l'adjoint de lui-même.}

Revenons maintenant aux relations d'inégalité entre les caractères d'une surface, et considérons $p_{g}$ et $p^{(1)}$, en supposant, par simplicité, que le système canonique soit irréductible $(p_{g} > 1)$.

M.\ \emph{Nöther}${}^{*)}$ a démontré qu'on a toujours
\[ p^{(1)} \geqq 2p_{g} - 3. \]
\emph{Si le minimum de $p^{(1)}$ est atteint, les courbes canoniques sont hyperelliptiques} (et \emph{viceversa}); elles se rencontrent deux à deux en
\[ \frac{p^{(1)} - 1}{2} = p_{g} - 2 \]
couples de points conjugués. Par suite si $p_{g} > 2$, on peut prendre comme équation de la surface (ou d'une transformée de celle-ci)
\[ z^{2} = f(x, y), \]
où $f$ est un polynôme. On peut même préciser davantage${}^{**)}$, car \emph{on peut s'arranger de manière que}

1) \emph{ou bien $f$ ait le degré 6 en $x$ (et le degré total $2n > 6$)},

2) \emph{ou bien $f$ ait le degré total 8 ou 10 (respect.\ $p_{g} = 3, 6$)}.

Dans le cas 1) les plans $y = const.$ découpent sur la surface un faisceau de courbe du genre 2.

Qu'on écarte maintenant les surfaces pour lesquelles
\[ p^{(1)} = 2p_{g} - 3; \]
on démontre alors que $p^{(1)}$ ne peut descendre au dessous de $3p - 6$, à savoir
\[ p^{(1)} \geqq 3p_{g} - 6; \]
et précisément${}^{***)}$:

\emph{Toute surface ayant}
\[ p^{(1)} = 3p_{g} - 6 \qquad (p_{g} > 3) \]
\emph{peut être représentée par une equation}
\[ f(x, y, z) = 0 \]
\emph{du degré 4 par rapport à $x$, $y$ (ou par une transformée de celle-ci); sauf dans le cas de certaines surfaces du genre $p_{g} = 4$ (surface générale}

\textbf{*)} \emph{Zur Theorie}... Math.\ Ann.\ 8, p.\ 252.

\textbf{**)} \emph{Enriques}, \emph{Sopra le superficie algebriche di cui le curve canoniche sono iperellittiche}. Rendic.\ Accad.\ Lincei 1896.

\textbf{***)} \emph{Castelnuovo}, \emph{Osservazioni interno alla geometria sopra una superficie}, Nota II, Rendic.\ Istit.\ Lombardo 1896.

\origpage{294}
\emph{du cinquième ordre), $p_{g} = 5$, $p_{g} = 7$.} Dans le cas général, les plans $z = const.$ découpent sur la surface un faisceau de courbes du genre 3 (et d'ordre 4).

\section*{Chapitre VI. Systèmes linéaires spéciaux et non-spéciaux. Extension aux surfaces du théorème de Riemann-Roch.}

Avant de quitter la théorie générale des surfaces, il faut que nous consacrions quelques mots aux questions qui se rapportent à l'extension du théorème de Riemann-Roch aux surfaces.

Mais d'abord examinons quelle est la vraie nature de ce théorème, lorsqu'on se borne aux courbes.

\subsection*{32. Séries spéciales et non-spéciales sur une courbe.}

Par rapport à la série canonique $g_{2p-2}^{p-1}$ appartenant à une courbe de genre $p$, les séries situées sur la même courbe se répartissent en deux familles${}^{*)}$. A la première famille appartiennent les séries qui sont contenues dans la série canonique, et même celle-ci; on les appelle \emph{séries spéciales}; les séries spéciales ont naturellement l'ordre $\leqq 2p - 2$ et la dimension $\leqq p - 1$. L'autre famille est formée par les séries qui ne sont pas contenues dans la série canonique; ce sont les \emph{séries non-spéciales}, dont l'ordre et la dimension n'ont pas de limite supérieure.

Si l'on veut voir clairement la différence qui passe entre ces deux familles, il convient de recourir aux séries complètes.

\emph{Soit $g_{n}^{r}$ une série complète sur une courbe du genre $p$; alors}

1) \emph{si $g_{n}^{r}$ est non-spéciale on a la relation}
\[ r = n - p, \]
qui permet de calculer un des trois caractères $n$, $r$, $p$ à l'aide des autres;

2) \emph{si, au contraire, $g_{n}^{r}$ est spéciale, l'égalité précédente est à remplacer par l'inégalité}
\[ r > n - p. \]
Et, inversement, si la dernière inégalité subsiste entre les caractères d'une série quelconque (complète ou incomplète), on peut affirmer que la série est spéciale.

La distinction que nous venons de faire, peut même se présenter sous la forme suivante. Regardons $n - p$ comme la \emph{dimension virtuelle} (théorique) d'une série complète d'ordre $n$ sur une courbe de genre $p$,

\textbf{*)} \emph{Brill} et \emph{Nöther}, \emph{Ueber die algebraischen Functionen}..., Math.\ Ann.\ 7.

\origpage{295}
et comparons la dimension virtuelle avec la \emph{dimension effective} $r$ que la série possède en effet. Alors nous voyons que si la série est non-spéciale les deux dimensions coïncident, et la série est, pour ainsi dire, \emph{régulière}; si, au contraire, la série est spéciale la dimension effective est supérieure à la dimension virtuelle, et la série est irrégulière.

Pour une série spéciale complète $g_{n}^{r}$, il y a lieu à considérer la différence entre les deux dimensions
\[ i = r - (n-p) = p - n + r, \]
On démontre que \emph{$i$ est précisément le nombre des groupes canoniques linéairement indépendants qui contiennent un groupe général de la série $g_{n}^{r}$}. C'est là le théorème connu qui porte le nom de \emph{Riemann-Roch}. Le nombre $i$ est un caractère de la série $g_{n}^{r}$, qu'on pourrait appeler le \emph{degré de spécialité}; pour une série non-spéciale on a $i = 0$, pour la série canonique $g_{2p-2}^{p-1}$, $i = 1$; etc.

En profitant de la notion des séries résiduelles (n${}^{\text{o}}$ 9), on peut énoncer le théorème de \emph{Riemann-Roch} sous la forme suivante, qui est due à MM.\ \emph{Brill} et \emph{Nöther}, et qui se prête le mieux aux applications géométriques:

\emph{Toute série spéciale complète $g_{n}^{r}$ située sur une courbe de genre $p$, a, par rapport à la série canonique, une série résiduelle $g_{n'}^{r'}$ dont les caractères ont les valeurs suivantes}
\[ n' = 2p - 2 - n, \quad r' = p - 1 - (n-r) = i - 1. \]

\subsection*{33. Quelques remarques sur les systèmes linéaires de courbes planes.}

La répartition des séries en deux familles s'étend tout de suite aux systèmes linéaires de courbes sur une surface, et donne lieu à la notion de \emph{systèmes spéciaux} et \emph{systèmes non-spéciaux}, dont les premiers sont contenus dans le système canonique (et ont par suite la dimension $\leqq p_{g} - 1$, le genre $\leqq p^{(1)}, \ldots$), tandis que les seconds n'y sont pas contenus. Or on saurait s'attendre de trouver pour les systèmes non-spéciaux une \emph{régularité} analogue à celle, que nous avons rencontrée dans les séries non-spéciales. Mais il en est tout autrement. C'est ce que nous allons reconnaître tout d'abord dans le cas le plus simple, sur le \emph{plan} $(p_{g} = 0)$, où tous les systèmes sont non-spéciaux.

Un système $|C|$ complet de courbes planes d'ordre $m$ est entièrement défini, dès que l'on a fixé les points-base du système, et les ordres de multiplicité que les courbes $C$ ont en ces points. Or pour calculer la dimension de $|C|$, du moins avec quelque approximation, on peut d'abord évaluer le nombre des conditions, que chaque point-base, séparé des autres, présente aux courbes d'ordre $m$; on additionnera

\origpage{296}
les nombres qui correspondent à tous les points-base, et on retranchera la somme du nombre $\frac{1}{2} m(m+3)$, qui exprime la dimension du système formé par toutes les courbes de l'ordre $m$. Désignons par $\varrho$ le résultat de ce calcul; est-ce que $\varrho$ nous fournira exactement la dimension $r$ de $|C|$? Certainement si les points-base présentent des conditions toutes indépendantes aux courbes d'ordre $m$; c'est ce qui arrive, par exemple, s'il n'y a pas de points-base, ou si ceux-ci se trouvent en des points généraux du plan; on a alors la relation $r = \varrho$. Mais il n'en est plus ainsi, si $\omega$ parmi les conditions imposées par les points-base sont une conséquence des autres; car on trouve alors $r = \varrho + \omega$. Un exemple de ce cas nous est fourni par le système des courbes d'ordre $m$, qui passent par les $m\mu$ intersections d'une d'entre elles avec une courbe générale d'ordre $\mu < m$; on a en effet
\[ r = \frac{(m-\mu)(m-\mu+3)}{2} + 1, \]
\[ \varrho = \frac{m(m+3)}{2} - m\mu, \]
et, par suite,
\[ \omega = r - \varrho = \frac{(\mu-1)(\mu-2)}{2}. \]

Quoique le nombre $\varrho$ puisse, dans bien de cas, différer de la dimension exacte $r$ d'un système linéaire, il convient pourtant d'y faire attention, car il est un nouvel invariant du système (par rapport aux transformations birationnelles). Nous appelons $\varrho$ la \emph{dimension virtuelle} du système $|C|$, $r$ la \emph{dimension effective}. En conséquence nous répartissons les systèmes linéaires complets de courbes planes (qui sont tous non-spéciaux) en deux familles${}^{*)}$:

1) la première famille comprend les \emph{systèmes réguliers}, pour lesquels la dimension effective est égale à la dimension virtuelle, $r = \varrho$; ils correspondent, en une certaine façon, aux séries non-spéciales (ou régulières) sur une courbe;

2) la seconde famille comprend les systèmes \emph{surabondants}, dont la dimension effective surpasse la dimension virtuelle, $r = \varrho + \omega$; $\omega > 0$ est la \emph{surabondance} du système. Il ne convient pas de regarder ces systèmes comme les analogues des séries spéciales; au contraire, ils en diffèrent par bien de propriétés. Par exemple, en cela que sur une courbe de genre zéro il n'y a pas de séries spéciales, tandis que sur le plan $(p_{g} = p_{n} = 0)$ nous trouvons des systèmes surabondants; en cela, encore, que l'ordre et la dimension d'une série spéciale ne peuvent surpasser certaines limites, tandis que le degré, la dimension, le genre

\textbf{*)} \emph{Castelnuovo}, \emph{Ricerche generali sopra i sistemi lineari di curve piane}, l.\ c.\ n${}^{\text{o}}$ 1.

\origpage{297}
d'un système surabondant peuvent être aussi élevés, que l'on désire. Il convient plutôt de regarder la \emph{surabondance}, comme une \emph{cause d'irrégularité} d'une autre nature que la \emph{spécialité}, considérée sur les courbes, bien qu'elle ait le même effet d'augmenter la dimension effective par rapport à la dimension virtuelle.

La distinction que nous avons faite entre les systèmes réguliers et surabondants, se traduit en une distinction entre les séries caractéristiques de ces systèmes. On a en effet le théorème${}^{*)}$:

\emph{Un système est régulier ou bien surabondant, selon que la série caractéristique du système est non-spéciale ou bien spéciale;} d'où il s'ensuit, par exemple, qu'un système surabondant de courbes de genre $\pi$ a le degré $\leqq 2\pi - 2$ et la dimension effective $\leqq \pi - 1$${}^{**)}$.

En général si l'on désigne (comme d'ordinaire) par $n$ et $r$ le degré et la dimension effective d'un système complet, on a pour un système régulier l'égalité
\[ \pi - n + r - 1 = 0, \tag{1} \]
tandis que pour les systèmes surabondants on a
\[ \pi - n + r - 1 = \omega > 0 \tag{2} \]
\[ (\omega \text{ étant la surabondance}); \]
dans tous les cas
\[ \varrho = n - \pi + 1 \tag{3} \]
est la dimension virtuelle du système.

\subsection*{34. Systèmes réguliers sur une surface.}

La distinction que nous venons de faire entre les systèmes \emph{réguliers} et \emph{surabondants} sur le plan (ou sur les surfaces rationnelles), peut maintenant se transporter sur une surface quelconque. Ici l'on rencontre pourtant une complication, qui dépend de la répartition des systèmes en spéciaux et non-spéciaux, dont nous avons parlé plus haut. Pour éviter cette difficulté nous nous bornons tout d'abord aux systèmes \emph{non-spéciaux} (à savoir, aux systèmes qui ne sont pas contenus dans le système canonique).

Soit $|C|$ un système non-spécial, complet, dont nous désignerons les caractères par $n$, $r$, $\pi$, suivant l'habitude. Envisageons l'expression
\[ \pi - n + r - 1 \]
que nous avons déjà considérée sur le plan. Elle peut prendre différentes

\textbf{*)} \emph{Castelnuovo}, l.\ c., n${}^{\text{o}}$ 18.

\textbf{**)} \emph{Segre}, \emph{Sui sistemi lineari}, Circolo Matem.\ di Palermo, t.\ I, 1887.

\origpage{298}
valeurs, selon le système $|C|$ auquel elle se rapporte; mais on démontre${}^{*)}$ que parmi toutes ces valeurs il y a \emph{un minimum, qui est précisément le genre numérique $p_{n}$ de la surface} (pour le plan on a $p_{n} = 0$). Eh bien, nous appelons \emph{régulier} tout système (non-spécial), dont les caractères satisfont à la relation
\[ \pi - n + r - 1 = p_{n}; \tag{4} \]
\emph{surabondant} tout système pour lequel on a
\[ \pi - n + r - 1 = p_{n} + \omega, \tag{5} \]
où $\omega > 0$ désigne la \emph{surabondance} du système. On va considérer tout d'abord les systemes réguliers, car ils sont les plus simples, et ils satisfont à des lois générales fort remarquables. En voici quelques unes${}^{**)}$.

\emph{La série caractéristique d'un système complet régulier a le défaut $p_{g} - p_{n}$}, c'est-à-dire le défaut maximum qui appartient aux systèmes complets (n${}^{\text{o}}$ 27).

\emph{Le système canonique découpe sur la courbe générale d'un système régulier une série complète}, dont la dimension est $p_{g} - 1$, et l'ordre est $2\pi - 2 - n$; c'est la série résiduelle de la série caractéristique par rapport à la série canonique.

En particulier: \emph{si la surface a le genre géometrique $p_{g} = 0$, la série caractéristique d'un système régulier est non-spéciale}; c'est ce qui arrive par exemple sur le plan, ou sur les surfaces réglées.

On peut construire un système régulier dont la dimension est aussi grande que l'on veut; à cet effet on commencera par construire dans l'espace ordinaire une surface $F$, qui soit en correspondance birationnelle avec la surface donnée, et qui n'ait pas de points multiples \emph{isolés} (n${}^{\text{o}}$ 15); ensuite on conduira les surfaces adjointes à $F$. On démontre que ces surfaces découpent sur $F$ un système régulier, aussitôt que leur ordre est assez élevé. Si, par exemple, la $F$ était une surface générale de l'ordre $n$, on verrait aisément que les surfaces générales d'un ordre quelconque $> n - 4$ découpent sur $F$ un système régulier non-spécial.

La définition de la \emph{surabondance} $\omega$ que nous avons donnée, est tout-à-fait arithmétique; nous sommes en mesure maintenant d'en reconnaître la signification géométrique et d'en justifier le nom.

Mais il faut d'abord que nous introduisons un nouveau caractère d'un système linéaire quelconque $|C|$ (spécial ou non-spécial) sur la surface;

\textbf{*)} Pour les surfaces régulières on trouvera démontré ce théorème dans les \emph{Ricerche di geometria}... de M.\ \emph{Enriques}, IV, 2; pour une surface quelconque voir \emph{Castelnuovo}, \emph{Alcuni risultati}..., n${}^{\text{o}}$ 8. Il est à remarquer que ce théorème fournit une nouvelle définition de $p_{n}$.

\textbf{**)} \emph{Castelnuovo}, l.\ c.

\origpage{299}
caractère qui devient la \emph{dimension virtuelle} du système (n${}^{\text{o}}$ 33), dès que la surface est rationnelle, et qui par suite sera désigné dans tous les cas par le même nom.

Construisons à cet effet un système \emph{régulier} $|D|$ assez ample pour contenir $|C|$; c'est ce que nous pouvons faire de plusieurs façons. Pour déduire le système $|C|$ du nouveau système $|D|$, il faudra en général retrancher de $|D|$ une courbe $\gamma$ convenable, et imposer certains points-base au système résiduel $|D-\gamma|$. Il se présente maintenant un moyen pour calculer la dimension $r$ de $|C|$; car il suffit de retrancher de la dimension de $|D|$ le nombre des conditions qu'il faut imposer aux courbes $D$, pour qu'elles contiennent la courbe $\gamma$ et les points-base nommés. Or si nous nous bornons d'abord à une première approximation, nous pouvons supposer: 1) que la série découpée par $|D|$ sur la courbe $\gamma$ soit complète et non-spéciale; 2) que les conditions imposées par les points-base nommés soient toutes indépendantes entre elles et des précédentes. En exécutant le calcul dans ces hypothèses nous parvenons à un résultat $\varrho$, qui ne coïncide pas généralement avec la dimension vraie, \emph{effective}, $r$ de $|C|$, mais qui pourtant ne manque pas d'intérêt. On démontre en effet que l'on parvient au même résultat $\varrho$, si, dans les mêmes hypothèses, on calcule la dimension de $|C|$ en partant d'un système régulier quelconque, autre que $|D|$, qui contienne $|C|$. Par conséquent $\varrho$ est un vrai caractère du système $|C|$; nous l'appellerons \emph{la dimension virtuelle} de $|C|$. Il s'exprime du reste, aisément, à l'aide des autres caractères $n$, $\pi$ de $|C|$, et de $p_{n}$, car on a
\[ \varrho = p_{n} + n - \pi + 1, \tag{6} \]
relation qui comprend, en particulier, la (3) du n${}^{\text{o}}$ 33.

Cela posé, revenons maintenant à notre système $|C|$ non-spécial, dont nous désignons par $r$ et $\varrho$ les dimensions effective et virtuelle. Si nous comparons ces caractères à l'aide des relations (4), (5) et (6), nous trouvons sur le champ que si $|C|$ est régulier, on a $\varrho = r$; si, au contraire $|C|$ est surabondant, avec la surabondance $\omega$, alors $\varrho < r$, et précisément
\[ r - \varrho = \omega. \]

Donc \emph{la surabondance d'un système non-spécial est la différence entre la dimension effective et la dimension virtuelle du système}, d'accord avec ce que nous avons reconnu sur le plan.

Une autre interprétation, celle-ci tout-à-fait géométrique, peut être donnée du caractère $\omega$ sur les surfaces régulières $(p_{g} = p_{n} = p)$. On y parvient en considérant le groupe des $n$ intersections variables de deux courbes générales du système donné $|C|$, et les courbes du système adjoint $|C'|$ qui passent par ce groupe. En effet tandis que le nombre de ces dernières courbes, qui sont linéairement indépendantes, est $2p$,

\origpage{300}
si le système $|C|$ est régulier, le même nombre est $2p + \omega$ si le système $|C|$ a la surabondance $\omega$. Par suite${}^{*)}$:

\emph{La surabondance d'un système non-spécial $|C|$ sur une surface régulière de genre (géom.\ $=$ num.) $p$, est l'excès sur $2p$ du nombre des courbes adjointes $C'$ linéairement indépendantes qui passent par les intersections variables de deux $C$.} On vérifiera le théorème tout de suite sur le plan.

\emph{Remarque.} --- Supposons que la surface $F$ (régulière ou non) soit donnée dans l'espace ordinaire, et ne possède que des singularités ordinaires. D'après le \emph{Restsatz} (n${}^{\text{o}}$ 12), tout système $|C|$ sur $F$ peut être construit à l'aide des surfaces adjointes (ou sous-adjointes), d'un ordre assez élevé, qui passent par une courbe fixe $\gamma$ et par les points-base de $|C|$. Alors la \emph{dimension virtuelle} $\varrho$ de $|C|$ (et par suite sa surabondance $\omega$) peut être calculée d'après les \emph{formules de postulation} de M.\ \emph{Nöther} (n${}^{\text{o}}$ 26).

\subsection*{35. Systèmes spéciaux.}

Les propriétés que nous venons d'énoncer sur les systèmes non-spéciaux, s'étendent aussi aux systèmes spéciaux, pourvu que l'on ait égard à un nouveau caractère qui se présente dans ce cas, et qui est tout-à-fait analogue au caractère $i$, que nous avons introduit dans l'étude des séries spéciales. Un système spécial complet $|C|$ est contenu, par définition, dans le système canonique $|K|$; il y a donc lieu à considérer le nombre $i$ des courbes canoniques linéairement indépendantes qui contiennent la courbe générale de $|C|$; ce nombre $i$, qui désigne, pour ainsi dire, le \emph{degré de spécialité} de $|C|$ est le nouveau caractère. Si $|C|$ est non-spécial, on a $i = 0$; si $|C|$ est le système canonique, alors $i = 1$; etc. Au lieu de $i$, on peut aussi considérer la dimension $r_{1}$ du système $|C_{1}| = |K-C|$ résiduel de $|C|$ par rapport au système canonique; on a naturellement
\[ r_{1} = i - 1. \]

A l'aide du caractère $i$ nous pouvons étendre de la façon suivante les formules que nous avons trouvées à propos des systèmes non-spéciaux. Désignons par $n$, $\pi$, $r$, $i$ les caractères d'un système spécial $|C|$, et par $p_{n}$ le genre numerique de la surface; on a alors
\[ \pi - n + r - 1 \geqq p_{n} - i, \]
et précisément si
\[ \pi - n + r - 1 = p_{n} - i + \omega, \tag{7} \]
on appelle $\omega \geqq 0$ la \emph{surabondance} du système; pour $i = 0$ on retombe

\textbf{*)} \emph{Enriques}, \emph{Ricerche di geometria}... IV.\ 2.

\origpage{301}
dans la formule des systèmes non-spéciaux. On voit d'ici qu'\emph{un système est certainement spécial si}
\[ \pi - n + r - 1 < p_{n}; \]
mais, à cause de $\omega$, on ne peut dire que cette relation soit aussi nécessaire; au contraire de ce qui arrivait sur les courbes, où l'inégalité $n - r < p$ pouvait servir de définition d'une série spéciale complète.

On peut écrire l'égalité (7) de façon à mettre en évidence la dimension $r$ du système $|K-C|$. On parvient au résultat suivant:

\emph{Sur une surface de genre numérique $p_{n}$, tout système spécial complet ayant les caractères $n$, $r$, $\pi$, a, par rapport au système canonique, un système résiduel dont la dimension est}
\[ r_{1} = p_{n} - (\pi-n+r) + \omega, \tag{8} \]
\emph{où $\omega$ (surabondance du premier système) désigne un nombre $\geqq 0$.}

Il faut regarder ce théorème, comme l'extension aux surfaces du théorème de \emph{Riemann-Roch} énoncé pour les courbes, car celui-ci nous donnait aussi la dimension $r_{1}$ de la série résiduelle d'une série $g_{n}^{r}$ par rapport à la série canonique $(r_{1} = p - 1 - (n-r))$. Pour les surfaces nous voyons que la formule est compliquée par la présence du nombre $\omega$, qui n'a pas son analogue dans la théorie des courbes${}^{*)}$${}^{**)}$.

L'égalité (7) peut encore se transformer en introduisant la dimension virtuelle $\varrho$ du système $|C|$, que nous avons définie plus haut pour un système quelconque à l'aide de la relation (6) du n${}^{\text{o}}$ 34. On trouve ainsi la relation
\[ r - \varrho = \omega - i, \]

\textbf{*)} L'extension du théorème de R.-R.\ aux surfaces est due à M.\ \emph{Nöther} (Comptes Rendus de l'Ac.\ des Sc.\ 1886), dans l'hypothèse sousentendue que la série caractéristique du système $|C|$ soit complète; nous savons à présent que cette hypothèse est vérifiée toujours sur une surface régulière $(p_{g}=p_{n}=p)$. M.\ \emph{Nöther} présente la relation du théorème sous la forme
\[ r_{1} \geqq p + n - \pi - r; \]
c'est M.\ \emph{Enriques} (\emph{Ricerche di geometria}... IV, 2) qui a montré la signification géométrique de la différence $\omega$ entre les deux membres de l'inégalité; voir à ce propos les lignes qui suivent ci-dessus.

\textbf{**)} Outre la dimension $r_{1}$ du système $|K-C|$, résiduel de $|C|$ par rapport au système canonique, on peut exprimer aussi les autres caractères $n_{1}$ (degré), $\pi_{1}$ (genre), $\omega_{1}$ (surabondance), du nouveau système à l'aide des caractères du système $|C|$. On n'a qu'à recourir aux relations du n${}^{\text{o}}$ 9. On trouvera ainsi
\[ \begin{aligned} \pi_{1} &= p^{(1)} - 3(\pi-1) + 2n, \\ n_{1} &= p^{(1)} - 1 - 4(\pi+1) + 3n, \\ \omega_{1} &= \omega, \end{aligned} \]
où $p^{(1)}$ est le genre du système canonique (\emph{Curvengeschlecht}). Voir à ce sujet \emph{Enriques}, l.\ c.

\origpage{302}
laquelle nous dit que \emph{pour un système spécial la surabondance s'obtient en ajoutant le caractère $i$ à la différence entre les dimensions effective et virtuelle du système.} Cette définition de $\omega$ se réduit à celle que nous avons exposée au sujet des systèmes non-spéciaux, lorsque $i = 0$.

Pour les surfaces régulières $(p_{g} = p_{n} = p)$ on peut même étendre aux systèmes spéciaux la seconde définition de la surabondance, que nous avons énoncée pour les systèmes non-spéciaux. On n'a qu'à remplacer, dans celle-ci, $\omega$ par $\omega - i$.

\subsection*{36. Aperçu sur les considérations qui précèdent.}

En ayant égard aux complications que peut présenter un système linéaire de courbes sur une surface algébrique, il ne paraîtra pas inutile que nous résumons ici brièvement les remarques, que nous venons de faire à ce sujet.

Si $n$ et $\pi$ désignent le degré et le genre d'un système complet $|C|$ sur une surface de genre numérique $p_{n}$, il faut faire attention d'abord à la dimension virtuelle de $|C|$, qui est fournie par l'expression
\[ \varrho = p_{n} + n - \pi + 1. \]

Ensuite on doit avoir égard à deux causes d'irrégularité, dont l'effet est de modifier la dimension effective $r$ de $|C|$ par rapport à la dimension virtuelle $\varrho$. Ces deux causes ont du reste des effets opposés, qui parfois peuvent s'éliminer. Elles peuvent agir soit séparément, soit en même temps. Elles sont:

1) la \emph{surabondance}, qui fait \emph{augmenter} la dimension du système $(r = \varrho + \omega)$; elle peut se présenter sur une surface quelconque, et pour des systèmes de dimension, de genre, et d'ordre aussi élevés que l'on veut; rien d'analogue n'a lieu pour les séries linéaires sur une courbe;

2) la \emph{spécialité}, à laquelle on doit, au contraire, une \emph{diminution} de la dimension $(r = \varrho - i)$; propriété d'autant plus remarquable que dans la théorie des séries linéaires sur une courbe, la cause analogue a pour effet d'augmenter la dimension $(r = n - \pi + i)$. Sur les surfaces, du reste (de même que sur les courbes), on n'a lieu à considérer cette cause d'irrégularité que pour les systèmes ayant leurs caractères $n$, $\pi$ et $r$ inférieurs à certaines limites (qui sont respectivement $p^{(2)}$, $p^{(1)}$ et $p_{g} - 1$).

Les seuls systèmes sur lesquels ni l'une ni l'autre de ces causes n'agissent, sont les \emph{systèmes réguliers} $(r = \varrho, \omega = i = 0)$.

\origpage{303}

\section*{Chapitre VII. Sur les surfaces rationnelles et sur quelques plans doubles.}

Nous consacrons ce dernier Chapitre de notre Monographie à l'étude de la classe des surfaces rationnelles; classe remarquable aussi parce que les recherches auxquelles elle a donné lieu, ont conduit, par extension naturelle, à des recherches sur les surfaces algébriques générales. Nous dirons ensuite quelques mots sur les plans doubles.

La théorie des surfaces rationnelles a commencé par l'étude de quelques surfaces, particulières au point de vue projectif, qui pouvaient se représenter birationnellement sur le plan; on cherchait de lire sur la représentation les principales propriétés de telles surfaces. Ainsi l'on trouve considérée tout d'abord la surface du second ordre (\emph{Plücker}, 1847; \emph{Cayley} et \emph{Chasles}, 1861), et quelques ans plus tard la surface générale du troisième ordre (\emph{Cremona}, \emph{Clebsch}, 1865); suivent la surface réglée du troisième ordre, la surface de \emph{Steiner}, etc.

Ayant reconnu sur ces exemples le profit que l'on pouvait tirer de la représentation d'une surface pour l'étude de celle-ci, ayant appris à vaincre les difficultés que pouvait présenter, au premier abord, la lecture des propriétés de la surface sur la représentation, le problème (projectif) d'étudier une surface d'après sa représentation connue, pouvait se regarder comme résolu d'avance dans tous les cas. Cependant un autre problème (appartenant celui-ci à la \emph{Géométrie sur une surface}) se présentait: déterminer si une surface donnée par quelqu'un de ses caractères est rationnelle, ou bien ne l'est pas.

C'est à \emph{Clebsch} et à M.\ \emph{Nöther} qu'on doit les propositions fondamentales sur ce sujet. \emph{Clebsch}${}^{*)}$ a montré comment on pouvait établir le caractère rationnel d'un grand nombre de surfaces, en représentant d'abord ces surfaces sur \emph{un plan double} (voir n${}^{\text{o}}$ 38), et en cherchant ensuite les conditions de rationnalité d'un plan double. M.\ \emph{Nöther}${}^{**)}$ a accompli ces recherches, et a donné le théorème capital${}^{***)}$, d'après lequel est rationnelle toute surface contenant un faisceau linéaire de courbes rationnelles.

On est parvenu plus tard à étendre ces propositions, ainsi qu'en ajouter d'autres d'un caractère très général, que nous allons exposer dans les pages suivantes. Mais il faut pourtant remarquer que la plupart des nouveaux résultats s'appuyent en conclusion sur les théorèmes de \emph{Clebsch} et \emph{Nöther}.

\textbf{*)} \emph{Ueber den Zusammenhang}... Math.\ Annalen, 3, 1870.

\textbf{**)} \emph{Ueber die ein-zweideutigen Ebenentransformationen}, Sitzungsber.\ d.\ phys.\ medicin.\ Soc.\ zu Erlangen, 1878; \emph{Ueber eine Classe}..., Math.\ Annalen, 33.

\textbf{***)} \emph{Ueber Flächen, welche Schaaren rationaler Curven besitzen}. Math.\ Ann., 3, 1870.

\origpage{304}

\subsection*{37. Quelques propriétés des surfaces rationnelles.}

Les propriétés générales des surfaces que nous venons d'étudier dans les Chapitres précédents, fournissent des propriétés des surfaces rationnelles dès qu'on particularise les genres de la surface $(p_{g} = p_{n} = 0)$. Ces mêmes propriétés, du reste, peuvent se deduire de l'étude directe des systèmes de courbes planes; et la dernière voie est naturellement préférable au point de vue de la simplicité.

D'une manière ou d'autre, on parviendra aux propositions suivantes que nous réunissons ici, bien que quelques unes d'entre elles aient été déjà énoncées.

\emph{Sur une surface rationnelle:}

a) \emph{la série caractéristique de tout système linéaire complet, est complète;}

b) \emph{la série (canonique) découpée sur la courbe générale d'un système linéaire par le système adjoint, est elle-même complète.}

Chacun de ces deux théorèmes revient à dire qu'\emph{une surface rationnelle a les deux genres $p_{g}$, $p_{n}$ égaux}; chacun peut se déduire de l'autre.

c) \emph{Sur une surface rationnelle il existe des systèmes linéaires dont la série caractéristique est non-spéciale, et parmi ceux-ci il y a des systèmes de genre $\pi$ dont le degré $n > 2\pi - 2$;} tel est par exemple, le système formé par toutes les courbes planes d'un ordre donné.

Le théorème c) nous dit que \emph{le genre géométrique $p_{g}$ $(= p_{n})$ d'une surface rationnelle est nul.} Il nous dit davantage: non seulement il n'y a pas de système canonique sur la surface, mais ni le système bicanonique, ni le système i-canonique ($i$ quelconque) n'existent non plus; de sorte que l'on a (en adoptant les mêmes désignations qu'au n${}^{\text{o}}$ 30)
\[ p_{g} = 0, \quad P_{2} = 0, \quad P_{3} = 0 \ldots . \]

Nous verrons bientôt que ces conditions ne sont pas, du reste, toutes indépendantes.

d) \emph{Sur une surface rationnelle il existe des systèmes linéaires de genre $\pi$ quelconque.} Pour $\pi = 0$ on trouve des systèmes ayant la dimension $r$, et le degré $n = r - 1$ aussi hauts que l'on veut; il n'en est plus ainsi pour $\pi > 0$, car \emph{un système de courbes de genre 1 a la dimension et le degré $\leqq 9$; tandis que pour un système de genre $\pi > 1$, la dimension ne peut surpasser $3\pi + 5$, ni le degré ne peut surpasser $4\pi + 4$;} mais on peut atteindre ces \emph{maxima}${}^{*)}$.

\textbf{*)} Pour $\pi = 1$ le théorème dépend de la réduction des systèmes de courbes elliptiques planes à des types bien connus; voir les travaux de MM.\ \emph{Bertini} (Annali di Matematica, s${}^{\text{e}}$ 2${}^{\text{a}}$.\ t.\ 8, 1877), \emph{Guccia} (1886), \emph{Martinetti} (1887) (Rendic.\ Circolo Mat.\ di Palermo, t.\ I), \emph{Jung} (Rendic.\ Istituto Lombardo 1887---88). Pour $\pi > 1$ les limites supérieures de $n$ et $r$ ont été assignées par M.\ \emph{Castelnuovo} (\emph{Massima dimensione}..., Annali di Matem.\ s${}^{\text{e}}$ 2${}^{\text{a}}$, t.\ 18, 1890).

\origpage{305}

Sans insister ultérieurement sur les propriétés des surfaces rationnelles, nous allons aborder maintenant une question très importante, qui est, en quelque sorte, l'inverse des précédentes: il s'agit de décider si une surface donnée est rationnelle, par l'examen de quelques caractères connus de la surface. Le problème, ainsi posé, n'est pas bien défini; il faudra particulariser les caractères de la surface que l'on regarde connus; c'est ce que nous allons faire de plusieurs façons.

\subsection*{38. On connaît le degré d'un système linéaire de courbes situées sur la surface.}

Souvent c'est la connaissance d'un système linéaire de courbes sur la surface, qui permet de décider la rationnalité de celle-ci; parfois même (dans les cas les plus simples) la connaissance d'un caractère du système (par exemple, du degré $n$, ou du genre $\pi$) suffit; d'après le caractère auquel on fait attention, on parvient à des classes différentes de résultats.

Supposons, tout d'abord, qu'il existe sur la surface un système $|C|$ dont le \emph{degré} soit $n > 0$; la dimension $r$ du système sera $\geqq 2$ et $\leqq n + 1$. Examinons successivement les premières valeurs qu'on peut attribuer à $n$.

1) Les premiers cas que nous rencontrons, correspondent à $n = 1$, $r = 2$, et $n = 2$, $r = 3$; ceux-ci conduisent à des surfaces rationnelles, ainsi qu'on voit sur le champ.

2) Arrêtons-nous plutôt sur le cas plus intéressant, où l'on a $n = 2$, $r = 2$. Si nous faisons correspondre à chaque courbe $C$ du système, une droite d'un plan, nous allons représenter deux points conjugués de la surface (intersections de deux $C$) sur un point général du plan; la surface de cette façon se représente sur un \emph{plan double}, c'est-à-dire sur une surface dont l'équation peut s'écrire sous la forme $z^{2} = R(x, y)$, $R$ étant une fonction rationnelle (que l'on peut supposer entière). C'est donc la question de la rationnalité d'un plan double que nous rencontrons maintenant.

Un plan double est entièrement défini par sa \emph{courbe de diramation} (\emph{Uebergangscurve}) $R(x, y) = 0$, qui est le lieu d'un point dont les deux points correspondants sur la surface coïncident. Quelle particularité doit présenter la courbe de diramation pour que le plan double soit rationnel (représentable sur le plan simple)? La question, abordée par

\origpage{306}
\emph{Clebsch}, a été entièrement résolue par M.\ \emph{Nöther}${}^{*)}$. Voici la réponse:

\emph{Afin qu'un plan double soit rationnel, il faut et il suffit que la courbe de diramation puisse être ramenée, par une transformation birationnelle du plan, à l'une des trois courbes qui suivent:}

a) \emph{courbe d'un certain ordre $n$ ayant un point multiple d'ordre $n - 2$;}

b) \emph{courbe du quatrième ordre;}

c) \emph{courbe du sixième ordre ayant deux points triples infiniment prochains.}

Les dégénérations de ces trois cas sont admises, excepté la scission a${}_{1}$) de la courbe d'ordre $n$ en $n$ droites d'un faisceau, b${}_{1}$) de la courbe du quatrième ordre en quatre droite d'un faisceau, c${}_{1}$) de la courbe du sixième ordre en six droites d'un faisceau, ou en trois coniques d'un faisceau avec deux contacts fixes; car ces dégénerations amènent à des plans doubles, qui représentent des surfaces réglées${}^{**)}$.

Venons maintenant aux valeurs $> 2$ du degré $n$, mais supposons en même temps $r > 2$, car jusqu'à present on ne sait rien sur la rationnalité des plans triples, $\ldots$ $(n = 3, \ldots, r = 2)$. Supposons plus précisément qu'on puisse construire une surface $F$ d'ordre $n$ de l'espace ordinaire, dont les sections planes appartiennent au système $|C|$ donné, qui a le degré $n$ et la dimension $r$.

Il s'agit d'examiner dans quels cas la surface $F$ est rationnelle. La question a été résolue entièrement pour les valeurs les plus petites de $n(> 2)$.

3) On sait par exemple que \emph{toute surface de l'ordre $n = 3$ est rationnelle, abstraction faite du cône général du troisième ordre.}

4) On sait aussi qu'\emph{une surface de l'ordre $n = 4$ ayant des lignes multiples} (au moins une droite double) \emph{est rationnelle, excepté les surfaces réglées de genre $> 0$. Parmi les surfaces du quatrième ordre qui n'ont pas de courbes multiples, il y a seulement quatre familles de surfaces rationnelles, qui ont été déterminées par} M.\ \emph{Nöther}${}^{***)}$. Une d'entre elles se compose des surfaces du quatrième ordre ayant un point triple.

\textbf{*)} Mémoires cités.

\textbf{**)} La correspondance qui passe entre les points d'un plan double rationnel et les couples de points d'un plan simple, donne lieu, sur celui-ci, à une \emph{involution de} $(\infty^{2})$ \emph{couples de points}. Or les types de ces involutions qui sont distinctes par rapport aux transformations birationnelles, ont été déterminés par M.\ \emph{Bertini} (\emph{Ricerche sulle trasformazioni univoche involutorie del piano}. Annali di Matem.\ s${}^{\text{e}}$ 2${}^{\text{a}}$, t.\ 8, 1877); ils correspondent naturellement aux trois types de plans doubles. La possibilité de l'existence d'autres types, en dehors des précédents, qui n'était pas exclue par le Mémoire de M.\ \emph{Bertini}, a été écartée par des recherches ultérieures (\emph{Kantor}, \emph{Castelnuovo}).

\textbf{***)} \emph{Ueber die rationalen Flächen vierter Ordnung}, Math.\ Annalen, 33, 1888.

\origpage{307}
Les trois autres familles correspondent à un point double singulier que la surface possède, par lequel la surface est projectée sur un des trois plans doubles rationnels que nous venons de citer; le premier de ces cas remarquables a été signalé par M.\ \emph{Cremona}.

5) Pour $n \geqq 5$ on n'a que des résultats incomplets; c'est pourquoi nous n'y insistons pas pour le moment, en nous réservant d'indiquer comment la détermination des surfaces rationnelles d'un ordre donné saurait maintenant s'aborder.

\subsection*{39. On connaît le genre d'un système linéaire de courbes situées sur la surface.}

Une autre série de recherches correspond à la classification des systèmes de courbes de la surface, faite d'après le \emph{genre} $\pi$. Pour les premières valeurs de $\pi$ on a des résultats qu'il est bon de rappeler.

1) Un théorème bien connu de M.\ \emph{Nöther} se rapporte au cas $\pi = 0$${}^{*)}$:

\emph{Toute surface contenant un système linéaire $\infty^{1}$ de courbes rationnelles $(\pi = 0)$, est rationnelle.}

En particulier: \emph{toute surface de l'espace ordinaire dont les sections planes sont rationnelles, est rationnelle}; et on peut ajouter (grâce à un théorème de M.\ \emph{Picard}${}^{**)}$) que \emph{la surface est réglée, à moins qu'elle ne soit la surface de Steiner} (du quatrième ordre, avec un point triple et trois droites doubles).

2) Si $\pi = 1$, la connaissance d'un système $\infty^{1}$ de courbes ne suffit pas pour en conclure la rationnalité de la surface (on peut le vérifier sur de nombreux exemples); mais si l'on connaît un système linéaire $\infty^{2}$ de courbes elliptiques $(\pi = 1)$, on peut énoncer un résultat intéressant${}^{***)}$:

\emph{Toute surface contenant un système linéaire $\infty^{2}$ de courbes elliptiques peut se transformer birationnellement en un plan, ou bien en une surface réglée elliptique.} Dans ce dernier cas il existe sur la surface donnée un faisceau elliptique de courbes rationnelles, qui découpent en un seul point variable chaque courbe du système $\infty^{2}$; c'est le faisceau correspondant à la série des génératrices de la surface réglée.

\textbf{*)} l.\ c.\ Math.\ Annalen, 3.

\textbf{**)} \emph{Sur les surfaces algébriques dont toutes les sections planes sont unicursales}, Journal für Mathem., 100, 1885; voir aussi \emph{Guccia}, Rendic.\ Circolo Mat.\ di Palermo, t.\ I, 1886.

\textbf{***)} \emph{Castelnuovo}, \emph{Sulle superficie algebriche che contengono una rete di curve iperellittiche}, Rendic.\ della R.\ Acc.\ d.\ Lincei, 1894; la démonstration est rapportée par M.\ \emph{Enriques} dans le Mémoire \emph{Sui sistemi lineari di superficie algebriche}..., n${}^{\text{o}}$ 8, Math.\ Annalen, 46.

\origpage{308}

En particulier${}^{*)}$: \emph{toute surface de l'espace ordinaire dont les sections planes sont de courbes elliptiques, est rationnelle, ou bien elle est une surface réglée.} Dans le premier cas l'ordre de la surface ne peut surpasser 9.

3) Pour $\pi = 2$ on a un résultat analogue au précédent, pourvu que l'on ait égard à la restriction $n > 2$${}^{**)}$:

\emph{Toute surface contenant un système linéaire de courbes de genre 2 dont la dimension $r \geqq 2$, et le degré $n > 2$, peut se transformer birationnellement en un plan, ou bien en une surface réglée dont le genre peut assumer les valeurs 2 ou 1.} Dans le cas de la surface réglée du genre 2, les génératrices de celle-ci sont découpées en un seul point par toute courbe du système donné. Si au contraire la surface réglée est elliptique, ses génératrices sont rencontrées en plus d'un point par les courbes du système $\infty^{2}$.

\emph{Toute surface de l'espace ordinaire dont les sections planes ont le genre 2, est rationnelle (et contient un faisceau de coniques), ou bien elle est réglée}${}^{***)}$.

4) Les derniers théorèmes s'étendent de suite à toute surface qui contient un système $\infty^{2}$ de courbes \emph{hyperelliptiques}, dont le genre $\pi$ est aussi haut que l'on veut${}^{\dagger)}$:

\emph{Toute surface contenant un système linéaire $\infty^{2}$ de courbes hyperelliptiques de genre $\pi$, dont la série caractéristique est non-spéciale, peut se transformer birationnellement en un plan, ou bien en une surface réglée dont le genre a une des deux valeurs $\pi$, 1.} Au sujet de ce dernier cas, il y a lieu de faire une remarque analogue à la précédente.

\emph{Toute surface de l'espace ordinaire dont les sections planes sont des courbes hyperelliptiques est rationnelle (et contient un faisceau de coniques), ou bien elle est réglée.}

5) Pour $\pi = 3$, le cas hyperelliptique excepté, on n'a jusqu'à présent que des résultats incomplets${}^{\dagger\dagger)}$.

\textbf{*)} \emph{Castelnuovo}, \emph{Sulle superficie algebriche le cui sezioni piane sono curve ellittiche}, Rendic.\ della R.\ Acc.\ d.\ Lincei, 1894; voir aussi \emph{Enriques}, \emph{Sui sistemi lineari}..., n${}^{\text{o}}$ 5.

\textbf{**)} \emph{Castelnuovo}, \emph{Sulle superficie algebriche che contengono una rete}..., l.\ c.; \emph{Enriques} l.\ c.\ n${}^{\text{o}}$ 8.

\textbf{***)} \emph{Enriques}, \emph{Sui sistemi lineari di superficie algebriche}..., Rendic.\ della R.\ Acc.\ d.\ Lincei 1893, Math.\ Annalen 39.

\textbf{†)} Voir pour ces résultats les Mémoires cités de MM.\ \emph{Castelnuovo} et \emph{Enriques}.

\textbf{††)} \emph{Castelnuovo}, \emph{Sulle superficie algebriche le cui sezioni sono curve di genere 3}, Atti dell'Acc.\ d.\ Scienze di Torino, vol.\ 25, 1890.

\origpage{309}

\subsection*{40. Quelques résultats qui se ramènent aux précédents.}

Si l'on connaît sur une surface un système linéaire assez ample de courbes, on peut parfois en déduire un autre qui se trouve dans les conditions considérées plus haut, à l'aide d'une méthode qui a été employée avec succès par M.\ \emph{Del Pezzo}${}^{*)}$, et qui consiste à imposer de nouveaux points-base (simples, doubles...) aux courbes du système. Au point de vue projectif, cela revient à considérer la surface de l'espace $S_{r}$ (à $r$ dimensions), dont les sections par les $S_{r-1}$ forment le système donné, et à projecter la surface par quelques-uns de ses points, de ses plans tangents, etc.

Par ce moyen on voit sur le champ que \emph{toute surface de l'espace $S_{r}(r \geqq 3)$ ayant l'ordre $n = r - 1$, ou $n = r$, est rationnelle, excepté le cône elliptique qui peut se présenter dans le second cas;} on peut même assigner des propriétés projectives intéressantes de ces surfaces${}^{**)}$.

Un autre résultat, qu'on a obtenu par la même méthode${}^{***)}$, complète en quelque sorte le théorème cité plus haut, d'après lequel un système linéaire de genre $\pi > 1$ sur une surface rationnelle, a la dimension $\leqq 3\pi + 5$:

\emph{Si un système linéaire de courbes de genre $\pi \geqq 1$ sur une surface a la dimension $r \geqq 3\pi + 5$, la surface est rationnelle, ou bien elle peut se représenter sur une surface réglée de genre $\pi$, dont les génératrices rencontrent en un seul point les courbes du système. C'est le second cas qui se présente toujours si $r > 3\pi + 5$ pour $\pi > 1$, ou $r > 9$ pour $\pi = 1$.}

On pourrait même, à ce qu'il paraît, obtenir des résultats plus expressifs en relation avec des valeurs inférieures de $r$.

\subsection*{41. Méthode des systèmes adjoints successifs.}

La méthode que nous venons d'exposer, ne se prête pas, dans la plupart des cas, à examiner si une surface, possédant un système donné de courbes, est rationnelle. L'imposition de nouveaux points-base au système amènera d'ordinaire à un nouveau système, dont la dimension résultera aussi petite que l'on voudra, mais dont le degré et le genre

\textbf{*)} \emph{Sulle superficie di ordine $n$ immerse nello spazio di $n + 1$ dimensioni} (Rendic.\ Accad.\ di Napoli, 1885); \emph{Sulle superficie dell'n${}^{\text{o}}$ ordine immerse nello spazio di $n$ dimensioni} (Rendic.\ Circolo Mat.\ di Palermo t.\ I, 1887).

\textbf{**)} \emph{Del Pezzo}, l.\ c.

\textbf{***)} \emph{Enriques}, \emph{Sulla massima dimensione}..., Atti dell'Acc.\ d.\ Scienze di Torino, vol.\ 29.

\origpage{310}
seront encore trop grands pour qu'on puisse en tirer quelque conclusion. Or un procédé de \emph{reduction} ne peut nous satisfaire, s'il ne conduit nécessairement à un système de genre ou de degré assez petits, toutes les fois que la surface considérée est rationnelle; car seulement dans ce cas il nous permettra de décider si une surface donnée est rationnelle, ou bien ne l'est pas. Un tel procédé nous est fourni par l'opération \emph{d'adjonction}, que nous avons définie dans un Chapitre précédent.

Soit $|C|$ un système linéaire de courbes sur la surface donnée; formons le système adjoint $|C'|$, ensuite le système $|C''|$ adjoint de celui-ci, que nous appelons le \emph{second adjoint} de $|C|$, et ainsi de suite. Or, par rapport à la succession des systèmes $|C|$, $|C'|$, $|C''|$, $\ldots$, deux cas peuvent se présenter.

1) Ou bien la succession a un terme; on parvient donc à un dernier système $|C^{(i)}|$ qui n'admet plus de système adjoint;

2) ou bien la succession est illimitée; chacun des systèmes qui la composent admet le système adjoint.

C'est naturellement le premier cas qui va se presenter, lorsque la surface considérée est rationnelle; on le voit sur le champ, si l'on opère sur des systèmes de courbes planes. Mais on ne peut pas dire inversement que la condition 1) soit toujours \emph{suffisante}, pour affirmer la rationnalité de la surface; car il faut ajouter encore la condition de \emph{régularité} de la surface $(p_{g} = p_{n})$. En effet le dernier système $|C^{(i)}|$ aura alors (on le démontre) le genre 0 ou 1, et on pourra appliquer les résultats du n${}^{\text{o}}$ 39.

En conclusion pour reconnaître le caractère rationnel d'une surface, il faut d'abord vérifier qu'elle est régulière (pour cela il suffit en général d'examiner si la série découpée par $|C'|$ sur la courbe $C$ est complète (n${}^{\text{o}}$ 28)).

Ensuite il faut vérifier si la succession $|C|$, $|C'|$, $|C''|$ $\ldots$ a un terme. Ce procédé exige, a ce qu'il paraît, un nombre d'opérations qu'on ne peut pas assigner \emph{a priori}, avant de conduire à une réponse définitive sur la nature de la surface proposée. C'est là un défaut sérieux, que nous tâcherons d'éviter. On y réussit en remarquant que tous les systèmes adjoints peuvent se déduire des deux premiers $|C|$ et $|C'|$ par de simples opérations d'addition et de soustraction; on a en effet (en vertu du théorème fondamental n${}^{\text{o}}$ 16)
\[ |C''| = |2C' - C|, \quad |C'''| = |3C' - 2C|, \quad \text{etc.} \]
C'est donc la considération des systèmes $|C|$ et $|C'|$ qui suffit pour trancher la question de la rationnalité de la surface. Voici comment il faudra procéder.

a) Il faut d'abord reconnaître que $|C'|$ ne contient pas $|C|$;

\origpage{311}

b) il faut encore que $|2C'|$ ne contienne pas $|2C|$; cette condition renferme la précédente, mais elle n'en découle pas.

Toute vérification ultérieure est superflue; en effet on démontre que si la dernière condition est remplie, la succession $|C|$, $|C'|$, $|C''|$ $\ldots$ a un terme; par suite la surface est rationnelle. On a donc le théorème suivant${}^{*)}$:

\emph{Sur une surface régulière on suppose connu un système linéaire $|C|$ de courbes} (du moins $\infty^{3}$), \emph{ainsi que son système adjoint $|C'|$; pour que la surface soit rationnelle, il faut et il suffit que le système $|2C'|$ ne contienne pas le système $|2C|$.}

Nous allons voir tout de suite quelques corollaires de ce théorème; mais remarquons d'abord que la condition nommée est certainement remplie lorsque \emph{le degré $n$ du système $|C|$ est supérieur à $2\pi - 2$ ($\pi$ étant le genre de $|C|$).}

\subsection*{42. Rationnalité d'une surface déduite de la connaissance de ses invariants numériques.}

La conséquence la plus importante que l'on peut tirer du théorème répond à la question suivante: la classe des surfaces rationnelles sera-t-elle entièrement définie par des valeurs particulières des nombres invariants $(p_{g}, p_{n}, P_{2}, P_{3} \ldots)$? \emph{Clebsch} s'était posé la question analogue pour la classe des courbes rationnelles; il l'a résolue en remarquant qu'une courbe est rationnelle si son genre est nul, et \emph{viceversa}. Nous allons voir que pour les surfaces la réponse est analogue, bien qu'il y ait à envisager plus d'un caractère.

Rappelons d'abord que pour une surface rationnelle on a
\[ p_{n} = p_{g} = P_{2} = P_{3} = \cdots = 0; \]
c'est ce qui résulte des définitions de ces caractères. Ces conditions ne sont pourtant pas toutes indépendantes, mais il ne faut penser non plus que l'une d'elle entraine de suite les autres (comme il arrive pour les courbes). En effet les surfaces réglées irrationnelles, par exemple, ont
\[ p_{g} = P_{2} = P_{3} = \cdots = 0, \]
mais $p_{n} < 0$; et l'on connaît de même des surfaces non-rationnelles avec $p_{n} = 0$. On était vraiment porté à croire que les deux conditions $p_{n} = p_{g} = 0$ suffiraient à caractériser la classe des surfaces rationnelles; et c'est seulement lorsqu'on a trouvé des exemples contraires, qu'on a vu qu'il en était tout autrement; ce sont les exemples que nous avons dejà cités au n${}^{\text{o}}$ 30 $(p_{n} = p_{g} = 0, P_{2} > 0)$.

C'est le théorème du paragraphe précédent qui nous fournit la

\textbf{*)} \emph{Castelnuovo}, \emph{Sulle superficie di genere zero}.

\origpage{312}
réponse à notre question; car les hypothèses qu'on y fait se traduisent par les relations $p_{g} = p_{n}$, $P_{2} = 0$ (qui découlent de $p_{n} = 0$, $P_{2} = 0$). On a donc le résultat${}^{*)}$:

\emph{Afin qu'une surface soit rationnelle, il faut et il suffit que le genre numérique $p_{n}$, et le bigenre $P_{2}$ s'évanouissent} $(p_{n} = P_{2} = 0)$.

\subsection*{43. Détermination des surfaces rationnelles d'un ordre donné.}

On peut faire une application du théorème dernier à la détermination de toutes les surfaces rationnelles de l'ordre $n$ dans l'espace ordinaire; cela, du moins, en théorie, car l'exécution pratique exige des notions qu'on ne possède pas pour le moment. La question est réduite à l'analyse des singularités (courbes et points multiples), qu'on peut attribuer à une surface de l'ordre $n$, ayant égard aux conditions: 1) que le genre numérique de la surface soit nul, 2) qu'il n'existe pas de surfaces biadjointes de l'ordre $2(n-4)$, en dehors de celles qui sont formées par la surface donnée avec une autre surface d'ordre $n - 8$. Ainsi la difficulté est ramenée à déterminer l'influence exercitée par une singularité connue sur le genre numérique d'une surface, et le comportement d'une surface biadjointe dans la même singularité. C'est ce qu'on peut faire dans les cas les plus simples; ainsi:

\emph{Si une surface $F$ de l'ordre $n$ possède une courbe double ayant l'ordre $d$ et le genre $\pi$, et possède en outre $t$ points triples, qui sont triples aussi pour la dite courbe (et n'a pas d'autres singularités), les conditions de rationnalité de la surface sont les deux suivantes:}

1) \emph{on doit avoir la relation}
\[ \frac{(n-1)(n-2)(n-3)}{6} - (n-4)d + 2t + \pi - 1 = 0, \]

2) \emph{il ne doit pas exister de surfaces d'ordre $2(n-4)$ qui passent doublement par la courbe double de $F$ (en dehors des surfaces qui contiennent la $F$).}

\subsection*{44. Surfaces les points desquelles ont des coordonnées qui s'expriment en fonctions rationnelles de deux paramètres.}

Le théorème fondamental du n${}^{\text{o}}$ 41 nous permet aussi de répondre à une question, qui se présente tout d'abord dans l'étude des surfaces rationnelles. La définition que nous avons donnée de ces surfaces, se traduit par les deux conditions suivantes:

\textbf{*)} \emph{Castelnuovo}, l.\ c.

\origpage{313}

1) les coordonnées $x$, $y$, $z$ du point général de la surface s'expriment en fonctions rationnelles de deux paramètres $\alpha$, $\beta$ (qui peuvent être regardés comme les coordonnées d'un point du plan),
\[ x = f_{1}(\alpha, \beta), \quad y = f_{2}(\alpha, \beta), \quad z = f_{3}(\alpha, \beta); \tag{1} \]

2) inversement les paramètres $\alpha$, $\beta$ s'expriment rationnellement en fonctions de $x$, $y$, $z$,
\[ \alpha = \varphi_{1}(x, y, z), \quad \beta = \varphi_{2}(x, y, z). \tag{2} \]

Or supposons que, pour une surface donnée, la première condition soit remplie, en d'autres termes supposons que les (1) subsistent; et tâchons de déduire les (2) des (1) par des opérations d'élimination. Si l'on peut effectuer le passage, ce qui arrive dans bien de cas, la surface est rationnelle, d'après la définition. Mais il peut arriver, au contraire, qu'on ne puisse pas déduire les (2) des (1), parce qu'à un point $(x, y, z)$ de la surface correspondent $n > 1$ points $(\alpha, \beta)$ du plan. Est-ce qu'il faudra conclure dans ce cas que la surface n'est pas rationnelle? Ou bien, au contraire, pourra-t-on introduire, au lieu de $\alpha$ et $\beta$, deux nouveaux paramètres $\alpha'$, $\beta'$ (fonctions rationnelles des premiers), de telle façon que $x$, $y$, $z$ s'expriment rationnellement à l'aide de $\alpha'$, $\beta'$, et \emph{viceversa} $\alpha'$, $\beta'$ s'expriment rationnellement à l'aide de $x$, $y$, $z$? Dans ce dernier cas la surface serait naturellement rationnelle, et la particularité qu'on vient de remarquer à l'égard de $\alpha$ et $\beta$, dépendrait seulement du choix spécial des paramètres.

Or on soupçonnait bien qu'il en était ainsi; car dans le même sens avait été résolue, par M.\ \emph{Lüroth}, la question analogue qui se rapporte aux courbes. Mais la démonstration de M.\ \emph{Lüroth}, à ce qu'il paraît, ne peut pas s'étendre aux surfaces.

Voici comment on réussit à démontrer le théorème.

A tout point $(x, y, z)$ de la surface correspondent (nous l'avons dit) $n > 1$ points $(\alpha, \beta)$ du plan; à toute courbe algébrique de la surface correspond une courbe algébrique du plan, et à un système linéaire de courbes sur celle-là, un systèmé linéaire de courbes sur celui-ci. Or des propriétés connues de ces derniers systèmes on peut déduire des propriétés des premiers systèmes. Ainsi l'on voit par exemple que sur la surface tout système linéaire complet a la série caractéristique complète, d'où il résulte que les deux genres $p_{g}$, $p_{n}$ de la surface sont égaux. On voit, de plus, que sur la surface existent des systèmes de genre $\pi$ avec le degré $> 2\pi - 2$, d'où il s'ensuit
\[ p_{g} = P_{2} = P_{3} = \cdots = 0. \]
Les deux conditions $(p_{n} = P_{2} = 0)$ pour la rationalité d'une surface sont donc remplies. C'est ce qui affirme le théorème suivant${}^{*)}$:

\textbf{*)} \emph{Castelnuovo}, \emph{Sulla razionalità delle involuzioni piane}, Math.\ Ann.\ 44, 1893.

\origpage{314}

\emph{Si les coordonnées du point général d'une surface peuvent s'exprimer rationnellement en fonctions de deux paramètres, la surface est rationnelle.}

On peut présenter ce théorème sous une autre forme, pourvu que l'on remarque que les groupes de $n$ points $(\alpha, \beta)$ correspondants aux points $(x, y, z)$ de notre surface, forment sur le plan une \emph{involution}, c'est-à-dire, un système $\infty^{2}$ de groupes, dont chaque groupe est entièrement déterminé par un de ses points:

\emph{Toute involution sur le plan est rationnelle.}

\subsection*{45. Rationnalité d'une surface déduite de la connaissance d'une série, pas linéaire, de courbes.}

Jusqu'ici nous nous sommes bornés à considérer des systèmes linéaires de courbes. Or il y a des cas, où la connaissance d'une série $\infty^{1}$ non-linéaire de courbes sur une surface, série telle que par le point général de la surface il passe plus d'une courbe de la série, peut suffire pour reconnaître la rationalité de celle-ci. Voici quelques exemples:

a) Un premier exemple regarde le cas où sont rationnelles les courbes de la série, ainsi que la série elle-même (en y considérant la courbe comme élément): on a alors le théorème suivant qui se déduit sur le champ du théorème sur les involutions planes (ainsi que celui-ci s'ensuit de celui-là):

\emph{Toute surface qui contient une série rationnelle de courbes rationnelles, est rationnelle}${}^{*)}$.

Ce théorème comprend, comme un cas particulier, le théorème de M.\ \emph{Nöther}, relatif aux surfaces ayant un système \emph{linéaire} $\infty^{1}$ de courbes rationnelles; pour le retrouver il suffit, en effet, de supposer que par tout point de la surface passe une seule courbe de la série.

b) Un second théorème envisage une série dont le degré (nombre des intersections variables de deux courbes) est 1:

\emph{Toute surface qui contient une série algébrique de courbes, dont le degré est 1, et telle que par chaque point de la surface passent $i > 2$ courbes de la série, est rationnelle}${}^{**)}$.

Si, les autres hypothèses restant les mêmes, on suppose $i = 2$, la surface peut se représenter birationnellement sur la variété formée

\textbf{*)} \emph{Castelnuovo}, \emph{Sulle superficie algebriche che contengono una rete di curve iperellitiche}, l.\ c.

\textbf{**)} \emph{Humbert}, Comptes Rendus, 1893; \emph{Sur quelques points de la théorie}..., Journal de Mathématiques 4${}^{\text{e}}$ s${}^{\text{e}}$, t${}^{\text{e}}$ X; \emph{Castelnuovo}, \emph{Sulla linearità delle involuzioni}..., Atti dell'Acc.\ d.\ Sc.\ di Torino, vol.\ 28, 1893.

\origpage{315}
par les couples de points d'une courbe, dont le genre est égal à celui de la série.

c) Au théorème qui précède se rattache le suivant${}^{*)}$:

\emph{Qu'il existe sur une surface deux séries algébriques de courbes, telles que deux courbes générales appartenant à de séries différentes, se rencontrent en un seul point variable; alors si par le point général de la surface il passe plus d'une courbe de chacune des séries, la surface est rationnelle, et les courbes des séries sont elles-mêmes rationnelles.} Si par le point général de la surface passe une seule courbe de la première série, et plus d'une de la seconde, la surface peut être transformée birationnellement en une surface réglée; si enfin par le point passe une seule courbe de chaque série, la surface peut se représenter sur la variété des couples de points de deux courbes.

\subsection*{46. Surfaces admettant un groupe continu de transformations birationnelles en elles-mêmes.}

Pour accomplir, autant que possible, cette revue des résultats obtenus au sujet des surfaces rationnelles, il faut que nous disons encore quelques mots sur les surfaces qui admettent un groupe continu de transformations birationnelles en elles-mêmes, au point de vue de leur caractère rationnel.

M.\ \emph{Picard}${}^{**)}$ a abordé, le premier, le problème des groupes de transformations d'une surface en elle-même, et il est parvenu à des résultats extrêmement remarquables, qui ne trouvent pas ici leur place, car ils concernent principalement certaines classes de surfaces irrationnelles. Pour notre but il suffit en effet d'énoncer les théorèmes suivants, qui donnent l'extension naturelle du résultat bien connu de M.\ \emph{Schwarz} relatif aux courbes:

\emph{Toute surface qui admet un groupe continu transitif de transformations birationnelles en elle-même, groupe dépendant de deux paramètres, est rationnelle, si les transformations ne sont pas échangeables deux à deux.}

\emph{Toute surface qui admet un groupe continu transitif de transformations birationnelles en elle-même, groupe dépendant de plus de deux paramètres, est rationnelle, ou bien peut se transformer en une surface réglée elliptique}${}^{***)}$.

\textbf{*)} \emph{Castelnuovo}, l.\ c.

\textbf{**)} Voir les Mémoires cités dans l'Introduction.

\textbf{***)} \emph{Castelnuovo} et \emph{Enriques}, Comptes Rendus de l'Ac.\ d.\ Sc.\ 1895; voir aussi dans le même recueil (1895) une Note de M.\ \emph{Painlevé}. Au sujet des énoncés que nous rapportons ici, il est bon de remarquer que nous appelons \emph{transitif} (d'après M.\ Lie) un groupe qui puisse changer un point en un autre point arbitrairement fixé; et que le second théorème dit un peu plus de celui qu'on trouve dans les Comptes Rendus.

\origpage{316}

\subsection*{47. Plans doubles du genre un.}

Dans le n${}^{\text{o}}$ 38 nous avons parlé des \emph{plans doubles}, c'est-à-dire des champs algébriques
\[ x, y, \sqrt{f(x, y)}, \]
où $f$ désigne un polynôme; nous avons énoncé les résultats de \emph{Clebsch} et \emph{Nöther} sur la nature de la courbe de diramation $f(x, y) = 0$, lorsque le plan double est rationnel, et en conséquence $\sqrt{f(x, y)}$ devient une fonction rationnelle de deux nouveaux paramètres, par une substitution rationnelle convenable sur $x$, $y$.

Le problème se pose maintenant de classifier les plans doubles supérieurs (d'après leurs courbes de diramation). Les résultats obtenus jusqu'ici sur ce sujet se bornent au cas des plans doubles, dont les genres $p_{n}$ et $P_{2}$ (et par suite $p^{(1)}$, $p_{g}$, $P_{3} \ldots$) ont la valeur 1. On a le théorème${}^{*)}$:

\emph{Les plans doubles dont les genres $(p_{n}, P_{2})$ ont la valeur 1, se répartissent en quatre classes, selon que leur courbe de diramation peut être ramenée, par une transformation birationnelle du plan, à l'une des courbes suivantes:}

1) \emph{courbe générale de l'ordre 6:}

2) \emph{courbe de l'ordre 8 douée de deux points du quatrième ordre;}

3) \emph{courbe de l'ordre 10 douée d'un point de l'ordre 7 et de deux points triples infiniment prochains à celui-là;}

4) \emph{courbe de l'ordre 12 douée d'un point de l'ordre 9 et de trois points triples infiniment prochains à celui-là.}

\emph{Florence}, 13.\ février 1896.

\textbf{*)} \emph{Enriques}, \emph{Sui piani doppi di genere uno}, l.\ c.

\end{document}
